%% dissertacao_tomo_sar.tex
%% Tese de Doutorado
%% Título: Imageamento Tomográfico Subsuperficial por Fusão Coerente
%%         Monoestática--Biestática com Veículos Aéreos Não Tripulados
%%         Projeto da Trajetória do Receptor Orientado por Métricas de Imagem
%% Idioma: Português
%% Compilar: pdflatex -interaction=nonstopmode dissertacao_tomo_sar.tex  (duas passagens)

\documentclass[12pt,a4paper,oneside]{book}

%% ─── Pacotes ─────────────────────────────────────────────────────────────────
\usepackage[utf8]{inputenc}
\usepackage[T1]{fontenc}
\usepackage[portuguese]{babel}
\usepackage{lmodern}
\usepackage{microtype}
\usepackage{geometry}
\geometry{left=3cm, right=2.5cm, top=3cm, bottom=2.5cm}

\usepackage{amsmath,amssymb,amsfonts,amsthm}
\usepackage{mathtools}
\usepackage{bm}
\usepackage{icomma}
\usepackage{siunitx}
\sisetup{output-decimal-marker={,}}

\usepackage{graphicx}
\graphicspath{{figures/}{figures/relatorio_1610/}{figures/relatorio_2809/}}

\usepackage{float}
\usepackage{caption}
\captionsetup{font=small,labelfont=bf}
\usepackage{subcaption}
\usepackage[table,xcdraw]{xcolor}
\usepackage{booktabs}
\usepackage{array}
\usepackage{multirow}
\usepackage{longtable}
\usepackage{tabularx}
\usepackage{enumitem}
\usepackage{setspace}
\onehalfspacing
\setlength{\emergencystretch}{3em}

\usepackage{tcolorbox}
\tcbuselibrary{theorems,skins,breakable}

\usepackage[unicode,colorlinks=true,
  linkcolor=blue!60!black,
  citecolor=green!50!black,
  urlcolor=cyan!60!black,
  pdftitle={Imageamento Tomografico Subsuperficial por Fusao Coerente Monoestatica-Biestatica},
  pdfauthor={David Atencia}]{hyperref}

\usepackage{fancyhdr}
\setlength{\headheight}{15pt}
\pagestyle{fancy}
\fancyhf{}
\fancyhead[L]{\small\nouppercase{\leftmark}}
\fancyhead[R]{\small\thepage}
\renewcommand{\headrulewidth}{0.4pt}

%% ─── Cores das tabelas comparativas ─────────────────────────────────────────
\definecolor{melhor}{HTML}{CDEFD9}
\definecolor{regular}{HTML}{FFF1C2}
\definecolor{pior}{HTML}{F8D3D0}

%% ─── Caixa de síntese ao final de cada estudo ───────────────────────────────
\newtcolorbox{sintese}[1][]{enhanced,breakable,colback=green!4,
  colframe=green!45!black,fonttitle=\bfseries\small,title={Síntese},#1}

%% ─── Teoremas adicionais ─────────────────────────────────────────────────────
\theoremstyle{remark}
\newtheorem*{observacao}{Observação}

%% ─── Operadores ──────────────────────────────────────────────────────────────
\DeclareMathOperator{\sinc}{sinc}
\DeclareMathOperator{\var}{var}
\DeclareMathOperator{\atantwo}{atan2}
\newcommand{\vect}[1]{\boldsymbol{#1}}
\newcommand{\uvec}[1]{\hat{\boldsymbol{#1}}}
\newcommand{\dd}{\mathrm{d}}
\newcommand{\jj}{\mathrm{j}}
\newcommand{\Real}{\mathbb{R}}
\newcommand{\dB}{\,\mathrm{dB}}
\newcommand{\citar}[1]{[\ref{ref:#1}]}
\newcolumntype{L}[1]{>{\raggedright\arraybackslash}p{#1}}

%% ─── Teoremas ────────────────────────────────────────────────────────────────
\theoremstyle{plain}
\newtheorem{teorema}{Teorema}[chapter]
\newtheorem{proposicao}{Proposição}[chapter]
\newtheorem{lema}{Lema}[chapter]
\newtheorem{corolario}{Corolário}[chapter]
\theoremstyle{definition}
\newtheorem{definicaom}{Definição}[chapter]
\newtheorem{hipotese}{Hipótese}[chapter]

%%═══════════════════════════════════════════════════════════════════════════════
\begin{document}
%%═══════════════════════════════════════════════════════════════════════════════

%%─── Capa ─────────────────────────────────────────────────────────────────────
\frontmatter
\begin{titlepage}
  \centering\vspace*{1.5cm}
  {\Large\textsc{Universidade Estadual de Campinas}\par}
  {\large\textsc{Faculdade de Engenharia Elétrica e de Computação}\par}
  \vspace{0.5cm}
  {\large\textsc{Programa de Pós-Graduação em Engenharia Elétrica}\par}
  \vfill
  {\LARGE\bfseries Imageamento Tomográfico Subsuperficial\\[0.4em]
   por Fusão Coerente Monoestática--Biestática\\[0.4em]
   com Veículos Aéreos Não Tripulados\par}
  \vspace{1cm}
  {\Large\itshape Projeto da Trajetória do Receptor Orientado por\\
   Métricas de Imagem e Supressão do Retorno de Superfície\par}
  \vfill
  {\large\textbf{David Atencia}\par}
  \vspace{0.6cm}
  {\normalsize Tese apresentada à Faculdade de Engenharia Elétrica e de
   Computação da Universidade Estadual de Campinas como parte dos requisitos
   exigidos para a obtenção do título de Doutor em Engenharia Elétrica.\par}
  \vspace{0.4cm}
  {\normalsize Orientador: Prof.\ Dr.\ Hugo E.\ Hernández-Figueroa\par}
  \vspace{0.4cm}
  {\normalsize Campinas, 2026\par}
\end{titlepage}

%%─── Dedicatória ──────────────────────────────────────────────────────────────
\cleardoublepage
\thispagestyle{empty}
\vspace*{\fill}
\begin{flushright}
  % TODO: Escrever dedicatória.
  \textit{[Dedicatória a ser escrita.]}
\end{flushright}
\vspace*{\fill}

%%─── Agradecimentos ───────────────────────────────────────────────────────────
\chapter*{Agradecimentos}
\addcontentsline{toc}{chapter}{Agradecimentos}

% TODO: Escrever agradecimentos (orientador, banca, agências de fomento,
%       colaboradores, família).

\textit{[Agradecimentos a serem escritos.]}

%%─── Resumo ───────────────────────────────────────────────────────────────────
\chapter*{Resumo}
\addcontentsline{toc}{chapter}{Resumo}

Esta tese estuda o imageamento tomográfico tridimensional de alvos enterrados
com um radar de abertura sintética (SAR) embarcado em dois veículos aéreos não
tripulados: um transceptor, que transmite e adquire continuamente o canal
monoestático, e um receptor remoto, que adquire simultaneamente o canal
biestático do mesmo pulso. A pergunta central é quando, onde e por que o canal
biestático acrescenta informação à imagem monoestática que o sistema já produz,
e como a trajetória do receptor deve ser projetada para maximizar essa
contribuição.

A investigação partiu da hipótese de que posicionar o receptor de modo que o
raio refratado incida na interface ar--solo no ângulo de Brewster maximizaria a
potência recebida e, por consequência, a qualidade da imagem. Essa hipótese foi
refutada. Demonstra-se que a transmitância de Brewster acrescenta cerca de
$1\dB$ ao enlace, valor dominado pela perda geométrica de propagação, e que a
potência média recebida, quando promediada sobre a orientação desconhecida do
alvo, é estatisticamente a mesma em qualquer posição biestática e na
configuração monoestática. O ângulo de pseudo-Brewster de um solo com perdas é
um mínimo largo, porém raso, que só suprime a reflexão especular coerente, e
esta só alcança o receptor na geometria diametral, inútil para imageamento.

O ganho do canal biestático é geométrico e aparece na reconstrução. Derivam-se
expressões fechadas da resolução tridimensional para geometrias Tx--Rx
arbitrárias com refração, que reproduzem a imagem com erro inferior a $1,6\,\%$,
e prova-se que um receptor à mesma altura do transmissor nunca melhora a
resolução a $-3\dB$. Em contrapartida, separações azimutais entre $60^\circ$ e
$150^\circ$ preenchem o disco horizontal do espaço de números de onda: a fusão
coerente normalizada reduz a razão integrada de lóbulos laterais em até
$7,7\dB$, o pico de lóbulo lateral em até $5,2\dB$ e eleva o contraste em até
$5,5\dB$, à custa de um lóbulo principal $16$ a $30\,\%$ mais largo. Com alvos
condutores anisotrópicos e orientação aleatória, a melhoria de contraste
(${+}1,8$ a ${+}2,7\dB$ em $90$ a $98\,\%$ das orientações) só se mantém com um
receptor de dupla polarização alinhado ao campo do transmissor; um receptor TM
co-polar perde cerca de $20\dB$ a $90^\circ$ por descasamento geométrico de
polarização. Identifica-se, ainda, um nulo biestático do espalhamento difuso
em polarização vertical, cuja posição se obtém em forma fechada
($\Delta\phi^\ast\approx140^\circ$) e que eleva a relação alvo--clutter de
alvos reais em $8$ a $20\dB$ quando o canal biestático é usado isoladamente.
Por fim, estabelecem-se os requisitos de sincronização: a fusão coerente exige
a calibração de uma fase entre canais com erro inferior a $30^\circ$, e uma
deriva de relógio superior a $1\,$ns ao longo da abertura anula o ganho.

Propõem-se dois modos de operação do receptor --- separação varrida de
$60^\circ$ a $150^\circ$ com fusão coerente, para imageamento, e separação fixa
próxima de $140^\circ$ em polarização TM, para detecção limitada por clutter
---, e um protocolo de validação experimental com dados reais.

\vspace{1em}
\noindent\textbf{Palavras-chave:} SAR biestático; SAR multiestático;
tomografia SAR; radar de penetração no solo; fusão coerente; espaço de números
de onda; ângulo de pseudo-Brewster; espalhamento por superfícies rugosas;
veículos aéreos não tripulados.

%%─── Abstract ─────────────────────────────────────────────────────────────────
\chapter*{Abstract}
\addcontentsline{toc}{chapter}{Abstract}

This thesis studies three-dimensional tomographic imaging of buried targets
with a synthetic aperture radar (SAR) carried by two unmanned aerial vehicles:
a transceiver, which transmits and continuously acquires the monostatic
channel, and a remote receiver, which simultaneously records the bistatic
channel of the same pulse. The central question is when, where and why the
bistatic channel adds information to the monostatic image that the system
already produces, and how the receiver trajectory should be designed to
maximize that contribution.

The investigation started from the hypothesis that placing the receiver so
that the refracted ray meets the air--soil interface at the Brewster angle
would maximize the received power and hence the image quality. This hypothesis
is refuted. The Brewster transmittance adds about $1\dB$ to the link, an amount
dominated by geometric spreading loss, and the mean received power, averaged
over the unknown target orientation, is statistically the same at every
bistatic position and in the monostatic configuration. The pseudo-Brewster
angle of a lossy soil is a broad but shallow minimum that only suppresses the
coherent specular reflection, which reaches the receiver only in the
diametral geometry, useless for imaging.

The gain of the bistatic channel is geometric and appears in the
reconstruction. Closed-form expressions of the 3D resolution are derived for
arbitrary Tx--Rx geometries with refraction; they reproduce the image within
$1.6\,\%$, and they prove that a receiver at the transmitter altitude never
improves the $-3\dB$ resolution. Azimuthal separations between $60^\circ$ and
$150^\circ$ instead fill the horizontal disk of the wavenumber domain: the
normalized coherent fusion lowers the integrated sidelobe ratio by up to
$7.7\dB$ and the peak sidelobe ratio by up to $5.2\dB$, and raises the contrast
by up to $5.5\dB$, at the cost of a $16$--$30\,\%$ wider main lobe. With
anisotropic conducting targets in random orientation, the contrast gain
($+1.8$ to $+2.7\dB$ in $90$--$98\,\%$ of the orientations) holds only with a
dual-polarized receiver aligned with the transmitter field; a co-polar TM
receiver loses about $20\dB$ at $90^\circ$ through a geometric polarization
mismatch. A bistatic null of the vertically polarized diffuse surface return
is also identified, located in closed form ($\Delta\phi^\ast\approx140^\circ$);
used as a stand-alone channel it raises the signal-to-clutter ratio of real
targets by $8$--$20\dB$. Finally, synchronization requirements are
established: coherent fusion needs one calibrated inter-channel phase with
error below $30^\circ$, and a clock drift above $1\,$ns over the aperture
cancels the gain.

Two receiver operating modes are proposed --- a separation swept from
$60^\circ$ to $150^\circ$ with coherent fusion for imaging, and a fixed
separation near $140^\circ$ in TM polarization for clutter-limited detection
--- together with a protocol for experimental validation with real data.

\vspace{1em}
\noindent\textbf{Keywords:} bistatic SAR; multistatic SAR; SAR tomography;
ground penetrating radar; coherent fusion; wavenumber domain; pseudo-Brewster
angle; rough-surface scattering; unmanned aerial vehicles.

%%─── Listas ───────────────────────────────────────────────────────────────────
\tableofcontents
\listoffigures
\listoftables

%%─── Lista de Acrônimos ───────────────────────────────────────────────────────
\chapter*{Lista de Acrônimos}
\addcontentsline{toc}{chapter}{Lista de Acrônimos}

{\small
\begin{longtable}{ll}
\toprule
Acrônimo & Descrição \\
\midrule
\endhead
ABCD   & Matriz de parâmetros de transmissão (transferência) \\
AIEM, IEM & \textit{(Advanced) Integral Equation Model} \\
BP     & \textit{Backprojection} (retroprojeção) \\
GNSS   & \textit{Global Navigation Satellite System} \\
GPR    & \textit{Ground Penetrating Radar} (radar de penetração no solo) \\
GPSDO  & Oscilador disciplinado por GPS \\
HH, VV & Polarização horizontal (TE) / vertical (TM) em transmissão e recepção \\
ISLR   & \textit{Integrated Sidelobe Ratio} (razão integrada de lóbulos laterais) \\
OCXO   & Oscilador a cristal com controle de temperatura \\
PEC    & Condutor elétrico perfeito \\
PO     & Óptica Física (\textit{Physical Optics}) \\
PSF    & \textit{Point Spread Function} (função de espalhamento de ponto) \\
PSLR   & \textit{Peak Sidelobe Ratio} (razão de pico de lóbulo lateral) \\
RCS    & \textit{Radar Cross Section} (seção reta radar) \\
RTK    & \textit{Real-Time Kinematic} (posicionamento GNSS centimétrico) \\
SAR    & \textit{Synthetic Aperture Radar} (radar de abertura sintética) \\
SCR    & \textit{Signal-to-Clutter Ratio} (relação alvo--clutter) \\
SNR    & \textit{Signal-to-Noise Ratio} (relação sinal--ruído) \\
SPM    & \textit{Small Perturbation Method} (método das pequenas perturbações) \\
TE     & Transverso-elétrico \\
TM     & Transverso-magnético \\
VANT   & Veículo aéreo não tripulado \\
\bottomrule
\end{longtable}}

%%─── Lista de Símbolos ────────────────────────────────────────────────────────
\chapter*{Lista de Símbolos}
\addcontentsline{toc}{chapter}{Lista de Símbolos}

{\small
\begin{longtable}{lll}
\toprule
Símbolo & Descrição & Unidade \\
\midrule
\endhead
$c$                                       & Velocidade da luz no vácuo                         & m/s       \\
$f_c,\ \lambda_0 = c/f_c$                 & Frequência portadora e comprimento de onda no ar    & Hz, m     \\
$B$                                       & Largura de banda                                    & Hz        \\
$k_0 = 2\pi f/c$                          & Número de onda no vácuo                             & rad/m     \\
$\varepsilon_r = \varepsilon'-\jj\varepsilon''$ & Permissividade relativa complexa do solo      & --        \\
$n = n'-\jj n''$                          & Índice de refração complexo do solo                 & --        \\
$r_{\text{TM}},\ r_{\text{TE}}$           & Coeficientes de reflexão de Fresnel                  & --        \\
$R_{\text{TM}},\ T_{\text{TM}}$           & Reflectância e transmitância TM (potência)          & --        \\
$\theta_B,\ \theta_{pB}$                  & Ângulos de Brewster e de pseudo-Brewster            & rad       \\
$\mathcal{L}$                             & Caminho óptico                                      & m         \\
$p$                                       & Invariante de Snell (parâmetro de raio)             & --        \\
$\mathbf{M}_i,\ Z_i$                      & Matriz ABCD e impedância TM normalizada da camada $i$ & --      \\
$\Gamma_{\text{TM}},\ \tau_{\text{TM}}$   & Coeficientes globais de reflexão e transmissão      & --        \\
$\rho,\ z$                                & Raio horizontal e altura da plataforma              & m         \\
$\psi$                                    & Ângulo de incidência no ar visto do alvo            & rad       \\
$\theta$                                  & Ângulo de refração no solo                          & rad       \\
$\Delta\phi$                              & Separação azimutal Tx--Rx                           & rad       \\
$\kappa$                                  & Razão entre os raios horizontais do Rx e do Tx      & --        \\
$\hat{\vect u}_T,\ \hat{\vect u}_R$       & Versores, no solo, do alvo para Tx e Rx             & --        \\
$\vect K$                                 & Vetor de número de onda biestático                  & rad/m     \\
$S$                                       & Fator de abertura horizontal biestática             & --        \\
$q$                                       & Fator de abertura vertical refratada                & --        \\
$x_{J_0} = 1,1264$                        & Argumento de $-3\dB$ da função de Bessel $J_0$      & --        \\
$\delta_x,\ \delta_y,\ \delta_z$          & Resoluções a $-3\dB$ da PSF em amplitude            & m         \\
$V_{-3}$                                  & Volume do lóbulo principal a $-3\dB$                & m$^3$     \\
$\sigma,\ \sigma^0_{pp}$                  & RCS do alvo e coeficiente de espalhamento da superfície & m$^2$, -- \\
$s,\ \ell$                                & Rugosidade rms e comprimento de correlação          & m         \\
$\Delta\phi^\ast$                         & Separação do nulo biestático VV                     & rad       \\
$I(\vect r)$                              & Imagem complexa reconstruída                        & --        \\
$\varphi_0$                               & Fase residual entre canais                          & rad       \\
$J$                                       & Função de custo                                     & --        \\
\bottomrule
\end{longtable}}

%%═══════════════════════════════════════════════════════════════════════════════
\mainmatter
%%═══════════════════════════════════════════════════════════════════════════════

%%═══════════════════════════════════════════════════════════════════════════════
\chapter{Introdução}
\label{cap:intro}
%%═══════════════════════════════════════════════════════════════════════════════

\section{Contexto e Motivação}
\label{sec:intro_contexto}

O radar de abertura sintética (SAR) é uma modalidade de radar coerente que
processa os ecos de múltiplas posições sequenciais da antena para sintetizar
uma abertura virtual muito superior à antena física, produzindo imagens de alta
resolução independentes das condições atmosféricas e da iluminação
solar~\citar{moreira2013}, \citar{cumming2005}, \citar{elachi2006}. Quando a plataforma descreve trajetórias circulares
ou helicoidais ao redor da cena, o sistema adquire projeções em todo o
intervalo azimutal e numa faixa de elevações, o que permite reconstruir a
distribuição tridimensional de refletividade: é a tomografia
SAR~\citar{knaell1995}, \citar{reigber2000}, \citar{zhu2010}, \citar{ponce2014}, \citar{goes2022}.
A tomografia biestática foi demonstrada, até aqui, para reconstrução de
superfície e de deformação urbana~\citar{ge2019}, \citar{chenz2024}, \citar{he2025},
e não para sensoriamento subsuperficial.
No regime subsuperficial --- detecção de dutos, cavidades, artefatos
enterrados e estruturas estratigráficas --- a onda atravessa a interface
ar--solo, propaga-se num dielétrico com perdas e retorna, de modo que a
refração, a atenuação e o retorno da própria superfície passam a limitar o
desempenho do sistema~\citar{garcia2019}, \citar{daniels2004}.

Na configuração \emph{biestática}, transmissor e receptor ocupam plataformas
distintas. A literatura não estabelece uma superioridade geral do SAR
biestático sobre o monoestático: a separação entre as plataformas cria graus
de liberdade geométricos que podem ser explorados para finalidades
específicas --- diversidade angular, linhas de base interferométricas ou
tomográficas, preenchimento de regiões adicionais do espaço de frequências
espaciais e supressão de determinados mecanismos de reflexão ---, mas também
introduz sincronização, erros de fase, resolução variante no espaço e
geometrias degeneradas~\citar{horne2002}, \citar{willis1995}, \citar{krieger2006}, \citar{wu2025}.
A miniaturização dos radares e a disponibilidade de VANTs com posicionamento
centimétrico tornaram viáveis sistemas multiestáticos coerentes de baixo
custo~\citar{grathwohl2023}, \citar{kanz2025}, o que coloca em primeiro plano a
pergunta de \emph{como} posicionar as plataformas.

\section{A Evolução da Pergunta de Pesquisa}
\label{sec:intro_evolucao}

Esta tese começou com uma hipótese eletromagnética simples. Em polarização TM
existe um ângulo de incidência --- o ângulo de Brewster $\theta_B$ --- no qual
a reflexão numa interface dielétrica sem perdas se anula e toda a energia é
transmitida~\citar{bornwolf1999}. Se o receptor fosse posicionado, a cada
instante, de modo que o raio refratado alvo$\to$Rx incidisse na interface em
$\theta_B$, a potência recebida seria máxima; formulou-se então a otimização
da trajetória do receptor como um compromisso entre essa potência e a
resolução tridimensional da imagem. Uma revisão bibliográfica dirigida não
encontrou nenhum trabalho que combinasse tomografia SAR biestática, refração
na interface e otimização da geometria guiada pelo ângulo de Brewster.

A análise quantitativa refutou essa hipótese em três etapas
(Capítulo~\ref{cap:potencia}). Primeiro, a transmitância de Brewster
acrescenta cerca de $1\dB$ ao enlace de ida e volta, ao passo que a perda
geométrica varia dezenas de decibéis com a posição do receptor: um otimizador
de potência aproxima o receptor do alvo em vez de buscar Brewster. Segundo,
para um alvo sobre o eixo das hélices, a potência é rigorosamente invariante à
separação azimutal entre transmissor e receptor. Terceiro, quando se
considera a orientação desconhecida de alvos anisotrópicos, a potência média
recebida é estatisticamente a mesma em qualquer posição biestática e na
configuração monoestática. A posição do receptor não pode, portanto, ser
escolhida pela potência.

A pergunta foi então reformulada a partir de uma constatação de arquitetura:
o transmissor é um \emph{transceptor}, que recebe o próprio eco em todos os
pulsos. O sistema sempre produz uma imagem monoestática, e o receptor remoto
só se justifica se acrescentar a ela informação que o canal monoestático não
possui. A questão central desta tese passa a ser:

\begin{quote}
\emph{Quando, onde e por que um canal biestático simultâneo melhora a imagem
tomográfica monoestática de alvos enterrados, e qual trajetória do receptor
maximiza essa melhoria sob restrições de voo, de polarização e de
sincronização?}
\end{quote}

O ângulo de Brewster não desaparece da análise, mas muda de papel: deixa de
ser um critério de potência e passa a ser investigado como mecanismo de
\emph{supressão do retorno da superfície}, cuja utilidade se mede pela relação
alvo--clutter e não pelo eco do alvo.

\section{O Sistema de Sensoriamento}
\label{sec:intro_sistema}

O sistema considerado é composto por dois VANTs (Fig.~\ref{fig:geom_sistema}).
O primeiro transporta o transceptor, que transmite pulsos de banda larga e
registra o canal monoestático complexo; sua trajetória é uma hélice cônica
fixa, projetada de antemão. O segundo transporta o receptor remoto, que
registra o canal biestático do mesmo pulso. Os dois canais são coerentes entre
si, de modo que podem ser combinados no domínio complexo após a formação de
imagem por retroprojeção com refração. O produto científico é a comparação
entre três imagens: monoestática, biestática e fundida.

\begin{figure}[H]
\centering
\includegraphics[width=\textwidth]{A_geometria.pdf}
\caption{Geometria do sistema: transceptor (canal monoestático) em hélice
cônica, receptor remoto na configuração diametral ($\Delta\phi=180^\circ$) e
numa posição intermediária ($\Delta\phi=90^\circ$). À direita, o perfil radial
comum das hélices e as linhas de incidência mínima, de pseudo-Brewster do solo
de referência e máxima.}
\label{fig:geom_sistema}
\end{figure}

A interface ar--solo é plana e horizontal; o solo é um dielétrico homogêneo
com permissividade complexa, e a extensão a $N$ camadas é tratada no
Capítulo~\ref{cap:multicamadas}. A superfície tem rugosidade moderada e
contribui com uma componente especular coerente e uma componente difusa.

\section{Objetivos}
\label{sec:intro_objetivos}

O \textbf{objetivo geral} desta tese é estabelecer, com rigor analítico e por
simulação eletromagnética coerente, as condições geométricas,
eletromagnéticas e de sincronização em que um canal biestático simultâneo
melhora a imagem tomográfica monoestática de alvos enterrados, e derivar dessas
condições critérios de projeto da trajetória do receptor.

Os \textbf{objetivos específicos} são:

\begin{enumerate}[label=\textbf{O\arabic*.},leftmargin=*]
  \item Derivar expressões fechadas da resolução tridimensional de um SAR com
  refração na interface para geometrias Tx--Rx arbitrárias e alvos dentro ou
  fora do eixo da trajetória, e estabelecer os limites em que o canal
  biestático pode resolver melhor que o monoestático.
  \item Reavaliar quantitativamente o critério de potência recebida e o papel
  do ângulo de Brewster no posicionamento do receptor, considerando o enlace
  completo, a geometria real e a orientação desconhecida do alvo.
  \item Construir um modelo de simulação coerente que gere, para o mesmo pulso,
  os canais monoestático e biestático com refração, solo com perdas,
  superfície rugosa e alvos anisotrópicos, e definir uma fusão normalizada que
  separe o ganho geométrico do ganho por número de amostras.
  \item Caracterizar como a separação Tx--Rx altera o suporte no espaço de
  números de onda, a PSF e as métricas de imagem, e identificar as separações
  de maior complementaridade.
  \item Avaliar a robustez da melhoria à anisotropia e à orientação do alvo e
  determinar a regra de polarização do receptor.
  \item Avaliar se a operação próxima do pseudo-Brewster suprime o retorno da
  superfície de forma robusta e qual é o seu custo em qualidade de imagem, e
  identificar mecanismos geométricos alternativos de supressão de clutter.
  \item Projetar a trajetória do receptor por uma função de custo baseada em
  métricas de imagem, de clutter e de voo, e quantificar os requisitos de
  sincronização e calibração entre canais.
  \item Definir o protocolo de validação experimental com dados reais.
\end{enumerate}

\section{Contribuições}
\label{sec:intro_contrib}

As contribuições originais deste trabalho são:

\begin{enumerate}
  \item \textbf{Resolução 3D biestática com refração em forma fechada}
  (Capítulo~\ref{cap:resolucao}). As larguras $\delta_x$, $\delta_y$ e
  $\delta_z$ são expressas em termos de dois fatores de abertura, $S$ e $q$,
  calculados a partir da geometria vista do alvo; as expressões reproduzem a
  imagem com erro máximo de $1,6\,\%$ em 34 geometrias. Prova-se que um
  receptor à mesma altura do transmissor nunca melhora $\delta_{xy}$ nem
  $\delta_z$ e que o ganho vertical máximo de qualquer geometria biestática
  está limitado a $4,6\,\%$ pela refração.

  \item \textbf{Refutação documentada do critério de Brewster para o
  posicionamento do receptor} (Capítulo~\ref{cap:potencia}). Mostra-se, por
  três vias independentes, que a potência recebida não discrimina posições do
  receptor e que o pseudo-Brewster só atua sobre a reflexão especular da
  geometria diametral.

  \item \textbf{Fusão coerente mono+biestática normalizada e mecanismo do
  ganho} (Capítulos~\ref{cap:modelo} e~\ref{cap:geometria}). A fusão a
  orçamento constante de amostras isola o ganho geométrico, que se explica
  pelo preenchimento do disco horizontal do espaço de números de onda por
  anéis de raio $|\cos(\Delta\phi/2)|$ vezes o monoestático.

  \item \textbf{Regra de polarização do receptor remoto}
  (Capítulo~\ref{cap:alvos}). Um receptor TM co-polar sofre um descasamento
  geométrico de cerca de $20\dB$ a $\Delta\phi=90^\circ$ com alvos reais,
  invisível para o alvo pontual escalar; um receptor de dupla polarização
  alinhado recupera de $13$ a $28\dB$ e torna a melhoria de contraste robusta
  à orientação.

  \item \textbf{Nulo biestático do clutter difuso em polarização vertical}
  (Capítulo~\ref{cap:superficie}). Deriva-se a curva fechada
  $\cos\Delta\phi^\ast = -\varepsilon\sin^2\theta/(\varepsilon-\sin^2\theta)$,
  que generaliza o ângulo de Brewster ao espalhamento biestático difuso, e
  demonstra-se que esse nulo, ao contrário do nulo HH a $90^\circ$, sobrevive
  com alvos reais e eleva a relação alvo--clutter em $8$ a $20\dB$.

  \item \textbf{Projeto da trajetória do receptor por métricas de imagem e
  requisitos de sincronização} (Capítulo~\ref{cap:trajetoria}). Uma fronteira
  de Pareto entre qualidade de imagem, clutter e custo de voo conduz a uma
  trajetória de separação varrida com custo cinemático quase nulo, e
  quantificam-se as tolerâncias de fase, relógio e posição.
\end{enumerate}

\section{Organização da Tese}
\label{sec:intro_org}

O \textbf{Capítulo~\ref{cap:fundamentos}} reúne os fundamentos
eletromagnéticos e geométricos: coeficientes de Fresnel em meios com perdas,
ângulos de Brewster e de pseudo-Brewster, traçado de raios pelo princípio de
Fermat, geometria helicoidal, modelo de sinal e retroprojeção com refração. O
\textbf{Capítulo~\ref{cap:resolucao}} deriva a resolução tridimensional em
forma fechada e a valida. O \textbf{Capítulo~\ref{cap:multicamadas}} estende o
modelo de transmissão a meios estratificados pelo formalismo ABCD. O
\textbf{Capítulo~\ref{cap:potencia}} formula a hipótese original de
otimização por potência e documenta sua refutação. O
\textbf{Capítulo~\ref{cap:modelo}} descreve o modelo de simulação coerente, a
fusão normalizada e as métricas. Os \textbf{Capítulos~\ref{cap:geometria}
a~\ref{cap:trajetoria}} apresentam os resultados: complementaridade
geométrica, alvos reais e polarização, supressão do retorno de superfície e
projeto da trajetória com requisitos de sincronização. O
\textbf{Capítulo~\ref{cap:dados_reais}} estabelece o protocolo e reserva o
espaço para a avaliação com dados reais. O
\textbf{Capítulo~\ref{cap:conclusoes}} responde às perguntas de pesquisa,
lista os resultados negativos e aponta as perspectivas. Os apêndices reúnem
derivações matemáticas e os parâmetros das simulações.

\section{Publicações}
\label{sec:intro_pub}

% TODO: Listar artigos publicados e submetidos originados desta tese.

\textit{[Lista de publicações a ser inserida.]}

%%═══════════════════════════════════════════════════════════════════════════════
\chapter{Fundamentos Teóricos}
\label{cap:fundamentos}
%%═══════════════════════════════════════════════════════════════════════════════

\section{Radar de Abertura Sintética: Princípios e Histórico}
\label{sec:fund_sar_intro}

O SAR é um sistema de imageamento baseado em radar coerente que explora o
movimento da plataforma portadora para sintetizar uma antena virtual de grande
abertura. A técnica, concebida por Carl Wiley em 1951 e demonstrada
experimentalmente por Sherwin et al.\ em 1953~\citar{moreira2013}, produz imagens
de alta resolução em azimute independentemente do comprimento físico da antena.
A resolução em azimute de um SAR monoestático convencional é
$\delta_{az} \approx D/2$, onde $D$ é o comprimento da antena real: a resolução
\emph{melhora} com antenas menores, pois a abertura sintética é inversamente
proporcional a~$D$.

Munson, O'Brien e Jenkins~\citar{munson1983} mostraram que o processamento SAR
em modo \textit{spotlight} é equivalente à tomografia computadorizada: cada
posição do radar mede uma fatia do espectro espacial da cena, e a imagem é a
transformada inversa da cobertura acumulada. Arikan e Munson~\citar{arikan1988}
estenderam a formulação ao caso biestático, no qual a fatia é determinada pela
bissetriz das direções de iluminação e de observação. Essa interpretação ---
a imagem como inversão do suporte no espaço de números de onda --- é o fio
condutor de toda a análise desta tese.

Na configuração biestática, transmissor e receptor encontram-se em plataformas
distintas. Esta separação amplia o espaço de geometrias disponíveis: ângulos de
iluminação e recepção independentes permitem observar componentes da resposta
do alvo inacessíveis ao radar monoestático~\citar{horne2002}, \citar{willis1995}.
No regime subsuperficial, a onda deve atravessar a interface ar--solo, o que
introduz a refração, a atenuação no meio e o retorno da superfície como fatores
limitantes adicionais~\citar{garcia2019}, \citar{daniels2004}.

\section{Propagação Eletromagnética em Meios Dielétricos}
\label{sec:fund_em}

\subsection{Equações de Maxwell em Meios Lineares Isotrópicos}
\label{ssec:maxwell}

Considerem-se dois meios lineares, isotrópicos, homogêneos e não magnéticos,
separados por uma interface plana em $z = 0$. O meio superior ($z > 0$) é o ar,
com $\varepsilon_1 = \varepsilon_0$ e índice $n_1 = 1$. O meio inferior
($z < 0$) é o solo, com permissividade relativa
$\varepsilon_r = \varepsilon' - \jj\varepsilon''$ e índice de refração
$n_2 = \sqrt{\varepsilon_r} = n' - \jj n''$. O caso sem perdas corresponde a
$\varepsilon''=0$.

As equações de Maxwell no domínio da frequência (convenção $e^{\jj\omega t}$)
são:
\begin{align}
  \nabla \times \vect{E} &= -\jj\omega\mu_0\,\vect{H}, \label{eq:maxwell_rot_e} &
  \nabla \times \vect{H} &= +\jj\omega\varepsilon\,\vect{E}, \\
  \nabla \cdot (\varepsilon\vect{E}) &= 0, \label{eq:maxwell_div_e} &
  \nabla \cdot \vect{H} &= 0 .
\end{align}
Tomando o rotacional da primeira equação e usando a identidade
$\nabla\times(\nabla\times\vect{E}) = \nabla(\nabla\cdot\vect{E}) - \nabla^2\vect{E}$,
obtém-se a equação de Helmholtz
\begin{equation}
  \nabla^2 \vect{E} + k^2 \vect{E} = 0, \qquad k = k_0 n = \frac{\omega n}{c}.
  \label{eq:helmholtz}
\end{equation}
Num meio com perdas, $k$ é complexo; a parte imaginária produz a atenuação
exponencial do campo, $e^{-\alpha\ell}$ com $\alpha = k_0 n''$, ao longo de um
percurso $\ell$ no solo.

\subsection{Condições de Contorno e Coeficientes de Fresnel}
\label{ssec:fresnel}

Uma onda plana incide no plano $xz$ com ângulo $\theta_i$ em relação à normal.
Na polarização TM (transverso-magnética, ``vertical''), o campo magnético é
perpendicular ao plano de incidência e o campo elétrico pertence a ele; na
polarização TE (transverso-elétrica, ``horizontal''), ocorre o inverso. A
continuidade das componentes tangenciais de $\vect E$ e $\vect H$ em $z=0$
impõe que a fase ao longo de $x$ seja a mesma para as ondas incidente, refletida
e transmitida, donde $\theta_r=\theta_i$ e a \textbf{lei de Snell}
\begin{equation}
  \boxed{n_1\sin\theta_i = n_2\sin\theta_t}.
  \label{eq:snell}
\end{equation}
As mesmas condições, aplicadas às amplitudes, fornecem os
\textbf{coeficientes de Fresnel}:
\begin{align}
  r_{\text{TM}} &= \frac{n_2\cos\theta_i - n_1\cos\theta_t}{n_2\cos\theta_i + n_1\cos\theta_t}, &
  t_{\text{TM}} &= \frac{2n_1\cos\theta_i}{n_2\cos\theta_i + n_1\cos\theta_t}, \label{eq:r_TM}\\
  r_{\text{TE}} &= \frac{n_1\cos\theta_i - n_2\cos\theta_t}{n_1\cos\theta_i + n_2\cos\theta_t}, &
  t_{\text{TE}} &= \frac{2n_1\cos\theta_i}{n_1\cos\theta_i + n_2\cos\theta_t}. \label{eq:r_TE}
\end{align}
Com $n_1=1$, é conveniente escrever $n_2\cos\theta_t=\sqrt{\varepsilon_r-\sin^2\theta_i}$,
expressão válida também para $\varepsilon_r$ complexo, em que $\cos\theta_t$ é
complexo e $\theta_t$ deixa de ter interpretação geométrica direta.

\subsection{Reflectância e Transmitância de Potência}
\label{ssec:pot_fresnel}

O vetor de Poynting médio $\vect{S} = \frac{1}{2}\,\text{Re}(\vect{E}\times\vect{H}^*)$
fornece o fluxo de potência por unidade de área. Para meios sem perdas, a
reflectância e a transmitância normalizadas pelo fluxo incidente são
\begin{equation}
  R_{\text{TM}} = |r_{\text{TM}}|^2, \qquad
  T_{\text{TM}} = \frac{n_2\cos\theta_t}{n_1\cos\theta_i}\,|t_{\text{TM}}|^2 = 1 - |r_{\text{TM}}|^2,
  \label{eq:T_TM}
\end{equation}
e a igualdade $R_{\text{TM}} + T_{\text{TM}} = 1$ é verificada algebricamente no
Apêndice~\ref{app:energia_fresnel}. Neste trabalho, a amplitude transmitida em
cada travessia da interface é representada por $\tau = \sqrt{1-|r|^2}$, que é a
raiz da transmitância de potência e incorpora o fator de impedância.

\begin{figure}[H]
\centering
\includegraphics[width=0.85\textwidth]{pv_fig01_fresnel_tm.pdf}
\caption{Reflectância $R_{\text{TM}}$ e transmitância $T_{\text{TM}}$ em função
do ângulo de incidência para $n_1=1$ e $n_2=2$ ($\varepsilon_r=4$, solo seco
sem perdas). O ângulo de Brewster $\theta_B \approx 63,4^\circ$ é o único zero
de $R_{\text{TM}}$.}
\label{fig:fresnel_tm}
\end{figure}

\section{Os Ângulos de Brewster e de Pseudo-Brewster}
\label{sec:brewster}

\begin{teorema}[Ângulo de Brewster]
\label{teo:brewster}
Para incidência TM entre meios sem perdas e não magnéticos, existe um único
ângulo $\theta_B \in (0,\pi/2)$ tal que $r_{\text{TM}}(\theta_B) = 0$, dado por
\begin{equation}
  \boxed{\theta_B = \arctan\!\left(\frac{n_2}{n_1}\right)}.
  \label{eq:brewster}
\end{equation}
\end{teorema}

\begin{proof}
Impondo $r_{\text{TM}} = 0$ em~\eqref{eq:r_TM}: $n_2\cos\theta_B = n_1\cos\theta_t^{(B)}$.
Elevando ao quadrado e usando a lei de Snell
$\sin\theta_t^{(B)} = (n_1/n_2)\sin\theta_B$:
\[
  n_2^2 - n_2^2\sin^2\theta_B = n_1^2 - \frac{n_1^4}{n_2^2}\sin^2\theta_B
  \;\Rightarrow\;
  \sin^2\theta_B = \frac{n_2^2}{n_1^2+n_2^2},
\]
de modo que $\tan\theta_B = n_2/n_1$. A unicidade decorre da monotonicidade de
$\sin^2\theta$ em $(0,\pi/2)$.
\end{proof}

\begin{corolario}
\label{cor:brewster_ortho}
No ângulo de Brewster, o raio refratado e o raio refletido seriam
perpendiculares: $\theta_B + \theta_t^{(B)} = \pi/2$.
\end{corolario}

\begin{proof}
$\sin\theta_t^{(B)} = (n_1/n_2)\cdot n_2/\sqrt{n_1^2+n_2^2} = \cos\theta_B$.
\end{proof}

A interpretação física é a da radiação dipolar: os dipolos induzidos no meio
inferior oscilam na direção do hipotético raio refletido, e um dipolo não
irradia ao longo do próprio eixo~\citar{bornwolf1999}. Para $n_2 = 2$:
$\theta_B \approx 63,43^\circ$ no ar e $\theta_t^{(B)} = 26,57^\circ$ no solo.
Na TE não existe zero de reflexão para meios não magnéticos.

O amortecimento da reflexão nas proximidades de $\theta_B$ é observável em
imagens SAR reais: Ouchi e Yang~\citar{ouchi2017}, \citar{ouchi2018} reportam
evidência direta em imagens TerraSAR-X de estruturas de concreto, e
Thirion-Lefevre e Guinvarc'h~\citar{thirion2018} mostram o efeito de duplo
ângulo de Brewster em superfícies dielétricas ortogonais. Em ambos os casos, o
fenômeno se manifesta como \emph{redução} de um mecanismo de reflexão em VV, e
não como aumento da potência de um alvo.

\begin{figure}[H]
\centering
\includegraphics[width=0.80\textwidth]{pv_fig08_brewster_geometry.pdf}
\caption{Geometria do ângulo de Brewster. Em $\theta_B$, a ortogonalidade entre
o raio refratado e o hipotético raio refletido impede a radiação na direção
refletida, anulando $r_{\text{TM}}$.}
\label{fig:brewster_geom}
\end{figure}

\subsection{Meios com Perdas: o Pseudo-Brewster}
\label{ssec:pseudo_brewster}

Quando $\varepsilon''>0$, o numerador de $r_{\text{TM}}$ não se anula para
nenhum ângulo real: a reflectância TM passa por um mínimo não nulo, no ângulo
de \emph{pseudo-Brewster}
\begin{equation}
  \theta_{pB} = \arg\min_{\theta\in[0,\pi/2)} |r_{\text{TM}}(\theta)|^2 .
  \label{eq:pseudo_brewster}
\end{equation}
A Tabela~\ref{tab:solos} apresenta os quatro solos usados nas simulações desta
tese, do caso ideal sem perdas a um solo úmido. Duas propriedades se destacam:
o mínimo torna-se mais raso com o aumento das perdas, e desloca-se para
incidências mais rasantes com o aumento de $\varepsilon'$.

\begin{table}[H]
\centering
\caption{Solos simulados e seus ângulos de pseudo-Brewster.}
\label{tab:solos}
\small
\begin{tabular}{@{}lcccc@{}}
\toprule
Solo & $\varepsilon_r$ & $\theta_{pB}$ & $\min|r_{\text{TM}}|^2$ & Atenuação a 425\,MHz \\
\midrule
S0 (sem perdas) & $4$            & $63,43^\circ$ & $-\infty$     & 0 \\
S1 (seco, referência) & $3,99-\jj0,40$ & $63,45^\circ$ & $-34,5\dB$    & $7,7\dB$/m \\
S2 (úmido)                            & $6,5-\jj0,9$   & $68,66^\circ$ & $-30,7\dB$    & $13,6\dB$/m \\
S3 (muito úmido)                      & $12-\jj2,5$    & $74,05^\circ$ & $-26,5\dB$    & $27,7\dB$/m \\
\bottomrule
\end{tabular}
\end{table}

\begin{figure}[H]
\centering
\includegraphics[width=\textwidth]{C_reflectancia.pdf}
\caption{Reflectâncias TM e TE dos quatro solos (esquerda) e reflectância TM
efetiva em função do erro de incidência e do erro de atitude (rolamento) para
os solos S0 e S1. No solo sem perdas o mínimo é um zero exato, mas estreito; no
solo com perdas, é largo e raso.}
\label{fig:reflectancia_solos}
\end{figure}

A Figura~\ref{fig:reflectancia_solos} antecipa o resultado central do
Capítulo~\ref{cap:superficie}: num solo sem perdas o zero de Brewster é exato,
mas um erro de $0,5^\circ$ eleva a reflectância de $-190$ para $-42\dB$; num
solo com perdas, o mínimo é tolerante a erros de alguns graus, mas só oferece
cerca de $17\dB$ de supressão em relação a uma incidência de $53^\circ$.

\section{Traçado de Raios e Princípio de Fermat}
\label{sec:fermat}

\subsection{Formulação Variacional}
\label{ssec:fermat_var}

O princípio de Fermat estabelece que o raio entre dois pontos torna o caminho
óptico \emph{estacionário}~\citar{bornwolf1999}.

\begin{definicaom}[Caminho Óptico]
\label{def:caminho_optico}
O caminho óptico entre os pontos $A$ e $B$ ao longo de uma trajetória
$\mathcal{C}$ é $\mathcal{L}[\mathcal{C}] = \int_{\mathcal{C}} n'(\vect{r})\,\dd s$.
\end{definicaom}

Para dois meios homogêneos separados pela interface $z=0$, com ponto aéreo
$A$ à altura $z_A>0$ e ponto subsuperficial $B$ à profundidade $|z_B|$, o único
grau de liberdade é o ponto de cruzamento $Q$ da interface, à distância
horizontal $u$ de $A$ no plano vertical que contém $A$ e $B$:
\begin{equation}
  \mathcal{L}(u) = n_1\sqrt{u^2 + z_A^2} + n'_2\sqrt{(d-u)^2 + z_B^2},
  \label{eq:fermat_L}
\end{equation}
onde $d$ é a separação horizontal entre $A$ e $B$. Em meios com perdas, a
geometria do raio é determinada pela parte real do índice, e a parte imaginária
entra apenas como atenuação ao longo do trecho no solo.

\subsection{Estacionaridade, Lei de Snell e Unicidade}
\label{ssec:fermat_snell}

A condição $\dd\mathcal{L}/\dd u = 0$ fornece
\begin{equation}
  g(u) \equiv \frac{n_1\,u}{\sqrt{u^2+z_A^2}} - \frac{n'_2\,(d-u)}{\sqrt{(d-u)^2+z_B^2}} = 0,
  \label{eq:fermat_estac}
\end{equation}
que é a lei de Snell~\eqref{eq:snell}. A derivada
\begin{equation}
  g'(u) = \frac{n_1 z_A^2}{(u^2+z_A^2)^{3/2}}
  + \frac{n'_2 z_B^2}{((d-u)^2+z_B^2)^{3/2}} > 0
  \label{eq:g_prime}
\end{equation}
é estritamente positiva, de modo que $g$ tem raiz única em $[0,d]$ e o ponto
estacionário é o mínimo global de $\mathcal{L}$.

\subsection{Solução Numérica e Caminho Geométrico}
\label{ssec:newton}

A equação $g(u)=0$ não tem solução fechada. A iteração de Newton--Raphson
$u^{(k+1)} = u^{(k)} - g(u^{(k)})/g'(u^{(k)})$ converge quadraticamente, pois
$g'>0$ e $g''$ é limitada; três a cinco iterações bastam para precisão
submilimétrica. O comprimento geométrico do raio,
\begin{equation}
  \ell = \underbrace{\sqrt{(u^*)^2+z_A^2}}_{\ell^{a}} + \underbrace{\sqrt{(d-u^*)^2+z_B^2}}_{\ell^{s}},
  \label{eq:path_geom}
\end{equation}
determina a perda por espalhamento geométrico; o trecho no solo $\ell^{s}$
determina a atenuação; e o caminho óptico $\ell^{a}+n'_2\ell^{s}$ determina o
retardo do eco e, portanto, a posição do alvo na imagem.

Para alvos rasos em comparação com a altura da plataforma, a lei de Snell
admite uma solução explícita com um passo de correção, útil nas expressões
analíticas do Capítulo~\ref{cap:resolucao}. Sendo $D$ a distância horizontal
plataforma--alvo, $h$ a altura da plataforma e $d_P$ a profundidade do alvo,
\begin{equation}
  \psi^{(0)}=\arctan\frac{D}{h},\qquad
  \psi=\arctan\!\left(\frac{D-d_P\,\sin\psi^{(0)}\big/\sqrt{n'^2_2-\sin^2\psi^{(0)}}}{h}\right),
  \label{eq:snell_explicita}
\end{equation}
que coincide com a solução iterativa em quatro algarismos para
$d_P/h\lesssim0,02$.

\section{Geometria Helicoidal}
\label{sec:fund_sar_bist}

O transceptor descreve uma hélice cônica:
\begin{equation}
  \vect p_{Tx}(t)=\big(\rho(t)\cos\phi(t),\ \rho(t)\sin\phi(t),\ z(t)\big),\qquad
  \phi(t)=\frac{2\pi N_v t}{T},
  \label{eq:helix}
\end{equation}
com raio e altura variando monotonamente entre os extremos
$(\rho_{\min},z_{\min})$ e $(\rho_{\max},z_{\max})$, $N_v$ voltas e duração $T$.
Para um alvo próximo do eixo, a incidência média vista do alvo é
$\psi_0 = \arctan(\rho_0/z_0)$, com $\rho_0$ e $z_0$ médios, e a excursão de
incidência ao longo da hélice é o que confere à abertura extensão vertical no
espaço de números de onda. A hélice usada como referência nesta tese tem duas
voltas, raio de $137,4$ a $165,5\,$m e altura de $131,5$ a $93,6\,$m; vista de
um alvo a $2\,$m de profundidade no eixo, a incidência varia de $46,1^\circ$ a
$60,4^\circ$ e a distância ao centro da cena é praticamente constante
($\approx189\,$m), porque a linha de base da hélice é perpendicular à linha de
visada.

O receptor remoto voa, a cada instante, à mesma altura do transmissor:
\begin{equation}
  \vect p_{Rx}(t) = \Big(\kappa(t)\rho(t)\cos[\phi(t)+\Delta\phi(t)],\;
  \kappa(t)\rho(t)\sin[\phi(t)+\Delta\phi(t)],\; z(t)\Big),
  \label{eq:rx_traj}
\end{equation}
onde $\Delta\phi(t)$ é a separação azimutal e $\kappa(t)$ a razão entre raios.
Com $\kappa=1$ e $\Delta\phi$ constante, o receptor percorre exatamente a mesma
hélice do transmissor, rigidamente girada de $\Delta\phi$ em azimute e
sincronizada com ele; a configuração monoestática corresponde a $\Delta\phi=0$,
e a configuração \emph{diametral}, a $\Delta\phi=180^\circ$.

\begin{figure}[H]
\centering
\includegraphics[width=0.86\textwidth]{pv_fig10_sistema_3d.pdf}
\caption{Geometria do SAR biestático helicoidal: transmissor e receptor em
hélices cônicas ao redor da cena subsuperficial, separados por um desfasamento
azimutal $\Delta\phi$.}
\label{fig:sistema3d}
\end{figure}

\section{Modelo de Sinal}
\label{ssec:signal_model}

O sinal transmitido é um pulso de banda $B$ centrado em $f_c$. Após a
compressão em alcance, o eco do pulso $m$ é representado no domínio da
frequência por $N_f$ amostras $f_k\in[f_c-B/2,\,f_c+B/2]$. Para um canal
(Tx, Rx) e um alvo de amplitude complexa de espalhamento $S_{mk}$, o eco é
\begin{equation}
  s_{mk} = \frac{\sqrt{P_t}\,\lambda_k}{(4\pi)^{3/2}}
  \frac{\sqrt{G_t G_r}}{R_T R_R}\,\tau_t\,\tau_r\,
  e^{-\alpha_k(\ell^{s}_T+\ell^{s}_R)}\,S_{mk}\,e^{-\jj2\pi f_k\tau_m},
  \qquad
  \tau_m=\frac{\ell^{a}_T+n'\ell^{s}_T+\ell^{a}_R+n'\ell^{s}_R}{c},
  \label{eq:eco}
\end{equation}
onde os índices $T$ e $R$ designam os trechos transmissor$\to$alvo e
alvo$\to$receptor, $R_{T,R}=\ell^a+\ell^s$ são os comprimentos geométricos,
$\tau_{t,r}$ as amplitudes transmitidas na interface na polarização de cada
trecho, $\alpha_k = k_0 n''$ a atenuação e $G_{t,r}$ os ganhos das antenas. O
canal monoestático é o caso $\vect p_{Rx}=\vect p_{Tx}$ do \emph{mesmo} pulso.
Para um alvo pontual isotrópico, $S_{mk}=\sqrt{\sigma}$; para alvos extensos,
$S_{mk}$ é calculado por Óptica Física (Seção~\ref{sec:mod_alvos}).

\subsection{Equação do Radar Biestático com Refração}
\label{ssec:radar_eq}

O módulo quadrático de~\eqref{eq:eco} é a equação do radar biestático com
refração~\citar{willis1995}, \citar{ulaby1981}:
\begin{equation}
  \boxed{P_r = \frac{P_t\,G_t\,G_r\,\sigma\,T_1\,T_2\,\lambda^2}{(4\pi)^3\,R_T^2\,R_R^2}\,
  e^{-2\alpha(\ell^{s}_T+\ell^{s}_R)}},
  \label{eq:radar_refratado}
\end{equation}
com $T_1=|\tau_t|^2$ e $T_2=|\tau_r|^2$ as transmitâncias de potência das duas
travessias da interface. Esta é a grandeza que a hipótese original desta tese
propunha maximizar (Capítulo~\ref{cap:potencia}).

\section{Retroprojeção com Refração}
\label{ssec:bp_algo}

A imagem é formada por retroprojeção (\textit{backprojection}) no domínio do
tempo~\citar{ding2002}, \citar{zhang2018}, \citar{sun2022}, \citar{zhao2025}: para cada voxel $\vect r$ da grade de
imagem, os ecos são compensados pelo retardo que um alvo em $\vect r$ teria e
somados coerentemente sobre pulsos e frequências,
\begin{equation}
  I(\vect r) = \frac{1}{\sum_m w_m\,N_f}\sum_m w_m\sum_k W_k\,s_{mk}\,
  e^{\,\jj 2\pi f_k \tau_m(\vect r)},
  \label{eq:bp}
\end{equation}
onde $\tau_m(\vect r)$ é o retardo óptico refratado do voxel, calculado pelo
traçado de raios exato da Seção~\ref{sec:fermat} para cada par
(pulso, voxel), $w_m$ é um peso por pulso e $W_k$ uma janela espectral. A
normalização pela soma dos pesos torna $I(\vect r)$ a \emph{média} coerente
das amostras, propriedade usada na fusão normalizada do
Capítulo~\ref{cap:modelo}. O custo é $O(N_p N_f N_v)$ para $N_p$ pulsos e
$N_v$ voxels; na implementação, o perfil em alcance é sobreamostrado $32$
vezes e interpolado, preservando a fase exata do retardo.

A escolha da janela espectral não é inócua. Com espectro retangular, a PSF em
alcance tem lóbulos laterais de $-13\dB$ que decaem lentamente; numa órbita de
altura constante, o reflexo especular da superfície no nadir do canal
monoestático tem retardo constante, e seus lóbulos distantes em alcance se
acumulam coerentemente sobre o eixo da órbita --- justamente onde está o
alvo. Com janela de Hann, esse vazamento cai cerca de $50\dB$ à custa de um
lóbulo principal vertical $1,6$ vez mais largo. Nas análises limitadas por
clutter usa-se, por isso, a janela de Hann; nas análises de forma da PSF,
o espectro retangular.

\section{Conclusões do Capítulo}
\label{sec:fund_concl}

Este capítulo estabeleceu a base sobre a qual os capítulos subsequentes são
construídos. Os coeficientes de Fresnel em meios com perdas mostram que o zero
de Brewster é uma idealização: num solo real há um pseudo-Brewster largo e
raso, e sua utilidade precisa ser medida no sistema completo. O princípio de
Fermat garante a unicidade do ponto de refração e fundamenta o traçado de raios
usado tanto na geração dos ecos quanto na retroprojeção. A geometria do
receptor foi parametrizada por uma separação azimutal e uma razão de raios, à
mesma altura do transmissor, e o modelo de sinal coerente permite gerar, para o
mesmo pulso, os canais monoestático e biestático.

%%═══════════════════════════════════════════════════════════════════════════════
\chapter{Resolução Espacial Tridimensional com Refração}
\label{cap:resolucao}
%%═══════════════════════════════════════════════════════════════════════════════

\section{Introdução e Estado da Arte}
\label{sec:res_arte}

A derivação de fórmulas de resolução para SAR tomográfico tem uma linhagem bem
estabelecida na literatura monoestática. Knaell e Cardillo~\citar{knaell1995}
introduziram a tomografia de radar para geração de imagens tridimensionais;
Reigber e Moreira~\citar{reigber2000} demonstraram a tomografia SAR
aerotransportada multilinha de base em banda L; Ponce et al.~\citar{ponce2014}
estenderam a técnica a SAR circular totalmente polarimétrico; e
Góes~\citar{goes2022} derivou expressões de $\delta_z$ e $\delta_{xy}$ para SAR
circular e espiral monoestático. No caso biestático, Homer et
al.~\citar{homer2005} apresentam um arcabouço geral para a resolução 3D com
transmissor, receptor e alvo não coplanares, e Garza e Qiao~\citar{garza2010}
analisam a PSF de imagens biestáticas reconstruídas por retroprojeção filtrada.
Wang et al.~\citar{wang2020} tratam a focalização com trajetória curva, e
Andre et al.~\citar{andre2025} analisam os artefatos de imagem multiestática 3D
com duas trajetórias independentes.

Nenhum desses trabalhos fornece expressões fechadas para a geometria desta
tese: transmissor em hélice cônica, receptor em posição arbitrária e alvo
enterrado sob uma interface refratora. Este capítulo preenche essa lacuna.
A estratégia é expressar o vetor de número de onda biestático em função de dois
fatores de abertura --- um horizontal, $S$, e um vertical, $q$ --- calculados a
partir da geometria vista do alvo, e obter as larguras a $-3\dB$ da PSF a
partir das estatísticas desses fatores ao longo da abertura.

\section{Vetor de Número de Onda Biestático com Refração}
\label{sec:k_bist}

Seja o alvo em $\vect P=(x_P;\,y_P;\,z_P)$, à profundidade $d_P=-z_P$ num solo de
índice $n=\mathrm{Re}\sqrt{\varepsilon_r}$, e sejam o transmissor e o receptor
em $\vect p_i=(x_i;\,y_i;\,h_i)$, $i\in\{T,R\}$. Para cada pulso $m$ e frequência
$f$, a fase do eco retroprojetado em $\vect P+\Delta\vect r$, linearizada em
$\Delta\vect r$, é $\vect K_{m,f}\cdot\Delta\vect r$, com
\begin{equation}
  \vect K_{m,f}=\frac{2\pi f}{c}\,n\,\big(\hat{\vect u}_{T,m}+\hat{\vect u}_{R,m}\big),
  \qquad
  \chi(\Delta\vect r)=\sum_{m}\sum_{f} e^{\,\jj\vect K_{m,f}\cdot\Delta\vect r},
  \label{eq:k_bist}
\end{equation}
onde $\hat{\vect u}_i$ é o versor, \emph{dentro do solo}, que aponta do alvo
para o ponto de refração da plataforma $i$, e $\chi$ é a PSF em amplitude para
pesos uniformes. O fator $n$ expressa o aumento do número de onda no
dielétrico; a soma dos dois versores é a forma biestática da relação de
Arikan e Munson~\citar{arikan1988}.

\paragraph{Geometria vista do alvo.} Para cada pulso,
\begin{equation}
  \begin{gathered}
  D_i=\sqrt{(x_i-x_P)^2+(y_i-y_P)^2},\qquad
  \phi'_i=\atantwo(y_i-y_P,\;x_i-x_P),\\
  D_i=h_i\tan\psi_i+d_P\tan\theta_i,
  \end{gathered}
  \label{eq:geo_alvo}
\end{equation}
com $\sin\psi_i=n\sin\theta_i$. A incidência no ar $\psi_i$ resulta
de~\eqref{eq:snell_explicita}, e
$\hat{\vect u}_i=(\sin\theta_i\cos\phi'_i;\ \sin\theta_i\sin\phi'_i;\ \cos\theta_i)$.

\section{Resolução Horizontal}
\label{sec:delta_xy}

Como $n\sin\theta_i=\sin\psi_i$, a componente horizontal
de~\eqref{eq:k_bist} \emph{não depende do solo}:
\begin{equation}
  \vect K_h=k_0\big[\sin\psi_T\,\hat{\vect e}(\phi'_T)+\sin\psi_R\,\hat{\vect e}(\phi'_R)\big],
  \qquad
  |\vect K_h|=2k_0\,S,
\end{equation}
\begin{equation}
  \boxed{S=\tfrac12\sqrt{\sin^2\psi_T+\sin^2\psi_R+2\sin\psi_T\sin\psi_R\cos(\phi'_R-\phi'_T)}}.
  \label{eq:S}
\end{equation}
No monoestático, $S=\sin\psi_T$; com $\psi_T=\psi_R=\psi$, $S=\sin\psi\,|\cos(\Delta\phi'/2)|$.

Se a abertura percorre $360^\circ$ em torno do alvo, o suporte horizontal é um
anel de raio $R_c=2k_0\langle S\rangle$ e a PSF horizontal é
$\chi(\rho)\propto J_0(R_c\rho)$ (Apêndice~\ref{app:bessel}). A largura a
$-3\dB$ resulta de $J_0(x_{J_0})=1/\sqrt2$, com $x_{J_0}=1,1264$:
\begin{equation}
  \boxed{\;\delta_x=\delta_y=\frac{2x_{J_0}}{R_c}=\frac{x_{J_0}\,\lambda_0}{2\pi\,\langle S\rangle}
  =\frac{0,1793\,\lambda_0}{\langle S\rangle}\;},
  \label{eq:delta_xy}
\end{equation}
onde $\langle\cdot\rangle$ é a média sobre os pulsos da abertura e
$\lambda_0=c/f_c$.

\begin{observacao}[Constante da largura de Bessel]
É comum usar $u_{-3}\approx1,20$ como argumento de meia potência de $J_0$.
Como $J_0(1,20)=0,67$, isto é, $-3,5\dB$, essa constante superestima
$\delta_{xy}$ em $6,5\,\%$. O valor correto para a largura a $-3\dB$ da PSF em
amplitude é $x_{J_0}=1,1264$, adotado em toda esta tese.
\end{observacao}

\begin{figure}[H]
\centering
\includegraphics[width=0.95\textwidth]{fig07_northoffset_kspace.pdf}
\caption{Cobertura do espaço de números de onda horizontal para três
separações com transmissor e receptor na mesma hélice. (a) $\Delta\phi=0^\circ$:
anel de raio máximo. (b) $\Delta\phi=90^\circ$: raio reduzido por $1/\sqrt{2}$.
(c) $\Delta\phi=180^\circ$: o anel colapsa na origem e não há resolução
horizontal para o alvo no eixo.}
\label{fig:northoffset}
\end{figure}

\section{Resolução Vertical}
\label{sec:delta_z}

A componente vertical de~\eqref{eq:k_bist} não depende do azimute:
\begin{equation}
  K_z=\frac{2\pi f}{c}\,q,\qquad
  \boxed{q=n(\cos\theta_T+\cos\theta_R)=\sqrt{n^2-\sin^2\psi_T}+\sqrt{n^2-\sin^2\psi_R}}.
  \label{eq:q}
\end{equation}
O suporte em $k_z$ é gerado por duas dispersões: a banda ($f\in f_c\pm B/2$) e a
variação de $q$ ao longo da abertura. Para distribuições aproximadamente
uniformes de cada uma, a PSF vertical é o produto de duas funções sinc.
Igualando a curvatura do lóbulo principal, a largura a $-3\dB$ é a de uma
única sinc com banda equivalente $\sqrt{(B\bar q)^2+12f_c^2\sigma_q^2}$:
\begin{equation}
  \boxed{\;\delta_z=\frac{c_w\,c}{\sqrt{(B\,\bar q)^2+12\,f_c^2\,\sigma_q^2}}\;},\qquad
  \bar q=\langle q\rangle,\quad \sigma_q^2=\var(q),
  \label{eq:delta_z}
\end{equation}
com $c_w=0,886$ para espectro retangular e $c_w=1,443$ para janela de Hann. O
primeiro termo do radicando é a resolução por banda, refratada pelo fator
$\bar q\approx2n\cos\theta_t$; o segundo é a contribuição tomográfica da
excursão de incidência.

\begin{observacao}[Relação com a soma linear das contribuições]
Uma forma usual na literatura de SAR helicoidal escreve a extensão vertical
como soma linear das contribuições de banda e de abertura, o que equivale a
$\delta_z=c/(B\bar q+f_c\Delta q)$, com $\Delta q$ a excursão total de $q$. A
expressão~\eqref{eq:delta_z} difere em dois pontos: as contribuições somam-se
em quadratura, porque correspondem a dispersões independentes do suporte, e a
constante $c_w$ converte a largura do suporte em largura a $-3\dB$. Na hélice de
referência, a soma linear subestima $\delta_z$ em $9\,\%$
(Seção~\ref{sec:res_valid}).
\end{observacao}

\section{Caso Helicoidal em Forma Explícita}
\label{sec:res_explicito}

Com o alvo no eixo e o receptor na mesma hélice desfasada de $\Delta\phi$,
$\psi_R=\psi_T$ e $\phi'_R-\phi'_T=\Delta\phi$ em todos os pulsos. Se a
incidência varre linearmente o intervalo $[\psi_1,\psi_2]$, as médias
de~\eqref{eq:delta_xy} e~\eqref{eq:delta_z} integram-se em forma fechada. Com
$Q(\psi)=\sqrt{n^2-\sin^2\psi}$:
\begin{align}
  \delta_x=\delta_y&=\frac{x_{J_0}\,\lambda_0}{2\pi\,|\cos(\Delta\phi/2)|}\;
  \frac{\psi_2-\psi_1}{\cos\psi_1-\cos\psi_2},\label{eq:dxy_eixo}\\
  \delta_z&=\frac{c_w\,c}{\sqrt{B^2\,[Q(\psi_1)+Q(\psi_2)]^2+4f_c^2\,[Q(\psi_1)-Q(\psi_2)]^2}} ,
  \label{eq:dz_eixo}
\end{align}
onde se usou que, para $q$ linear no tempo, $12\sigma_q^2=\Delta q^2$ com
$\Delta q=2[Q(\psi_1)-Q(\psi_2)]$. A expressão~\eqref{eq:dz_eixo} não contém
$\Delta\phi$: \textbf{a resolução vertical é invariante à separação azimutal}
quando as duas plataformas voam à mesma altura e no mesmo raio.

\paragraph{Aplicação numérica.} Para a hélice de referência vista de um alvo
em $(0;\,0;\,-2)\,$m, $\psi_1=46,11^\circ$ e $\psi_2=60,37^\circ$. Com
$f_c=425\,$MHz ($\lambda_0=0,7059\,$m), $B=50\,$MHz e $n=2$:
$\langle\sin\psi\rangle=0,7991$, donde $\delta_{xy}^{\text{mono}}=0,158\,$m e
$\delta_{xy}(90^\circ)=0,158/\cos45^\circ=0,224\,$m; e $Q(\psi_1)=1,8656$,
$Q(\psi_2)=1,8012$, donde $B\bar q=1,833\times10^{8}\,$Hz,
$f_c\Delta q=0,548\times10^{8}\,$Hz e $\delta_z=1,389\,$m para qualquer
$\Delta\phi$. A banda domina a resolução vertical: a contribuição tomográfica
acrescenta apenas $4\,\%$ à banda equivalente.

\section{Alvos Fora do Eixo}
\label{sec:res_offaxis}

As expressões~\eqref{eq:delta_xy} e~\eqref{eq:delta_z} valem para qualquer
$\vect P$; fora do eixo, mudam apenas as médias. A primeira ordem em $r_P/\rho$,
com $r_P=|\vect P_h|$ e $\rho$ o raio da órbita, para transmissor e receptor na
mesma hélice,
\begin{gather}
  \Delta\phi'(t)\simeq\Delta\phi-\frac{2r_P}{\rho}\,\sin\frac{\Delta\phi}{2}\,
  \cos\!\Big(\phi_T+\frac{\Delta\phi}{2}-\phi_P\Big),\label{eq:dphi_aparente}\\
  q(t)\simeq q_0(t)-2\,|Q'(\psi)|\,\frac{\cos^2\psi}{h}\,r_P\,
  \cos\frac{\Delta\phi}{2}\,\cos\!\Big(\phi_T+\frac{\Delta\phi}{2}-\phi_P\Big),
  \label{eq:q_aparente}
\end{gather}
com $Q'(\psi)=-\sin\psi\cos\psi/\sqrt{n^2-\sin^2\psi}$.
Daí decorrem duas consequências. A separação aparente oscila em torno de
$\Delta\phi$ com média nula, de modo que $\langle S\rangle$ e $\delta_{xy}$ não
mudam em primeira ordem: a lei $\delta_{xy}^{\text{mono}}/|\cos(\Delta\phi/2)|$
vale também fora do eixo. A diferença de incidência entre os lados próximo e
distante acrescenta a $q$ uma oscilação que eleva $\sigma_q$ e melhora
ligeiramente $\delta_z$; no canal biestático essa oscilação é escalada por
$|\cos(\Delta\phi/2)|$, porque transmissor e receptor estão em lados distintos e
suas variações se compensam. Este é o mecanismo do efeito
\textit{near-range} descrito por Góes~\citar{goes2022} para o caso monoestático,
agora estendido ao biestático.

\section{Limites do Ganho Biestático em Resolução}
\label{sec:res_limites}

\begin{proposicao}[Limite horizontal]
\label{prop:limite_xy}
Para qualquer geometria,
\begin{equation}
  \delta_{xy}^{\text{bi}}\ \ge\ \frac{x_{J_0}\,\lambda_0}{\pi\,\langle\sin\psi_T+\sin\psi_R\rangle},
\end{equation}
e $\delta_{xy}^{\text{bi}}<\delta_{xy}^{\text{mono}}$ exige
$\langle\sin\psi_R\rangle>\langle\sin\psi_T\rangle$.
\end{proposicao}

\begin{proof}
Pela desigualdade triangular em~\eqref{eq:S},
$S\le\tfrac12(\sin\psi_T+\sin\psi_R)$; substituindo em~\eqref{eq:delta_xy},
obtém-se a cota. Como $\delta_{xy}^{\text{mono}}$ corresponde a
$S=\sin\psi_T$, a desigualdade estrita exige que a média de
$\tfrac12(\sin\psi_T+\sin\psi_R)$ supere a de $\sin\psi_T$.
\end{proof}

\begin{proposicao}[Limite vertical]
\label{prop:limite_z}
O termo de banda de~\eqref{eq:delta_z} melhora em relação ao monoestático
somente se $\psi_R<\psi_T$, e a melhoria está limitada por
\begin{equation}
  \frac{\delta_z^{\text{mono}}}{\delta_z^{\text{bi}}}\bigg|_{\text{banda}}\le\frac{Q(\psi_T)+n}{2Q(\psi_T)} .
\end{equation}
\end{proposicao}

\begin{proof}
$Q$ é decrescente em $\psi$ e $Q\le n$, com igualdade no nadir. O termo de
banda é proporcional a $\bar q=Q(\psi_T)+Q(\psi_R)\le Q(\psi_T)+n$, contra
$2Q(\psi_T)$ no monoestático.
\end{proof}

Para $\psi_T=53,2^\circ$ e $n=2$, a cota vertical é $1,046$: a refração
concentra os raios no solo em torno da vertical ($\cos\theta_t\ge0,92$), e
nenhuma posição do receptor pode melhorar o termo de banda em mais de
$4,6\,\%$.

\begin{corolario}[Receptor à mesma altura e no mesmo raio]
\label{cor:mesma_altura}
Se $\psi_R=\psi_T$ em todos os pulsos, então
$\delta_{xy}^{\text{bi}}=\delta_{xy}^{\text{mono}}/|\cos(\Delta\phi/2)|\ge\delta_{xy}^{\text{mono}}$
e $\delta_z^{\text{bi}}=\delta_z^{\text{mono}}$.
\end{corolario}

As condições das Proposições~\ref{prop:limite_xy} e~\ref{prop:limite_z} são
opostas --- $\psi_R>\psi_T$ para o plano horizontal e $\psi_R<\psi_T$ para o
vertical --- e ambas requerem $\Delta\phi\approx0$, isto é, renunciar à
diversidade angular que justifica o receptor remoto. A consequência é
fundamental para o restante da tese: \emph{o canal biestático não deve ser
projetado para melhorar a resolução a $-3\dB$}. Se ele acrescenta informação
à imagem monoestática, esta deve aparecer em outras propriedades da PSF.

\section{Inclinação da Hélice}
\label{sec:helix_opt}

A contribuição tomográfica de~\eqref{eq:delta_z} é proporcional à excursão
$\Delta q\approx2|Q'(\psi_0)|\,\Delta\psi$. Para uma hélice de comprimento de
arco $L$ e inclinação $\beta=\arctan(\Delta z/\Delta\rho)$ vista a uma
distância $R_0$ sob a incidência $\psi_0$, a excursão de incidência é
$\Delta\psi\approx L|\cos(\beta-\psi_0)|/R_0$, máxima para $\beta=\psi_0$:
a linha de base da hélice deve ser perpendicular à linha de visada. É o caso da
hélice de referência, cuja distância ao centro da cena é praticamente
constante. Como a contribuição tomográfica soma-se em quadratura à de banda,
porém, o ganho dessa condição é pequeno quando $B\bar q\gg f_c\Delta q$.

\section{Validação}
\label{sec:res_valid}

As expressões fechadas foram comparadas com as larguras medidas em imagens
formadas por retroprojeção com refração exata (Seção~\ref{ssec:bp_algo}), em
cortes finos (passo de $2$ a $5\,$mm) pelo pico. A
Tabela~\ref{tab:res_verif} reúne dezesseis das 34 geometrias avaliadas: dois
alvos (no eixo e descentrado em $(6;\,4;\,-2)\,$m), cinco separações e três
geometrias do receptor fora da mesma hélice.

\begin{table}[H]
\centering\footnotesize\setlength{\tabcolsep}{3.4pt}
\caption{Verificação das expressões~\eqref{eq:delta_xy} e~\eqref{eq:delta_z}
frente às resoluções medidas na imagem. $\bar\psi_R$: incidência média do
receptor; $\varepsilon$: erro relativo da forma fechada.}
\label{tab:res_verif}
\begin{tabular}{@{}llccccccccc@{}}
\toprule
Alvo & Receptor & $\bar\psi_R$ & $\langle S\rangle$ & $\delta_x$ & $\delta_y$ & $\delta_{xy}^{\,\mathrm{fech}}$ & $\varepsilon_x/\varepsilon_y$ [\%] & $\delta_z$ & $\delta_z^{\,\mathrm{fech}}$ & $\varepsilon_z$ [\%] \\
 & & [°] & & [m] & [m] & [m] & & [m] & [m] & \\
\midrule
eixo & monoestático & 53,2 & 0,799 & 0,159 & 0,159 & 0,158 & $-0,4/-0,5$ & 1,387 & 1,389 & $+0,1$ \\
 & $\Delta\phi=45^{\circ}$ & 53,2 & 0,738 & 0,171 & 0,173 & 0,171 & $+0,2/-1,1$ & 1,387 & 1,389 & $+0,1$ \\
 & $\Delta\phi=90^{\circ}$ & 53,2 & 0,565 & 0,223 & 0,227 & 0,224 & $+0,4/-1,3$ & 1,387 & 1,389 & $+0,1$ \\
 & $\Delta\phi=135^{\circ}$ & 53,2 & 0,306 & 0,413 & 0,418 & 0,414 & $+0,2/-1,0$ & 1,387 & 1,389 & $+0,1$ \\
 & $\Delta\phi=165^{\circ}$ & 53,2 & 0,104 & 1,216 & 1,221 & 1,213 & $-0,3/-0,6$ & 1,387 & 1,389 & $+0,1$ \\
 & $\Delta\phi=0$, $\kappa=1,3$ & 60,1 & 0,832 & 0,153 & 0,153 & 0,152 & $-0,5/-0,5$ & 1,406 & 1,409 & $+0,2$ \\
 & $\Delta\phi=0$, $60\,$m abaixo & 70,6 & 0,870 & 0,146 & 0,146 & 0,146 & $-0,6/-0,6$ & 1,423 & 1,426 & $+0,2$ \\
 & $\Delta\phi=0$, $\kappa=0,2$ & 15,1 & 0,530 & 0,238 & 0,239 & 0,239 & $+0,4/-0,1$ & 1,364 & 1,371 & $+0,5$ \\
\midrule
descentrado & monoestático & 53,2 & 0,799 & 0,159 & 0,160 & 0,158 & $-0,2/-0,9$ & 1,378 & 1,380 & $+0,1$ \\
 & $\Delta\phi=45^{\circ}$ & 53,2 & 0,738 & 0,171 & 0,174 & 0,171 & $+0,4/-1,5$ & 1,382 & 1,384 & $+0,2$ \\
 & $\Delta\phi=90^{\circ}$ & 53,2 & 0,565 & 0,223 & 0,227 & 0,224 & $+0,4/-1,5$ & 1,386 & 1,388 & $+0,2$ \\
 & $\Delta\phi=135^{\circ}$ & 53,2 & 0,306 & 0,414 & 0,418 & 0,414 & $-0,0/-0,9$ & 1,388 & 1,390 & $+0,2$ \\
 & $\Delta\phi=165^{\circ}$ & 53,2 & 0,104 & 1,213 & 1,219 & 1,213 & $-0,0/-0,5$ & 1,387 & 1,389 & $+0,1$ \\
 & $\Delta\phi=0$, $\kappa=1,3$ & 60,1 & 0,832 & 0,153 & 0,154 & 0,152 & $-0,3/-0,9$ & 1,399 & 1,402 & $+0,2$ \\
 & $\Delta\phi=0$, $60\,$m abaixo & 70,6 & 0,869 & 0,146 & 0,147 & 0,146 & $-0,3/-1,1$ & 1,416 & 1,420 & $+0,3$ \\
 & $\Delta\phi=0$, $\kappa=0,2$ & 15,4 & 0,531 & 0,237 & 0,240 & 0,238 & $+0,5/-0,8$ & 1,358 & 1,364 & $+0,4$ \\
\bottomrule
\end{tabular}
\end{table}

Nas 34 geometrias com $\Delta\phi<180^\circ$, $|\varepsilon_x|\le0,6\,\%$,
$|\varepsilon_y|\le1,6\,\%$ e $|\varepsilon_z|\le0,5\,\%$. A largura $\delta_y$ é até
$2\,\%$ maior que $\delta_x$ porque a abertura de duas voltas com raio variável
não é perfeitamente isotrópica. As quatro últimas linhas de cada bloco
confirmam as Proposições~\ref{prop:limite_xy} e~\ref{prop:limite_z}: com o
receptor mais afastado ($\kappa=1,3$) ou $60\,$m mais baixo, $\delta_{xy}$
melhora $4$ e $8\,\%$, mas $\delta_z$ piora $1,4$ e $2,6\,\%$; com o receptor
próximo do nadir ($\kappa=0,2$), $\delta_z$ melhora $1,7\,\%$, mas $\delta_{xy}$
piora $50\,\%$. Na mesma hélice, $\delta_z$ é idêntica para todas as separações,
como prevê o Corolário~\ref{cor:mesma_altura}.

A forma de soma linear das contribuições verticais forneceria, na hélice de
referência, $\delta_z=c/(B\bar q+f_c\Delta q)=1,26\,$m, $9\,\%$ abaixo da largura
medida ($1,387\,$m). A retroprojeção foi também confrontada com uma
implementação independente em banda X ($10\,$GHz), com concordância de
$0,8\,\%$ na largura horizontal ($6,66$ contra $6,71\,$mm) e de $1\,\%$ na
vertical ($0,200$ contra $0,202\,$m).

Os limites de validez das expressões são três. (i) \emph{Abertura completa}:
\eqref{eq:delta_xy} supõe $360^\circ$ de azimute e dispersão de $|\vect K_h|$
pequena frente à média; falha quando $\Delta\phi\to180^\circ$, em que
$\langle S\rangle\to0$ e a largura é fixada pela curvatura residual da frente de
onda ($17\,$m para o alvo no eixo). (ii) \emph{Pesos uniformes}: o ganho de
antena, $1/R^2$ e a transmitância variam pouco ao longo da hélice; com janela
de Hann, usa-se $c_w=1,443$. (iii) \emph{Solo homogêneo}: as perdas atenuam,
mas não alteram a geometria de $\vect K$.

\section{Conclusões do Capítulo}
\label{sec:res_concl}

As resoluções $\delta_x$, $\delta_y$ e $\delta_z$ de um SAR com refração, para
qualquer geometria Tx--Rx, são funções de dois fatores de abertura vistos do
alvo: a média de $S$ fixa a resolução horizontal e as estatísticas de $q$ fixam
a vertical. As expressões reproduzem a imagem com erro inferior a $1,6\,\%$ e
mostram que a resolução não pode ser o objetivo do canal biestático: à mesma
altura, ele piora a resolução horizontal como $1/|\cos(\Delta\phi/2)|$ e
preserva a vertical; em qualquer geometria, o ganho vertical está limitado a
poucos por cento pela refração. O valor do receptor remoto deve ser buscado na
forma do suporte no espaço de números de onda, e não na sua extensão máxima.

%%═══════════════════════════════════════════════════════════════════════════════
\chapter{Transmissão em Meios Estratificados com $N$ Camadas}
\label{cap:multicamadas}
%%═══════════════════════════════════════════════════════════════════════════════

\section{Introdução e Estado da Arte}
\label{sec:multi_arte}

O formalismo de matrizes de transferência (ABCD) para propagação de ondas
planas através de pilhas de camadas dielétricas é uma ferramenta clássica,
sistematizada por Pozar~\citar{pozar2011} no contexto de redes de duas portas e
derivável diretamente das condições de contorno de Maxwell conforme Born e
Wolf~\citar{bornwolf1999}. Na literatura de GPR, Daniels~\citar{daniels2004}
descreve o tratamento padrão de meios estratificados, e Ishimaru~\citar{ishimaru1991}
fornece o arcabouço de propagação em meios aleatórios estratificados. Banda et
al.~\citar{banda2016} aplicam tomografia SAR polarimétrica em banda P a gelo
estratificado, mostrando a relevância de modelos multicamada em sensoriamento
tomográfico.

Este capítulo aplica o formalismo ao caso de incidência oblíqua variável,
determinada dinamicamente pela geometria helicoidal, com impedâncias de entrada
e de saída assimétricas e a última camada tratada como semi-infinita. Seu papel
na tese é duplo: fornecer o traçado de raios e a transmitância para cenários
estratificados, e mostrar que, num meio estratificado, o máximo de transmissão
deixa de coincidir com o ângulo de Brewster da primeira interface --- mais um
argumento contra o uso do ângulo de Brewster como critério geométrico de
projeto.

\section{Geometria do Modelo de $N$ Camadas}
\label{sec:multi_geom}

Considera-se uma pilha de camadas horizontais; a $i$-ésima tem índice de
refração $n_i$ e espessura $d_i$, e o meio de incidência (ar) tem $n_0=1$. O
alvo encontra-se na camada $L$, à profundidade total
$z_P = -\!\sum_{i=1}^{L-1}d_i - \Delta z_L$, onde $\Delta z_L$ é a profundidade
dentro da camada $L$. A camada $L$ é tratada como \textbf{semi-infinita}: toda
a energia transmitida propaga-se em direção ao alvo sem reflexão de retorno
por uma face inferior, o que representa a situação física de um alvo imerso
num estrato espesso.

\section{O Invariante de Snell Generalizado}
\label{sec:snell_inv}

Em meios estratificados com interfaces horizontais, a componente horizontal do
vetor de onda é conservada em cada interface:
\begin{equation}
  p = n_0\sin\theta_0 = n_1\sin\theta_1 = \cdots = n_L\sin\theta_L .
  \label{eq:snell_inv}
\end{equation}
O escalar $p$ é o \emph{parâmetro de raio}. O deslocamento horizontal acumulado
pelo raio até o alvo, a partir do ponto de entrada no solo, é
\begin{equation}
  x(p) = \sum_{i=1}^{L-1} \frac{p\,d_i}{\sqrt{n_i^2-p^2}}
         + \frac{p\,\Delta z_L}{\sqrt{n_L^2-p^2}} .
  \label{eq:horiz_p}
\end{equation}

\begin{proposicao}
\label{prop:unicidade_p}
A função $x(p)$ é estritamente crescente em $p \in [0, \min_i n_i)$, o que
garante a existência e unicidade de $p^*$ tal que $x(p^*) = x_{\text{alvo}}$:
\begin{equation}
  \frac{\dd x}{\dd p} = \sum_{i=1}^{L-1} \frac{n_i^2\,d_i}{(n_i^2-p^2)^{3/2}}
                        + \frac{n_L^2\,\Delta z_L}{(n_L^2-p^2)^{3/2}} > 0 .
\end{equation}
\end{proposicao}

O parâmetro $p^*$ é determinado por bisseção sobre $x(p)-x_{\text{alvo}}=0$,
incluindo o trecho no ar, com convergência em
$\lceil\log_2((p_{\max}-p_{\min})/\varepsilon)\rceil$ iterações.

\section{Formalismo de Matrizes de Transferência ABCD}
\label{sec:abcd}

A propagação TM em cada camada é descrita pela equação de Helmholtz
\begin{equation}
  \frac{\dd^2 H_y}{\dd z^2} + k_{z,i}^2\,H_y = 0, \qquad
  k_{z,i} = \frac{\omega}{c}\sqrt{n_i^2-p^2} .
  \label{eq:helmholtz_layer}
\end{equation}
As condições de contorno TM (continuidade de $H_y$ e de
$(1/n_i^2)\,\partial_z H_y$) nas faces da camada $i$ são encapsuladas pela
\textbf{matriz de transferência}
\begin{equation}
  \mathbf{M}_i = \begin{pmatrix}
    \cos(k_{z,i}\,d_i) & \jj\,Z_i\,\sin(k_{z,i}\,d_i) \\[4pt]
    \dfrac{\jj}{Z_i}\,\sin(k_{z,i}\,d_i) & \cos(k_{z,i}\,d_i)
  \end{pmatrix}, \qquad
  Z_i = \frac{n_i}{\cos\theta_i} = \frac{n_i^2}{\sqrt{n_i^2-p^2}},
  \label{eq:ABCD_matrix}
\end{equation}
derivada no Apêndice~\ref{app:abcd_deriv}. Para meios sem perdas,
$\det\mathbf{M}_i = 1$. A matriz total das $L-1$ camadas finitas é
\begin{equation}
  \begin{pmatrix} A & B \\ C & D \end{pmatrix}
  = \mathbf{M}_1 \cdot \mathbf{M}_2 \cdots \mathbf{M}_{L-1},
  \qquad AD - BC = 1 .
  \label{eq:ABCD_total}
\end{equation}

\section{Coeficientes Globais e Transmitância}
\label{sec:gamma_global}

As impedâncias normalizadas de entrada e de saída são
\begin{equation}
  Z_{\text{in}} = \frac{1}{\sqrt{1-p^2}}, \qquad
  Z_{\text{out}} = \frac{n_L^2}{\sqrt{n_L^2-p^2}},
  \label{eq:Zin_Zout}
\end{equation}
em geral distintas. O coeficiente de reflexão global, visto do ar, e o
coeficiente de transmissão são
\begin{equation}
  \boxed{\Gamma_{\text{TM}} = \frac{A Z_{\text{out}} + B
  - Z_{\text{in}}(C Z_{\text{out}} + D)}{%
  A Z_{\text{out}} + B + Z_{\text{in}}(C Z_{\text{out}} + D)}},
  \qquad
  \tau_{\text{TM}} = \frac{2\,Z_{\text{out}}}{A\,Z_{\text{out}} + B
  + C\,Z_{\text{in}}\,Z_{\text{out}} + D\,Z_{\text{in}}},
  \label{eq:Gamma_TM}
\end{equation}
e a transmitância de potência é
\begin{equation}
  \boxed{T_{\text{TM}} = |\tau_{\text{TM}}|^2 \cdot \frac{Z_{\text{in}}}{Z_{\text{out}}}
  = 1 - |\Gamma_{\text{TM}}|^2},
  \label{eq:T_TM_multi}
\end{equation}
onde o fator $Z_{\text{in}}/Z_{\text{out}}$ corrige a diferença entre os
semiespaços de entrada e de saída na contabilização do fluxo de
potência~\citar{bornwolf1999}. Os dois coeficientes
de~\eqref{eq:Gamma_TM} têm o mesmo denominador, e a igualdade
de~\eqref{eq:T_TM_multi} é verificada numericamente com erro inferior a
$10^{-15}$ em todos os cenários da Seção~\ref{sec:multi_sim}.

\begin{lema}
\label{lema:reducao_fresnel}
Para $L=1$ (alvo na primeira camada), $\Gamma_{\text{TM}}$ reduz-se ao
coeficiente de Fresnel TM~\eqref{eq:r_TM}.
\end{lema}

\begin{proof}
$L=1 \Rightarrow [A,B;C,D]=[1,0;0,1]$, e
\[
  \Gamma = \frac{Z_{\text{out}} - Z_{\text{in}}}{Z_{\text{out}} + Z_{\text{in}}}
  = \frac{n_2/\cos\theta_t - 1/\cos\theta_i}{n_2/\cos\theta_t + 1/\cos\theta_i}
  = \frac{n_2\cos\theta_i - \cos\theta_t}{n_2\cos\theta_i + \cos\theta_t}
  = r_{\text{TM}},
\]
usando $p = \sin\theta_i = n_2\sin\theta_t$, de modo que
$\sqrt{n_2^2-p^2} = n_2\cos\theta_t$ e $\sqrt{1-p^2} = \cos\theta_i$.
\end{proof}

\section{Ângulo de Transmissão Ótimo para a Pilha}
\label{sec:theta_opt}

Para uma interface única, o máximo de $T_{\text{TM}}$ ocorre em
$\theta_B=\arctan(n_2/n_1)$. Para uma pilha, as reflexões múltiplas internas
criam interferências que deslocam o máximo, que passa a ser determinado
numericamente:
\begin{equation}
  \theta_{\text{opt}} = \arg\max_{\theta_0 \in [0,\pi/2)} T_{\text{TM}}(f,\theta_0;\{n_i,d_i\}) .
  \label{eq:theta_opt}
\end{equation}

\begin{observacao}[Inacessibilidade do ângulo de Brewster interno]
Para $n_0=1$, $n_1=\sqrt{2}$, $n_2=2$, o ângulo de Brewster da interface
$n_1$--$n_2$, medido na primeira camada, é
$\theta_B^{(12)} = \arctan(\sqrt{2}) \approx 54,7^\circ$. O ângulo
correspondente no ar exigiria $\sin\theta_0 = \sqrt{2}\sin 54,7^\circ = 1,15>1$:
o ângulo de Brewster da interface interna é inacessível a partir do ar, e o
máximo de transmitância da pilha resulta de interferência entre as
reflexões internas, sem correspondência com nenhum ângulo de Brewster
clássico.
\end{observacao}

\section{Resultados de Simulação}
\label{sec:multi_sim}

A transmitância~\eqref{eq:T_TM_multi} foi calculada a $425\,$MHz para quatro
pilhas (Fig.~\ref{fig:ms_transmitancia}, Tabela~\ref{tab:ms_transmitancia}):
(a) interface única com $n=2$; (b) camada de $n_1=1,5$ com espessura de um quarto
do comprimento de onda no meio sobre um substrato de $n_2=2,5$; (c) a mesma
camada com meia onda; e (d) um perfil de três estratos de índices crescentes
($1,4$; $2,0$; $2,8$), com espessuras de $0,10$ e $0,30\,$m.

\begin{figure}[H]
\centering
\includegraphics[width=0.82\textwidth]{ms_fig01_transmitancia_abcd.pdf}
\caption{Transmitância TM de quatro pilhas a $425\,$MHz em função da
incidência no ar. Os pontos marcam o máximo de cada curva; a linha pontilhada
marca o ângulo de Brewster da interface única ($63,4^\circ$).}
\label{fig:ms_transmitancia}
\end{figure}

\begin{table}[H]
\centering
\caption{Ângulo de máxima transmitância das pilhas da
Fig.~\ref{fig:ms_transmitancia}.}
\label{tab:ms_transmitancia}
\begin{tabular}{@{}lccc@{}}
\toprule
Pilha & $\theta_{\text{opt}}$ & $T_{\text{TM}}(\theta_{\text{opt}})$ & $T_{\text{TM}}(0^\circ)$ \\
\midrule
(a) interface única, $n=2$                 & $63,4^\circ$ & $1,000$ & $0,889$ \\
(b) camada de quarto de onda               & $0^\circ$    & $0,997$ & $0,997$ \\
(c) camada de meia onda                    & $61,7^\circ$ & $0,971$ & $0,816$ \\
(d) três estratos de índice crescente      & $58,4^\circ$ & $1,000$ & $0,934$ \\
\bottomrule
\end{tabular}
\end{table}

Três observações decorrem dos resultados. Primeiro, a camada de quarto de onda,
cujo índice é próximo da média geométrica dos meios adjacentes, funciona como
transformador de impedância: a transmitância é máxima na incidência
\emph{normal} e praticamente plana até $50^\circ$, de modo que o ângulo de
Brewster perde qualquer significado. Segundo, a camada de meia onda é
transparente na incidência normal, mas não oblíqua: o máximo, de $0,971$,
desloca-se para $61,7^\circ$ e não alcança a unidade. Terceiro, um perfil de
índice crescente suaviza o contraste de impedância e desloca o máximo para
incidências menores ($58,4^\circ$). Em todos os casos a variação de
$T_{\text{TM}}$ entre a incidência de projeto da hélice ($46$ a $60^\circ$) e o
máximo é inferior a $0,5\dB$ por travessia.

\section{Conclusões do Capítulo}
\label{sec:multi_concl}

O formalismo ABCD com última camada semi-infinita é uma generalização
consistente do caso de Fresnel (Lema~\ref{lema:reducao_fresnel}), e a
Proposição~\ref{prop:unicidade_p} assegura a unicidade do parâmetro de raio no
traçado através de $N$ camadas. Os resultados mostram que, em meios
estratificados, o ângulo de máxima transmissão depende da espessura elétrica
das camadas e da frequência, e pode estar a dezenas de graus do ângulo de
Brewster da primeira interface. Como a transmitância varia apenas frações de
decibel na faixa de incidências de interesse, ela não é um critério
discriminante para a geometria de aquisição --- conclusão que o
Capítulo~\ref{cap:potencia} estende ao enlace completo.

%%═══════════════════════════════════════════════════════════════════════════════
\chapter{O Critério de Potência e o Ângulo de Brewster: uma Reavaliação}
\label{cap:potencia}
%%═══════════════════════════════════════════════════════════════════════════════

\section{Introdução}
\label{sec:pot_intro}

Este capítulo apresenta a hipótese com que a investigação começou, a
metodologia de otimização construída para explorá-la e as evidências que levaram
a abandoná-la. O registro é deliberadamente completo: um resultado negativo bem
fundamentado delimita o espaço de projeto tanto quanto um positivo, e a
literatura revisada não o documenta. A hipótese pode ser enunciada da seguinte
forma.

\begin{hipotese}[Posicionamento por Brewster]
\label{hip:brewster}
Para cada posição do transmissor, existe uma posição do receptor que maximiza a
potência recebida de um alvo enterrado, e essa posição é aquela em que o raio
refratado alvo$\to$receptor incide na interface no ângulo de Brewster. A
trajetória ótima do receptor resulta de um compromisso entre essa posição e a
resolução tridimensional da imagem.
\end{hipotese}

Uma revisão bibliográfica dirigida, que incluiu deliberadamente o caso
monoestático, não encontrou trabalho que combine tomografia SAR com otimização
da geometria de aquisição guiada pelo ângulo de Brewster. Os antecedentes mais
próximos são o GPR biestático de par fixo de Orfeo et al.~\citar{orfeo2019},
que avalia ângulo e altura das antenas para um sensor pontual sem síntese de
abertura, e as observações passivas do amortecimento de Brewster em imagens SAR
monoestáticas de superfície~\citar{ouchi2017}, \citar{ouchi2018}, \citar{thirion2018}.
Restrições de configuração entre plataformas para resolução foram tratadas por
Chen et al.~\citar{chen2022}, sem critério de potência nem reconstrução 3D. A
lacuna existia; a questão era se ela continha uma contribuição.

\section{Formulação da Otimização Multiobjetivo}
\label{sec:otim_form}

Para uma grade de $N_{tg}$ alvos, a potência total recebida
segue~\eqref{eq:radar_refratado} somada sobre os alvos, e, para uma posição
fixa do transmissor, a posição horizontal do receptor é escolhida por
\begin{equation}
  \min_{\vect{r}_{Rx} \in \mathcal{D}}\; J(\vect{r}_{Rx};\alpha) =
  -(1-\alpha)\log_{10}\!\left[P_r(\vect{r}_{Rx})\right]
  + \alpha\,\log_{10}\!\left[\delta_{xy}(\vect{r}_{Rx})\,\delta_z(\vect{r}_{Rx})\right],
  \label{eq:cost}
\end{equation}
com $\alpha \in [0,1]$ o peso de compromisso potência--resolução. A escala
logarítmica torna comensuráveis grandezas que diferem em muitas ordens de
grandeza ($P_r\sim10^{-15}\,$W e $\delta_{xy}\delta_z\sim10^{-4}\,$m$^2$), dá ao
termo de potência a interpretação de perda de percurso em decibéis e mantém o
gradiente bem condicionado. Como $J$ não é convexa --- a potência depende de
$\vect r_{Rx}$ através de $R_R^2$ e de $T_2(\theta_R)$ ---, a solução foi obtida
por otimização local com múltiplos pontos de partida geometricamente motivados
(posições oposta, perpendiculares e de Brewster, esta à distância horizontal
$d_B = z_{Rx}\tan\theta_B + |z_P|/\tan\theta_B$ do centroide), seguida de seleção
do melhor resultado e de interpolação monótona por partes cúbicas de Hermite
da sequência de posições ótimas, que garante uma trajetória contínua com
velocidade contínua.

\begin{figure}[H]
\centering
\includegraphics[width=0.78\textwidth]{pv_fig07_pareto.pdf}
\caption{Fronteira de Pareto potência--resolução da formulação~\eqref{eq:cost}
para interface única. Cada ponto corresponde à posição ótima do receptor para
um valor de $\alpha$.}
\label{fig:pareto}
\end{figure}

Num meio de duas camadas ($n_1=2$, $n_2=3$, $d_1=d_2=1\,$m, $400\,$MHz), um estudo de
sensibilidade cruzando $\alpha\in\{0;\,0,3;\,0,6;\,1\}$ com grades de alvos de
$30$, $60$ e $100\,$m de lado (125 alvos cada) mostrou um compromisso monotônico
e fortemente não linear (Tabela~\ref{tab:alpha_grid_resumo},
Fig.~\ref{fig:tradeoff_multi}): de $\alpha=0$ a $0,3$ o custo em potência é de
$1,4$ a $1,6\dB$ e o produto de resoluções cai mais de $98\,\%$, porque em
$\alpha=0$ o receptor ignora a resolução e assume geometrias quase diametrais;
acima de $\alpha=0,6$ o ganho de resolução é marginal.

\begin{table}[H]
\centering
\caption{Métricas médias ao longo da trajetória do transmissor no estudo de
sensibilidade $\alpha$--grade (meio de duas camadas, 125 alvos).}
\label{tab:alpha_grid_resumo}
\small
\begin{tabular}{@{}ccrrrr@{}}
\toprule
Grade (m) & $\alpha$ & Potência (dBm) & $\delta_{xy}$ (cm) &
$\delta_z$ (cm) & $\delta_{xy}\delta_z$ (cm$^2$) \\
\midrule
30  & 0     & $-24,6$ & 1748,7 & 89,8 & $1,58\times10^{5}$ \\
30  & 0,3   & $-26,2$ & 24,4   & 88,8 & 2169 \\
30  & 0,6   & $-28,2$ & 17,5   & 93,2 & 1628 \\
30  & 1     & $-31,0$ & 15,6   & 98,5 & 1540 \\
\midrule
60  & 0     & $-27,9$ & 1154,1 & 89,8 & $1,04\times10^{5}$ \\
60  & 0,3   & $-29,2$ & 21,8   & 89,3 & 1943 \\
60  & 0,6   & $-31,3$ & 17,2   & 93,9 & 1611 \\
60  & 1     & $-33,8$ & 15,6   & 98,5 & 1540 \\
\midrule
100 & 0     & $-31,0$ & 2196,8 & 89,7 & $1,98\times10^{5}$ \\
100 & 0,3   & $-32,4$ & 25,5   & 89,3 & 2275 \\
100 & 0,6   & $-34,5$ & 17,1   & 94,3 & 1608 \\
100 & 1     & $-36,8$ & 15,6   & 98,5 & 1540 \\
\bottomrule
\end{tabular}
\end{table}

\begin{figure}[H]
\centering
\includegraphics[width=0.78\textwidth]{pv_fig11_tradeoff_multicamada.png}
\caption{Curva de compromisso potência--resolução para o meio de duas camadas;
cada ponto de uma série corresponde a um valor de $\alpha$.}
\label{fig:tradeoff_multi}
\end{figure}

Esses resultados eram internamente coerentes, mas deixavam em aberto a questão
decisiva: \emph{a potência recebida distingue de fato geometrias boas de
geometrias ruins?} As três seções seguintes respondem que não.

\section{Primeira Evidência: o Alcance Domina a Transmitância}
\label{sec:pot_alcance}

A equação do radar~\eqref{eq:radar_refratado} é proporcional a
$T_1T_2/(R_T^2R_R^2)$. A transmitância de uma travessia varia entre cerca de
$0,33$ (incidência rasante) e $1$ (Brewster), isto é, no máximo $4,8\dB$; o termo
$R_R^2$ varia dezenas de decibéis com a posição do receptor. Na solução de
potência pura ($\alpha=0$), o receptor ótimo ficou, em média, a $20^\circ$ do
ângulo de Brewster no solo. A Tabela~\ref{tab:brewster_vs_alcance} decompõe a
diferença para uma posição representativa do transmissor.

\begin{table}[H]
\centering
\caption{Receptor ótimo em potência frente à posição de Brewster, para a mesma
posição do transmissor.}
\label{tab:brewster_vs_alcance}
\begin{tabular}{@{}lrr@{}}
\toprule
& Ótimo em potência & Posição de Brewster \\
\midrule
Distância horizontal ao alvo & $26,4\,$m & $229,5\,$m \\
Ângulo de refração no solo & $6,35^\circ$ & $26,57^\circ$ \\
Transmitância $T_2$ & $0,894$ & $1,000$ \\
Comprimento do trecho alvo$\to$Rx & $122,2\,$m & $260,0\,$m \\
Potência recebida & $\mathbf{-35,57\,dBm}$ & $-40,27\,$dBm \\
\midrule
Contribuição da transmitância & \multicolumn{2}{r}{$-0,48\dB$ (favorece Brewster)} \\
Contribuição do espalhamento geométrico & \multicolumn{2}{r}{$+6,56\dB$ (favorece a proximidade)} \\
\bottomrule
\end{tabular}
\end{table}

A posição próxima vence por $4,7\dB$: mover o receptor até Brewster multiplica o
trecho de retorno por $2,1$ e custa $6,6\dB$ de espalhamento, enquanto a
transmitância só pode devolver $0,5\dB$. Como contraprova, uma otimização que
maximiza \emph{apenas} as transmitâncias, retirando o ganho de antena e
$1/R^2$ do custo, localiza o ângulo de Brewster no solo com desvio médio de
$0,04^\circ$ em 275 posições --- o mecanismo de busca está correto --- mas
recebe $3,6\dB$ \emph{menos} potência que a linha de base.

O mesmo desequilíbrio aparece no projeto do transmissor. O termo tomográfico da
resolução vertical favorece incidências próximas de $45^\circ$, e o ângulo de
Brewster ar--solo é sempre maior que $45^\circ$ quando $n_2>n_1$: os dois
ótimos nunca coincidem. Uma estratégia em duas etapas --- primeiro a incidência
média do transmissor, depois o receptor --- baseava-se na fatoração exata
$P_r=K\,f_{Tx}\,f_{Rx}$, da qual decorreria um ganho aditivo independente da
segunda etapa. Com a geometria real e o receptor reotimizado, porém, a
aditividade falha: com peso de resolução $0,3$, a estratégia em duas etapas
piorou as quatro métricas (potência $-1,41\dB$, $\delta_{xy}$ $+6,7\,\%$,
$\delta_z$ $+2,4\,\%$ e custo); com peso nulo, ganhou $0,81\dB$, contra
$0,02\dB$ previstos pela transmitância, ganho que provém majoritariamente do
alcance. A fatoração vale para uma posição fixa do receptor, mas o receptor não
permanece fixo quando o transmissor muda.

\section{Segunda Evidência: Invariância Azimutal da Potência}
\label{sec:pot_azimute}

Quando o alvo está sobre o eixo vertical das hélices, todos os fatores
de~\eqref{eq:radar_refratado} dependem apenas do raio e da altura instantâneos
das plataformas: os comprimentos refratados, as incidências e as
transmitâncias, o ganho do transmissor (cuja trajetória não muda) e o ganho do
receptor apontado ao alvo. A separação azimutal apenas gira a hélice do
receptor em torno desse eixo e, portanto, \emph{não altera a potência}.

\begin{proposicao}[Invariância azimutal]
\label{prop:inv_azimutal}
Para um alvo pontual isotrópico sobre o eixo de rotação comum das hélices,
$P_r(t)$ é independente da separação azimutal $\Delta\phi$, enquanto
$\delta_{xy}\propto1/|\cos(\Delta\phi/2)|$.
\end{proposicao}

A simulação confirma a proposição bit a bit: com o receptor em $\Delta\phi=0$
e em $\Delta\phi=180^\circ$, as séries temporais de potência coincidem
($-75,09\,$dBm de média em ambas), e a resolução horizontal, em banda X,
passa de $0,015$ para $0,19\,$m. A separação azimutal é um grau de liberdade
energeticamente gratuito que só afeta a imagem.

\section{Terceira Evidência: a Orientação do Alvo Anula a Vantagem}
\label{sec:pot_orientacao}

A terceira evidência vem do alvo. Calculou-se a potência recebida de cinco
sólidos condutores (cubo, caixa, cilindro, esfera e cone) enterrados a $2\,$m
num solo de $n=2$, por Óptica Física biestática avaliada no meio, a $425\,$MHz,
ao longo de um arco de raio constante, comparando a configuração monoestática
com a configuração biestática especular (receptor simétrico ao transmissor em
relação à vertical do alvo).

\begin{figure}[H]
\centering
\includegraphics[width=\textwidth]{fig_resumen_objetos.png}
\caption{Diferença de energia entre as configurações biestática especular e
monoestática para cinco alvos. (a) Alvo em orientação nominal. (b) Monte Carlo
com orientação aleatória (média sobre 14 posições fora do nadir e intervalo de
confiança de $95\,\%$). Note-se a diferença de escala vertical entre os
painéis.}
\label{fig:pot_objetos}
\end{figure}

Em orientação nominal, os alvos com uma face horizontal (cubo, caixa e cilindro)
recebem de $6,2$ a $9,6\dB$ mais energia na configuração especular, porque a
bissetriz transmissor--receptor permanece vertical e mantém a reflexão especular
dessa face em toda a trajetória, enquanto o monoestático só a observa no nadir.
A vantagem é, contudo, uma propriedade do \emph{alinhamento} entre a geometria e
uma face do alvo. Com $10^4$ orientações aleatórias por posição, a diferença
agregada cai para valores entre $-0,37$ e $+0,21\dB$
(Tabela~\ref{tab:mc_potencia}), e cada configuração supera a outra em
aproximadamente metade das repetições.

\begin{table}[H]
\centering
\caption{Diferença de energia (especular menos monoestática) com orientação
aleatória do alvo: $10^4$ orientações por posição, 14 posições.}
\label{tab:mc_potencia}
\begin{tabular}{@{}lrc@{}}
\toprule
Alvo & $\Delta$ [dB] (IC $95\,\%$) & $\Pr\{\text{especular}>\text{mono}\}$ \\
\midrule
Cubo     & $-0,26$ \ $[-0,37;\,-0,16]$ & 0,52 \\
Caixa    & $-0,18$ \ $[-0,26;\,-0,10]$ & 0,50 \\
Cilindro & $-0,26$ \ $[-0,30;\,-0,22]$ & 0,50 \\
Esfera   & $-0,37$ \ $[-0,37;\,-0,36]$ & 0,38 \\
Cone     & $+0,21$ \ $[+0,15;\,+0,27]$ & 0,57 \\
\bottomrule
\end{tabular}
\end{table}

A explicação é geométrica. Pelo teorema de equivalência
monoestático--biestático~\citar{kell1965}, a RCS biestática aproxima a
monoestática observada ao longo da bissetriz, com uma redução que depende do
ângulo biestático $\beta$. Para orientações uniformes, a probabilidade de que
uma normal de face aponte para o radar (monoestático) ou ao longo da bissetriz
(biestático) é a mesma; e o pico especular biestático, menor por
$\cos^2(\beta/2)$, tem um lóbulo proporcionalmente mais largo no espaço de
orientações. As duas dependências compensam-se na média, de acordo com o
limite de óptica geométrica, no qual a RCS de uma esfera independe do ângulo
biestático.

Nesse mesmo estudo, a transmitância de Brewster acrescentou apenas $+1,02\dB$
ao enlace de ida e volta em relação ao nadir, idêntica nas duas
configurações. O único efeito expressivo do ângulo de Brewster foi outro: na
configuração especular com interface lisa e sem perdas, ele anula a reflexão
especular da superfície, que no restante do arco supera o eco do alvo em $8$ a
$17\dB$, e eleva a relação alvo--superfície em mais de $40\dB$ num intervalo
estreito em torno de $\theta_B$.

\section{Discussão: a Refutação da Hipótese}
\label{sec:pot_discussao}

As três evidências convergem. A potência recebida não é um critério
discriminante para a posição do receptor: depende sobretudo do alcance, é
invariante à separação azimutal para alvos centrados e, promediada sobre a
orientação desconhecida do alvo, é a mesma em qualquer geometria. O ângulo de
Brewster acrescenta cerca de $1\dB$ ao enlace e não determina a posição do
receptor. A Hipótese~\ref{hip:brewster} está refutada.

Duas consequências orientam o restante da tese. A primeira é que a geometria
do receptor deve ser escolhida por seu efeito sobre a \emph{imagem} ---
suporte no espaço de números de onda, lóbulos laterais, contraste,
condicionamento --- e não sobre a energia; nenhuma das funções de custo dos
capítulos seguintes contém um termo de potência do alvo. A segunda é que o
ângulo de Brewster muda de papel: ele só é relevante como mecanismo de
\emph{supressão do retorno da superfície}, e sua utilidade deve ser medida pela
relação alvo--clutter, num solo com perdas, com rugosidade, erros angulares e
banda finita. O efeito observado em interface lisa e sem perdas --- um zero
exato num intervalo angular estreito --- é precisamente o tipo de mínimo que
pode não sobreviver a essas condições, o que é examinado no
Capítulo~\ref{cap:superficie}.

%%═══════════════════════════════════════════════════════════════════════════════
\chapter{Modelo de Simulação Coerente Monoestático--Biestático}
\label{cap:modelo}
%%═══════════════════════════════════════════════════════════════════════════════

\section{Requisitos e Convenções}
\label{sec:mod_req}

Para responder à pergunta reformulada no Capítulo~\ref{cap:intro}, o modelo de
simulação deve satisfazer seis requisitos: (i) gerar sinais complexos,
preservando amplitude, fase e retardo, em todos os canais; (ii) gerar os canais
monoestático e biestático para o \emph{mesmo} pulso, com a mesma trajetória do
transmissor, o mesmo alvo e a mesma física; (iii) incluir a refração exata, o
solo com perdas, a superfície rugosa e alvos com RCS anisotrópica e
dependente da polarização; (iv) normalizar explicitamente a fusão, para que o
aumento do número de amostras não seja confundido com ganho geométrico;
(v) em estudos de Monte Carlo, sortear uma orientação do alvo por
\emph{aquisição} e mantê-la fixa ao longo da abertura e nos dois canais; e
(vi) nunca usar a potência do eco do alvo como critério de seleção da geometria.

Foi desenvolvido para isso um simulador próprio, em que cada componente foi
validado contra soluções analíticas e contra uma implementação independente
(Seção~\ref{sec:mod_valid}).

\section{Geometria e Sinal}
\label{sec:mod_geom}

O transceptor percorre a hélice de referência da Seção~\ref{sec:fund_sar_bist},
discretizada em $1600$ pulsos (800 nos estudos de Monte Carlo); o receptor
segue~\eqref{eq:rx_traj}. Os parâmetros do sistema estão no
Apêndice~\ref{app:parametros}: portadora de $425\,$MHz, banda de $50\,$MHz
amostrada em 64 frequências, potência de pico de $100\,$kW, antenas de $6\,$dBi
com feixe gaussiano de $14^\circ\times50^\circ$ apontado ao centro da cena,
pulso de $1,3\,\mu$s e temperatura de ruído de $300\,$K. O eco de cada canal
segue~\eqref{eq:eco}, com a geometria do raio obtida pelo princípio de Fermat
com a parte real do índice, a atenuação pela parte imaginária e as
transmitâncias de amplitude $\tau=\sqrt{1-|r|^2}$ na polarização de cada trecho.
O alvo de referência é um ponto isotrópico de $1\,$m$^2$ em $(0;\,0;\,-2)\,$m.

\section{Modelo da Superfície}
\label{sec:mod_superficie}

A superfície contribui com duas componentes sujeitas à mesma equação do radar.

\paragraph{Componente especular coerente.} Para um plano infinito, a fase
estacionária localiza um único ponto especular por pulso, e a teoria de
imagens fornece a reflexão exata com o coeficiente de Fresnel do ângulo de
incidência nesse ponto, multiplicado pelo fator de rugosidade de
Ament~\citar{ament1953},
\begin{equation}
  \chi_s=\exp\!\Big[-\tfrac12k_0^2s^2(\cos\theta_i+\cos\theta_s)^2\Big],
\end{equation}
onde $s$ é a rugosidade rms. O ponto especular do par Tx--Rx encontra-se a meio
caminho entre os nadires das plataformas e sua incidência é
\begin{equation}
  \theta_{esp}=\arctan\!\big(\tan\theta_c\,\sin(\Delta\phi/2)\big),
  \label{eq:theta_esp}
\end{equation}
com $\theta_c$ a incidência das plataformas vista do centro da superfície. No
canal monoestático, a especular é o reflexo do nadir ($\theta_{esp}=0$); a
$90^\circ$ e $135^\circ$, com $\theta_c=63,45^\circ$, ela ocorre em $54,8^\circ$ e
$61,6^\circ$; \emph{apenas na configuração diametral} ela coincide com
$\theta_c$. Escolher $\theta_c=\theta_{pB}$ só coloca a reflexão especular no
pseudo-Brewster nessa configuração.

\paragraph{Componente difusa.} A superfície é discretizada em facetas de
$0,15\times0,15\,$m numa região de $\pm12\,$m em torno do alvo, cada uma com
amplitude complexa $a_p\sqrt{\sigma^0_{pp}\,\Delta A}$, em que
$a_p\sim\mathcal{CN}(0,1)$ é \emph{congelada} durante toda a abertura e
compartilhada pelos canais monoestático e biestático, como numa superfície
física. O coeficiente $\sigma^0_{pp}$ é o do método das pequenas perturbações
biestático de primeira ordem~\citar{ulaby1982}, \citar{fung1994}:
\begin{equation}
  \sigma^0_{pp} = 8k_0^4s^2\cos^2\theta_i\cos^2\theta_s\,|\alpha_{pp}|^2\,
  W\!\big(k_0[\sin\theta_s\cos\phi_s-\sin\theta_i],\;k_0\sin\theta_s\sin\phi_s\big),
  \label{eq:spm}
\end{equation}
\begin{equation}
  \alpha_{hh}=\frac{(\varepsilon-1)\cos\phi_s}{(\cos\theta_i+q_i)(\cos\theta_s+q_s)},\qquad
  \alpha_{vv}=\frac{(\varepsilon-1)\,[\varepsilon\sin\theta_i\sin\theta_s-\cos\phi_s\,q_iq_s]}
                  {(\varepsilon\cos\theta_i+q_i)(\varepsilon\cos\theta_s+q_s)},
  \label{eq:alpha_pp}
\end{equation}
com $q=\sqrt{\varepsilon-\sin^2\theta}$, $\phi_s$ o azimute de espalhamento medido a
partir da direção especular e $W$ o espectro de potência da rugosidade
(gaussiano, com comprimento de correlação $\ell$). A rugosidade de referência é
moderada: $s=2\,$cm e $\ell=20\,$cm ($k_0s=0,18$). Em retroespalhamento,
$\alpha_{hh}=r_{\text{TE}}$; na direção especular, $\alpha_{vv}=0$ exatamente em
Brewster.

\paragraph{Erro de atitude.} Um rolamento $\gamma$ da antena em torno da linha
de visada mistura as polarizações; no modelo co-polar,
$r_{\text{ef}}=\cos^2\gamma\,r_{\text{TM}}+\sin^2\gamma\,r_{\text{TE}}$.

\section{Modelo dos Alvos}
\label{sec:mod_alvos}

Os alvos são condutores perfeitos modelados por Óptica Física
biestática~\citar{knott1993}, com integral exata por triângulo da superfície
iluminada, avaliada \emph{no solo} (comprimento de onda $\lambda_0/n'$) nas
direções refratadas e em cada frequência. A biblioteca inclui uma esfera de raio
$0,5\,$m, um cubo de $1\,$m, uma caixa de $1\times0,5\times0,3\,$m, um cilindro
vertical de raio $0,3\,$m e altura $1\,$m, um tubo horizontal de raio $0,15\,$m e
comprimento $2\,$m, um tubo inclinado e um cilindro inclinado, além de uma cena
com três objetos de RCS e orientações distintas. A Óptica Física reproduz as RCS
de referência no nadir com diferença nula e a placa especular,
$4\pi A^2\cos^2\theta/\lambda^2$, com erro máximo de $0,1\dB$.

\paragraph{Polarização.} Os vetores de polarização $\hat v$ (TM) e $\hat h$ (TE)
são definidos de forma contínua em função da direção, de modo que os canais
monoestático e biestático convergem quando $\Delta\phi\to0$. Consideram-se dois
receptores. O \emph{receptor co-polar} recebe na mesma polarização nominal do
transmissor (TM--TM ou TE--TE), definida no próprio plano de incidência. O
\emph{receptor alinhado}, de dupla polarização, combina seus dois canais para
receber a componente do campo do transmissor transversal à direção de
recepção; no monoestático coincide com TM--TM, e a transmitância de subida
pondera TM e TE segundo essa projeção.

\section{Formação de Imagem e Fusão Normalizada}
\label{sec:mod_fusao}

Cada canal é reconstruído por~\eqref{eq:bp} com retardo refratado exato por
voxel e pulso. Como cada imagem é a média coerente de suas amostras, definem-se
duas fusões:
\begin{itemize}
  \item \textbf{fusão completa}, $I_{\text{fus}}=\tfrac12(I_{\text{mono}}+I_{\text{bi}})$,
  com os dois canais completos ($2N$ amostras), que corresponde ao que o sistema
  efetivamente adquire;
  \item \textbf{fusão a orçamento constante}, com o canal monoestático nos pulsos
  ímpares e o biestático nos pares, totalizando as mesmas $N$ amostras da imagem
  monoestática.
\end{itemize}
As métricas de forma são invariantes à escala. A SNR é obtida propagando o ruído
térmico, com variância $kT/(T_pM_{\text{ef}})$ e $M_{\text{ef}}=M$ para os canais
isolados e para a fusão a orçamento constante, e $2M$ para a fusão completa.
Desse modo, o ganho de $3\dB$ por duplicar o número de amostras aparece
\emph{apenas} na SNR da fusão completa e é declarado como tal. Em todos os
estudos, as duas fusões produziram imagens com correlação superior a
$0,99999$: \textbf{toda variação de forma observada é geométrica}. Todas as
imagens são exibidas com a mesma referência (pico monoestático em $0\dB$) e a
mesma faixa dinâmica ($30\dB$).

\section{Métricas}
\label{sec:mod_metricas}

\begin{table}[H]
\centering\small
\caption{Métricas de avaliação da imagem, da informação, do clutter e da
operação.}
\label{tab:metricas}
\begin{tabular}{@{}L{3.3cm}L{11.0cm}@{}}
\toprule
$\delta_{x,y,z}$ & Largura a $-3\dB$ da PSF em amplitude em cortes pelo pico,
                   com interpolação em decibéis; censurada quando não há
                   cruzamento dentro da janela \\
$V_{-3}$ & Volume da componente conexa $|I|^2\ge\tfrac12$pico que contém o pico \\
Anisotropia & $\sqrt{\lambda_{\max}/\lambda_{\min}}$ do tensor de segundo momento do lóbulo a $-3\dB$ \\
PSLR, ISLR & Lóbulo principal definido pelo elipsoide com semieixos no primeiro nulo
             em $x$, $y$ e $z$, numa janela fixa de $2\times2\times5\,$m; ISLR é a razão
             entre a energia fora e dentro do lóbulo \\
Contraste & Pico dividido pela potência média fora do lóbulo principal \\
Localização & Distância entre o pico, refinado por ajuste parabólico, e a posição verdadeira \\
Separabilidade & Fração de alvos com máximo local e queda de pelo menos $3\dB$ para cada vizinho \\
Suporte em $k$ & Área ocupada no plano $(k_x,k_y)$ (células de $0,25\,$rad/m) \\
Condicionamento & Posto efetivo~\citar{roy2007} e coerência mútua da matriz de aquisição
                  linearizada $A_{iv}=e^{-\jj\vect K_i\cdot\vect r_v}$ sobre
                  $7\times7\times5$ voxels, com Gram normalizado por canal \\
SCR & Potência do alvo dividida pela potência esperada da superfície no voxel do
      alvo (na imagem) ou no enlace; PSCR: pico do alvo dividido pelo pico da
      superfície em toda a grade \\
SNR & $|I_{\text{pico}}|^2T_pM_{\text{ef}}/kT$ \\
Voo & Separação mínima entre VANTs, razão de velocidades, comprimento relativo e aceleração máxima \\
\bottomrule
\end{tabular}
\end{table}

\section{Validação do Modelo}
\label{sec:mod_valid}

\begin{table}[H]
\centering\small
\caption{Validações do simulador.}
\label{tab:mod_valid}
\begin{tabular}{@{}L{6.2cm}L{3.9cm}L{3.9cm}@{}}
\toprule
\textbf{Prova} & \textbf{Referência} & \textbf{Simulador} \\
\midrule
RCS por Óptica Física no nadir: cubo, caixa, cilindro, esfera e cone
 & 20,04; 14,02; 8,94; $-1,54$; $-15,86\,$dBsm & idênticas \\
Placa especular, $\theta=15$, $30$ e $45^\circ$
 & 19,74; 18,79; 17,03\,dBsm & 19,73; 18,88; 17,01\,dBsm \\
Retroprojeção frente a implementação independente, $10\,$GHz, corte em $x$ & $6,66\,$mm & $6,71\,$mm \\
Idem, eixo $z$ (monoestático / biestático com $\Delta\phi=0$) & $0,200$ / $0,206\,$m & $0,202$ / $0,207\,$m \\
PSF horizontal monoestática a $425\,$MHz frente ao anel de Bessel & $0,157\,$m & $0,1575\,$m \\
PSF biestática a $90^\circ$ frente a $\delta_{xy}/\cos(\Delta\phi/2)$ & $0,222\,$m & $0,221\,$m \\
Soma de facetas de Kirchhoff (plano liso) frente à teoria de imagens & $-0,102\,$dBm & $-0,095\,$dBm \\
Potência difusa média (40 realizações) frente a $\int\sigma^0\dd A$ & $0\dB$ & $-0,24$ a $+0,45\dB$ \\
$\alpha_{hh}$ em retroespalhamento frente a $r_{\text{TE}}$; $\alpha_{vv}$ especular em Brewster
 & $-0,5104$; $0$ & $-0,5104$; $0$ \\
\bottomrule
\end{tabular}
\end{table}

\section{Conclusões do Capítulo}
\label{sec:mod_concl}

O simulador gera, para o mesmo pulso e com a mesma física, os canais
monoestático e biestático complexos de alvos anisotrópicos enterrados num solo
com perdas sob uma superfície rugosa, e reconstrói as imagens por
retroprojeção com refração exata. A fusão a orçamento constante isola o ganho
geométrico, e as métricas cobrem a forma da PSF, os lóbulos, a informação, o
clutter e a operação. Cada componente foi validado contra soluções analíticas e
contra uma implementação independente.

%%═══════════════════════════════════════════════════════════════════════════════
\chapter{Complementaridade Geométrica no Espaço de Números de Onda}
\label{cap:geometria}
%%═══════════════════════════════════════════════════════════════════════════════

\section{Introdução}
\label{sec:geo_intro}

Este capítulo responde às três primeiras perguntas de pesquisa com o alvo
pontual isotrópico, que isola a geometria de qualquer efeito de polarização ou
de anisotropia: como a separação Tx--Rx altera o suporte no espaço de números
de onda e a PSF; qual separação produz a melhor imagem biestática isolada; e
qual separação faz o canal biestático acrescentar mais informação à imagem
monoestática. Em todos os casos o receptor voa na mesma hélice do transmissor
($\kappa=1$), desfasado de $\Delta\phi$, e o solo é o S1.

\section{As Configurações de Referência}
\label{sec:geo_ref}

A configuração monoestática é a referência permanente, porque o transceptor a
adquire sempre. A configuração diametral ($\Delta\phi=180^\circ$, mesma altura e
mesmo raio) é a referência biestática extrema: aproximadamente especular e,
em princípio, a de maior diversidade angular.

\begin{figure}[H]
\centering
\includegraphics[width=0.85\textwidth]{A_kspace_ref.pdf}
\caption{Onde se perde o suporte na configuração diametral. O canal
monoestático cobre um anel horizontal de raio $2k_0\sin\psi\approx14\,$rad/m;
na diametral, as direções de ida e de volta no solo são simétricas em relação à
vertical, $\vect K_h=0$ em todos os pulsos, e todo o suporte cai sobre o eixo
$k_z$.}
\label{fig:geo_kspace_ref}
\end{figure}

\begin{table}[H]
\centering\small
\caption{Métricas da PSF das configurações de referência. $\delta_x^{\text{longo}}$:
largura medida num corte de $\pm25\,$m. Fusões a orçamento constante, salvo
indicação.}
\label{tab:geo_ref}
\begin{tabular}{@{}lrrrrrrrr@{}}
\toprule
Canal & $\delta_x$ & $\delta_x^{\text{longo}}$ & $\delta_z$ & $V_{-3}$ & PSLR & ISLR & Contr. & SNR \\
      & [m] & [m] & [m] & [l] & [dB] & [dB] & [dB] & [dB] \\
\midrule
Monoestático              & 0,158 & 0,159 & 1,39 & 19,2 & $-7,9$ & 6,4 & 21,3 & 78,0 \\
Diametral                 & $>2$ & \textbf{17,2} & 1,39 & $>5331$ & $0,0$ & --- & 8,5 & 78,0 \\
Mono + diametral          & 0,235 & --- & 1,39 & 41,1 & $-3,8$ & 11,9 & 11,4 & 78,0 \\
Mono + diametral ($2N$)   & 0,235 & --- & 1,39 & 41,1 & $-3,8$ & 11,9 & 11,4 & 81,0 \\
\midrule
Biestático $90^\circ$     & 0,221 & 0,223 & 1,39 & 38,8 & $-7,9$ & 5,1 & 19,6 & 78,0 \\
Mono + bi $90^\circ$      & 0,182 & --- & 1,39 & 25,9 & $-9,9$ & 1,0 & 25,3 & 78,0 \\
\bottomrule
\end{tabular}
\end{table}

A configuração diametral reconstrói mal, como prevê o
Capítulo~\ref{cap:resolucao}: conserva a resolução vertical ($1,39\,$m), mas perde
toda a horizontal; o único foco lateral residual vem da curvatura da frente de
onda ($\delta_x=17\,$m a $190\,$m de distância). Sua fusão com o canal
monoestático é \emph{pior} que o monoestático isolado: o canal diametral
acrescenta um pedestal plano à PSF, eleva os lóbulos (PSLR de $-7,9$ para
$-3,8\dB$; ISLR de $6,4$ para $11,9\dB$) e alarga o lóbulo principal em $49\,\%$.
A potência média recebida é a mesma nos dois canais ($-68,87\,$dBm), confirmando
na formação de imagem a conclusão do Capítulo~\ref{cap:potencia}. As fusões
completa e a orçamento constante são idênticas em forma, e diferem apenas nos
$3\dB$ de SNR atribuíveis ao número de amostras.

\section{Varredura da Separação Azimutal}
\label{sec:geo_varredura}

\begin{figure}[H]
\centering
\includegraphics[width=\textwidth]{B_geometria.pdf}
\caption{Geometria da varredura. Acima, planta do transmissor e do receptor
para cada $\Delta\phi$: o receptor percorre a mesma hélice do transmissor,
girada de $\Delta\phi$ e sincronizada (marcas nos mesmos pulsos; linha
tracejada: linha de base). Abaixo, raio e altura das duas plataformas,
idênticos em todos os casos. As espiras aparecem separadas em planta porque,
num mesmo azimute, as plataformas passam em instantes e raios distintos.}
\label{fig:geo_B}
\end{figure}

A separação foi varrida de $0^\circ$ a $180^\circ$ em passos de $15^\circ$, sem
otimizador, para que a física se revelasse antes de qualquer função de custo.
Os anéis da Fig.~\ref{fig:geo_kspace} são vetores de número de onda, e não
trajetórias: como as plataformas voam à mesma altura e no mesmo raio, a
componente horizontal do canal biestático vale
$2k_0\sin\psi\,|\cos(\Delta\phi/2)|$, e o anel se contrai com a separação,
embora a órbita do receptor seja sempre a mesma.

\begin{figure}[H]
\centering
\includegraphics[width=\textwidth]{B_kspace.pdf}
\caption{Suporte no espaço de números de onda do canal monoestático, de três
separações intermediárias e da diametral. Acima, plano $(k_x,k_y)$; abaixo,
$(|k_h|,k_z)$. O canal biestático é um anel de raio $|\cos(\Delta\phi/2)|$ vezes
o monoestático, com a mesma extensão em $k_z$.}
\label{fig:geo_kspace}
\end{figure}

\begin{figure}[H]
\centering
\includegraphics[width=\textwidth]{B_resoluciones.pdf}
\caption{Resoluções em $x$, $y$ e $z$ em função de $\Delta\phi$. Curva
tracejada: $\delta_{xy}^{\text{mono}}/|\cos(\Delta\phi/2)|$, forma
fechada~\eqref{eq:dxy_eixo}.}
\label{fig:geo_res}
\end{figure}

\begin{figure}[H]
\centering
\includegraphics[width=\textwidth]{B_pslr_islr.pdf}
\caption{PSLR, ISLR e contraste em função de $\Delta\phi$: imagem monoestática
(azul), fusão completa (verde) e fusão a orçamento constante (pontilhado,
coincidente).}
\label{fig:geo_lobulos}
\end{figure}

\begin{figure}[H]
\centering
\includegraphics[width=\textwidth]{B_informacion.pdf}
\caption{Informação: área ocupada no plano $(k_x,k_y)$, posto efetivo e
coerência mútua da matriz de aquisição linearizada.}
\label{fig:geo_info}
\end{figure}

\begin{table}[H]
\centering\small
\caption{Varredura da separação, alvo no eixo: canal biestático isolado e
fusão a orçamento constante. Monoestático: $\delta_x=0,157\,$m, $V=19,2\,$l,
PSLR $-7,88\dB$, ISLR $6,37\dB$, contraste $21,3\dB$, posto efetivo $28,4$.}
\label{tab:geo_varredura}
\begin{tabular}{@{}r|rrr|rrrrrr@{}}
\toprule
 & \multicolumn{3}{c|}{\textbf{Biestático isolado}} & \multicolumn{6}{c}{\textbf{Fusão mono+bi}} \\
$\Delta\phi$ & $\delta_x$ & $V$ & ISLR & $\delta_x$ & $V$ & PSLR & ISLR & Contr. & Posto \\
\midrule
 15 & 0,158 & 19,7 & 6,36 & 0,158 & 19,6 & $-7,90$ & 6,35 & 21,3 & 28,3 \\
 30 & 0,162 & 20,8 & 6,31 & 0,160 & 20,2 & $-7,92$ & 6,21 & 21,4 & 28,1 \\
 45 & 0,169 & 22,6 & 6,19 & 0,163 & 21,0 & $-8,00$ & 5,71 & 21,7 & 28,2 \\
 60 & 0,180 & 25,6 & 5,87 & 0,167 & 21,9 & $-8,28$ & 4,54 & 22,6 & 28,9 \\
 75 & 0,199 & 29,4 & 5,58 & 0,174 & 23,4 & $-8,79$ & 2,53 & 24,2 & 30,4 \\
 90 & 0,221 & 38,8 & 5,10 & 0,182 & 25,9 & $-9,91$ & \textbf{1,04} & \textbf{25,3} & 32,4 \\
105 & 0,258 & 52,0 & 4,37 & 0,192 & 28,0 & $-12,19$ & 1,60 & 24,3 & \textbf{33,8} \\
120 & 0,315 & 75,9 & 3,45 & 0,204 & 30,3 & $\mathbf{-13,09}$ & 2,50 & 22,8 & 33,7 \\
135 & 0,412 & 131 & 1,62 & 0,214 & 34,8 & $-10,24$ & 2,32 & 22,2 & 31,5 \\
150 & 0,612 & 285 & $-1,13$ & 0,224 & 37,6 & $-9,48$ & 3,47 & 20,4 & 26,7 \\
165 & 1,216 & 1118 & $-9,6$ & 0,232 & 40,8 & $-5,25$ & 7,30 & 16,2 & 20,0 \\
180 & 17,2 & $>5331$ & --- & 0,235 & 41,1 & $-3,81$ & 11,95 & 11,4 & 16,0 \\
\bottomrule
\end{tabular}
\end{table}

\begin{figure}[H]
\centering
\includegraphics[width=\textwidth]{B_imagenes.pdf}
\caption{Imagens monoestática, biestática e fundida para
$\Delta\phi=45^\circ$, $90^\circ$, $135^\circ$ e $180^\circ$, com a mesma escala
(pico monoestático em $0\dB$, faixa de $30\dB$), numa janela de $2\times2\,$m em
torno do alvo. Os anéis são os lóbulos laterais da PSF do tipo $J_0$.}
\label{fig:geo_imagens}
\end{figure}

\paragraph{Lei geométrica.} A resolução horizontal do canal biestático segue
$\delta_{xy}^{\text{mono}}/|\cos(\Delta\phi/2)|$ com erro inferior a $1\,\%$ até
$150^\circ$; a vertical é a mesma em todos os casos ($1,386\,$m), porque a
extensão em $k_z$ é idêntica (desvio padrão de $1,174\,$rad/m em todas as
separações).

\paragraph{Melhor imagem biestática isolada.} É a de menor separação viável.
Toda separação degrada o lóbulo principal do canal isolado e não há ótimo
intermediário. Com uma separação mínima de segurança de $15\,$m entre as
plataformas ($\Delta\phi\gtrsim6^\circ$), o canal biestático mais nítido é,
na prática, um segundo canal monoestático e não acrescenta informação.

\paragraph{Maior complementaridade.} Ocorre para
$\Delta\phi\approx90^\circ$--$120^\circ$. O anel interior do canal biestático
preenche o disco do espaço de números de onda que o anel monoestático deixa
vazio: a área ocupada no plano $(k_x,k_y)$ passa de $294$ para
$475\,$rad$^2$/m$^2$ ($75^\circ$), e o posto efetivo da matriz de aquisição, de
$28,4$ para $33,8$ ($105^\circ$). Na imagem fundida, o ISLR cai $5,3\dB$
($90^\circ$), o PSLR melhora $5,2\dB$ ($120^\circ$) e o contraste sobe $4,0\dB$
($90^\circ$).

\paragraph{Custo.} O lóbulo principal horizontal alarga-se: $\delta_x$ aumenta
$16\,\%$ a $90^\circ$ e $30\,\%$ a $120^\circ$, e o volume, $35$ e $58\,\%$. A fusão
não melhora a resolução a $-3\dB$; melhora os lóbulos, o contraste e o
condicionamento. A anisotropia 3D diminui ($8,9\to7,4$ a $90^\circ$) apenas porque
o lóbulo horizontal se alarga.

\paragraph{Invariância na cena.} Para o alvo descentrado em $(6;\,4;\,-2)\,$m, os
resultados são os mesmos (diferenças inferiores a $0,3\dB$ no PSLR): com duas
voltas completas, a PSF é praticamente invariante na cena.

\section{Interpretação: Preenchimento do Disco de Números de Onda}
\label{sec:geo_interp}

A PSF horizontal de um anel de raio $R_c$ é $J_0(R_c\rho)$, cujos lóbulos
laterais decaem lentamente, como $\rho^{-1/2}$ --- daí o PSLR de $-7,9\dB$ e o
ISLR positivo da imagem monoestática. A fusão de dois anéis concêntricos de
raios $R_c$ e $R_c|\cos(\Delta\phi/2)|$ produz
$\tfrac12[J_0(R_c\rho)+J_0(R_c|\cos(\Delta\phi/2)|\rho)]$: os primeiros lóbulos
das duas funções de Bessel ocorrem em raios diferentes e se cancelam
parcialmente, enquanto o lóbulo principal se alarga por ser a média de dois
lóbulos de larguras distintas. O ótimo de lóbulos resulta de um compromisso
entre duas tendências: separações pequenas produzem anéis quase coincidentes,
sem cancelamento; separações grandes produzem um anel interior tão pequeno que
o canal biestático se torna um pedestal. Varrer $\Delta\phi$ ao longo da
abertura gera vários anéis e preenche o disco de forma mais uniforme, o que
motiva a trajetória do Capítulo~\ref{cap:trajetoria}.

\section{Separabilidade de Alvos Próximos}
\label{sec:geo_separabilidade}

Com alvos separados por $1\,$m, todos os canais, exceto o diametral, os resolvem.
A separabilidade foi então medida com dezesseis pontos num plano horizontal a
$2\,$m de profundidade, com espaçamento $s$ de $0,15$ a $0,30\,$m
($\delta_x^{\text{mono}}=0,157\,$m) e fases aleatórias fixas por ponto
(Tabela~\ref{tab:geo_sep}).

\begin{table}[H]
\centering\small
\caption{Fração de alvos resolvidos (queda de pelo menos $3\dB$ para cada
vizinho), média de três realizações.}
\label{tab:geo_sep}
\begin{tabular}{@{}lrrrr@{}}
\toprule
Canal & $s=0,15\,$m & $0,20\,$m & $0,25\,$m & $0,30\,$m \\
\midrule
Monoestático & $8\,\%$ & $8\,\%$ & $29\,\%$ & $\mathbf{73\,\%}$ \\
Mono + bi $45^\circ$ & $6\,\%$ & $2\,\%$ & $25\,\%$ & $63\,\%$ \\
Mono + bi $90^\circ$ & $6\,\%$ & $0\,\%$ & $10\,\%$ & $27\,\%$ \\
Mono + bi $120^\circ$ & $2\,\%$ & $2\,\%$ & $10\,\%$ & $13\,\%$ \\
Mono + bi $135^\circ$ & $2\,\%$ & $0\,\%$ & $13\,\%$ & $15\,\%$ \\
Biestático $90^\circ$ isolado & $0\,\%$ & $0\,\%$ & $4\,\%$ & $17\,\%$ \\
\bottomrule
\end{tabular}
\end{table}

\begin{figure}[H]
\centering
\includegraphics[width=\textwidth]{D2_separabilidad.pdf}
\caption{Separabilidade horizontal em função do espaçamento e imagens com
$s=0,20\,$m.}
\label{fig:geo_sep}
\end{figure}

A separação de alvos a cerca de duas larguras de resolução é governada pelo
lóbulo principal, e a fusão o alarga: a imagem monoestática separa melhor que
qualquer fusão, e a fusão a $90^\circ$--$135^\circ$ resolve de $2,7$ a $5,6$ vezes
menos alvos a $s=0,30\,$m. Os lóbulos laterais mais baixos da fusão não
compensam o alargamento quando os alvos estão próximos.

\begin{sintese}
O canal biestático à mesma altura não melhora nenhuma resolução: a horizontal
piora como $1/|\cos(\Delta\phi/2)|$ e a vertical não muda. Com
$\Delta\phi\approx90^\circ$--$120^\circ$, porém, ele preenche o disco do espaço
de números de onda, e a fusão reduz o ISLR em $5,3\dB$, melhora o PSLR em até
$5,2\dB$ e eleva o contraste em $4\dB$, com um lóbulo principal $16$ a $30\,\%$
mais largo. A configuração diametral piora a imagem monoestática, e a separação
de alvos próximos é melhor no canal monoestático.
\end{sintese}

%%═══════════════════════════════════════════════════════════════════════════════
\chapter{Alvos Reais: Anisotropia, Polarização e Orientação}
\label{cap:alvos}
%%═══════════════════════════════════════════════════════════════════════════════

\section{Introdução}
\label{sec:alv_intro}

O ganho de lóbulos e de contraste do Capítulo~\ref{cap:geometria} foi obtido com
um espalhador pontual isotrópico e escalar, cuja resposta não depende da
direção nem da polarização. Este capítulo verifica se ele sobrevive com alvos
condutores extensos, de RCS anisotrópica, em orientação nominal e em orientação
aleatória, e mostra que a resposta depende de uma propriedade do receptor que o
alvo pontual não revela: a polarização.

\section{Biblioteca de Alvos em Orientação Nominal}
\label{sec:alv_nominal}

Os alvos da Seção~\ref{sec:mod_alvos} foram centrados em $(0;\,0;\,-2)\,$m no
solo S1 e reconstruídos numa grade de $3\times3\times4\,$m. Para alvos extensos,
as métricas de PSF são avaliadas sobre a resposta do objeto, e acrescenta-se a
fração de energia dentro da região do objeto (caixa envolvente com margem de
$0,3\,$m).

\begin{figure}[H]
\centering
\includegraphics[width=0.92\textwidth]{D_mip_objetos.pdf}
\caption{Projeção de máxima intensidade, ao longo de $y$, de cada alvo: canais
monoestático e biestático e fusão para $\Delta\phi=90^\circ$ e $135^\circ$.
Escala comum por alvo (pico monoestático em $0\dB$, faixa de $30\dB$).}
\label{fig:alv_mip}
\end{figure}

\begin{table}[H]
\centering\small
\caption{Fusão a orçamento constante frente ao monoestático em orientação
nominal (receptor TM co-polar). ``Pico bi'': pico do canal biestático relativo ao
monoestático. Diferenças em dB, exceto a energia na região do objeto
($\Delta E$, pontos percentuais).}
\label{tab:alv_nominal}
\begin{tabular}{@{}lr|r|rrrrrr@{}}
\toprule
Alvo & $\Delta\phi$ & Pico bi & $\Delta V$ & $\Delta$PSLR & $\Delta$ISLR & $\Delta$Contr. & $\Delta$SNR & $\Delta E$ \\
\midrule
ponto & 90 & $+0,2$ & $+1,1$ & $\mathbf{-2,0}$ & $\mathbf{-4,0}$ & $\mathbf{+2,6}$ & $+0,1$ & $+4,1$ \\
 & 135 & $+0,3$ & $+2,4$ & $-2,4$ & $-2,7$ & $-0,4$ & $+0,1$ & $+0,8$ \\
esfera & 90 & $-14,9$ & $+0,2$ & $+0,7$ & $-1,5$ & $+1,3$ & $-4,8$ & $-0,2$ \\
 & 135 & $-3,5$ & $+0,9$ & $+2,6$ & $+3,0$ & $-4,3$ & $-5,8$ & $-8,5$ \\
cubo & 90 & $-19,3$ & $+0,1$ & $+0,3$ & $+0,9$ & $-0,8$ & $-6,7$ & $+1,0$ \\
 & 135 & $+12,6$ & $+11,7$ & $+3,1$ & $+2,5$ & $-4,8$ & $+5,1$ & $-6,3$ \\
caixa & 90 & $-10,5$ & $+1,2$ & $+0,5$ & $-0,5$ & $-0,7$ & $-6,3$ & $+1,5$ \\
 & 135 & $+17,4$ & $+8,5$ & $-2,8$ & $-7,6$ & $-1,1$ & $+11,7$ & $-8,2$ \\
cilindro vertical & 90 & $-16,9$ & $+0,1$ & $-1,0$ & $+0,2$ & $-0,5$ & $-6,3$ & $+0,4$ \\
 & 135 & $+16,7$ & $+5,7$ & $-2,0$ & $-3,8$ & $-3,3$ & $+11,9$ & $-11,7$ \\
tubo horizontal & 90 & $-19,6$ & $+0,2$ & $0,0$ & $-0,8$ & $+0,6$ & $-5,4$ & $+2,1$ \\
 & 135 & $+6,2$ & $+4,6$ & $-0,2$ & $-8,3$ & $+0,7$ & $+1,7$ & $-4,5$ \\
tubo inclinado & 90 & $-19,9$ & $+0,1$ & $+0,2$ & $-0,2$ & $+0,1$ & $-5,8$ & $+0,8$ \\
 & 135 & $-5,7$ & $+3,7$ & $+0,4$ & $+1,0$ & $-3,2$ & $-6,8$ & $-4,2$ \\
cilindro inclinado & 90 & $-16,9$ & $+0,2$ & $-0,1$ & $-0,9$ & $+0,7$ & $-5,1$ & $+1,3$ \\
 & 135 & $-4,2$ & $+3,1$ & $+3,2$ & $+0,6$ & $-3,3$ & $-5,7$ & $-5,1$ \\
\midrule
todos & 180 & \multicolumn{7}{l}{fusão sempre pior: $\Delta V=+6$ a $+28\dB$, $\Delta$contraste $-8$ a $-17\dB$} \\
\bottomrule
\end{tabular}
\end{table}

Três observações decorrem da Tabela~\ref{tab:alv_nominal}. \textbf{(i)} O ganho
de lóbulos é próprio do espalhador isotrópico: somente o ponto, e parcialmente
a esfera a $90^\circ$, o reproduz. Nos objetos com faces planas, a resposta é
dominada pelas reflexões especulares das faces, que já fixam o PSLR
monoestático entre $-0,5$ e $-4,7\dB$ e não dependem do preenchimento do espaço
de números de onda. \textbf{(ii)} A $90^\circ$, o eco biestático dos objetos é
$10$ a $20\dB$ mais fraco que o monoestático, inclusive na esfera, e a fusão de
pesos iguais resulta praticamente na imagem monoestática com $5$ a $7\dB$ menos de
SNR: \emph{a fusão de pesos iguais dilui o sinal}. \textbf{(iii)} A $135^\circ$ e
$180^\circ$ ocorre o oposto com cubo, caixa, cilindro vertical e tubo horizontal:
a bissetriz quase vertical mantém a especular da face superior, o canal
biestático domina ($+6$ a $+17\dB$) e a fusão herda seu lóbulo largo ($\Delta V$
de até $+12\dB$ a $135^\circ$ e $+28\dB$ a $180^\circ$). Mais eco biestático não
implica imagem melhor.

\section{O Descasamento Geométrico de Polarização}
\label{sec:alv_polarizacao}

A queda de $10$ a $20\dB$ do canal biestático a $90^\circ$ tem origem geométrica.
Dentro do solo, os raios formam cerca de $23^\circ$ com a vertical, de modo que o
vetor de polarização TM é quase horizontal e aponta radialmente para cada
plataforma. A $\Delta\phi=90^\circ$, o vetor TM do receptor é quase ortogonal ao
campo que as faces do objeto espalham na direção do receptor, e o acoplamento
co-polar escala aproximadamente com $\cos\Delta\phi$. O alvo pontual escalar
não tem polarização e não revela o efeito.

\begin{proposicao}[Descasamento co-polar]
\label{prop:descasamento}
Sejam $\theta$ o ângulo dos raios no solo com a vertical e $\Delta\phi$ a
separação azimutal entre transmissor e receptor à mesma altura. Para um
espalhador cuja matriz de espalhamento é aproximadamente proporcional à
identidade no plano transversal (reflexão especular de uma face), o acoplamento
entre os vetores TM do transmissor e do receptor é
\begin{equation}
  \hat v_T\cdot\hat v_R=\cos^2\theta\cos\Delta\phi+\sin^2\theta ,
\end{equation}
que, para $\theta\approx23^\circ$, cai de $1$ em $\Delta\phi=0$ para $0,15$ em
$\Delta\phi=90^\circ$ ($-16,3\dB$ em potência).
\end{proposicao}

\begin{proof}
No solo, os versores de polarização TM das duas plataformas, vistos do alvo,
são $\hat v_i=(\cos\theta\cos\phi'_i;\ \cos\theta\sin\phi'_i;\ -\sin\theta)$,
$i\in\{T,R\}$. O produto escalar é
$\cos^2\theta\cos(\phi'_R-\phi'_T)+\sin^2\theta$. Com $\theta=23^\circ$,
$\sin^2\theta=0,153$.
\end{proof}

O modelo simplificado reproduz a queda a $45^\circ$ ($-2,5\dB$, contra $-2,4$ a
$-2,6\dB$ observados com orientação aleatória) e a ordem de grandeza da queda a
$90^\circ$ ($-16\dB$, contra $10$ a $20\dB$). Para separações grandes ele deixa de
valer, porque a resposta passa a ser dominada pela reflexão especular ao longo
da bissetriz, cuja polarização não é a do plano de incidência de cada
plataforma.
Avaliou-se então o receptor alinhado de dupla polarização
(Seção~\ref{sec:mod_alvos}), que recebe a componente do campo do transmissor
transversal à direção de recepção.

\begin{figure}[H]
\centering
\includegraphics[width=\textwidth]{D3_rx_alineado.pdf}
\caption{Pico do canal biestático relativo ao monoestático com receptor TM
co-polar e com receptor alinhado de dupla polarização (orientação nominal).}
\label{fig:alv_alinhado}
\end{figure}

\begin{table}[H]
\centering\small
\caption{Receptor TM co-polar frente ao receptor alinhado a
$\Delta\phi=90^\circ$ (fusão a orçamento constante, diferenças em relação ao
monoestático).}
\label{tab:alv_alinhado}
\begin{tabular}{@{}lrr|rr|rr@{}}
\toprule
 & \multicolumn{2}{c|}{Pico bi [dB]} & \multicolumn{2}{c|}{$\Delta$SNR fusão [dB]} & \multicolumn{2}{c}{$\Delta$contraste fusão [dB]} \\
Alvo & co-polar & alinhado & co-polar & alinhado & co-polar & alinhado \\
\midrule
esfera & $-14,9$ & $-1,9$ & $-4,8$ & $-1,5$ & $+1,3$ & $+2,0$ \\
cubo & $-19,3$ & $+8,4$ & $-6,7$ & $+0,2$ & $-0,8$ & $-3,7$ \\
caixa & $-10,5$ & $+9,3$ & $-6,3$ & $+3,6$ & $-0,7$ & $-0,5$ \\
cilindro vertical & $-16,9$ & $+9,8$ & $-6,3$ & $+0,5$ & $-0,5$ & $-5,5$ \\
tubo horizontal & $-19,6$ & $+1,8$ & $-5,4$ & $+0,7$ & $+0,6$ & $+2,3$ \\
tubo inclinado & $-19,9$ & $-2,7$ & $-5,8$ & $-1,2$ & $+0,1$ & $+3,0$ \\
cilindro inclinado & $-16,9$ & $+0,5$ & $-5,1$ & $-0,1$ & $+0,7$ & $+2,1$ \\
\bottomrule
\end{tabular}
\end{table}

Com o receptor alinhado, o canal biestático a $90^\circ$ recupera de $13$ a
$28\dB$, e a fusão deixa de diluir a SNR ($-1,5$ a $+3,6\dB$, contra $-4,8$ a
$-6,7\dB$). O contraste da fusão melhora de $2$ a $3\dB$ na esfera, nos tubos e no
cilindro inclinado; nos objetos com uma face horizontal, a especular biestática
passa a dominar e o contraste piora. \textbf{Consequência de projeto: o receptor
remoto deve ser de dupla polarização, ou girar sua antena com $\Delta\phi$.} Um
receptor TM co-polar anula a complementaridade exatamente na faixa de
separações que o espaço de números de onda recomenda.

\section{Orientação Aleatória}
\label{sec:alv_montecarlo}

A orientação de um alvo enterrado é, em geral, desconhecida. Num Monte Carlo
\emph{por aquisição}, sorteia-se em cada repetição uma rotação uniforme em
$SO(3)$, gerada por um quatérnion gaussiano normalizado~\citar{shoemake1992};
a mesma rotação é aplicada ao alvo durante toda a abertura e nos dois canais
(desenho pareado). Avaliaram-se cubo (1000 orientações), caixa, cilindro e tubo
(200 cada), para $\Delta\phi\in\{45^\circ,90^\circ,135^\circ,180^\circ\}$, com
800 pulsos e os dois receptores. As medianas e as probabilidades de melhoria
foram estimadas com intervalos de confiança por \textit{bootstrap}~\citar{efron1993}.

\begin{table}[H]
\centering\small
\caption{Monte Carlo em $SO(3)$ por aquisição com receptor TM co-polar: mediana
de (fusão a orçamento constante $-$ monoestático) e probabilidade de melhoria
$P_{\text{mel}}$.}
\label{tab:alv_mc_tm}
\resizebox{\textwidth}{!}{\begin{tabular}{@{}lrr|rrrrrrrr@{}}
\toprule
Alvo & $\Delta\phi$ & $\bar P_{bi}/\bar P_{mono}$ & \multicolumn{2}{c}{$\Delta$SNR [dB]} & \multicolumn{2}{c}{$\Delta$Contr. [dB]} & \multicolumn{2}{c}{$\Delta$PSLR [dB]} & \multicolumn{2}{c}{$\Delta V$ [dB]} \\
 & & [dB] & med. & $P_{\text{mel}}$ & med. & $P_{\text{mel}}$ & med. & $P_{\text{mel}}$ & med. & $P_{\text{mel}}$ \\
\midrule
cubo & $45^\circ$ & $-2,6$ & $-1,4$ & $7\,\%$ & $+0,3$ & $81\,\%$ & $+0,00$ & $49\,\%$ & $+0,1$ & $38\,\%$ \\
cubo & $90^\circ$ & $-20,2$ & $-5,7$ & $0\,\%$ & $+0,3$ & $78\,\%$ & $+0,02$ & $47\,\%$ & $+0,1$ & $39\,\%$ \\
cubo & $135^\circ$ & $-3,8$ & $-4,4$ & $13\,\%$ & $-0,0$ & $48\,\%$ & $+0,07$ & $42\,\%$ & $+0,9$ & $43\,\%$ \\
cubo & $180^\circ$ & $-0,0$ & $-2,9$ & $31\,\%$ & $-2,8$ & $10\,\%$ & $+0,11$ & $34\,\%$ & $-0,4$ & $53\,\%$ \\
\midrule
caixa & $45^\circ$ & $-2,4$ & $-1,2$ & $2\,\%$ & $+0,4$ & $98\,\%$ & $-0,06$ & $71\,\%$ & $+0,1$ & $33\,\%$ \\
caixa & $90^\circ$ & $-19,8$ & $-5,4$ & $0\,\%$ & $+0,6$ & $91\,\%$ & $-0,00$ & $52\,\%$ & $+0,1$ & $40\,\%$ \\
caixa & $135^\circ$ & $-5,2$ & $-4,8$ & $10\,\%$ & $-0,5$ & $37\,\%$ & $+0,06$ & $46\,\%$ & $+0,6$ & $44\,\%$ \\
caixa & $180^\circ$ & $-2,9$ & $-2,4$ & $28\,\%$ & $-5,2$ & $3\,\%$ & $+0,19$ & $34\,\%$ & $+0,7$ & $43\,\%$ \\
\midrule
cilindro & $45^\circ$ & $-2,6$ & $-1,1$ & $0\,\%$ & $+0,5$ & $100\,\%$ & $-0,00$ & $50\,\%$ & $+0,2$ & $25\,\%$ \\
cilindro & $90^\circ$ & $-20,3$ & $-5,5$ & $0\,\%$ & $+0,6$ & $99\,\%$ & $-0,06$ & $66\,\%$ & $+0,1$ & $30\,\%$ \\
cilindro & $135^\circ$ & $-3,1$ & $-5,1$ & $7\,\%$ & $-1,3$ & $21\,\%$ & $+0,48$ & $26\,\%$ & $+3,5$ & $22\,\%$ \\
cilindro & $180^\circ$ & $+0,9$ & $-1,6$ & $28\,\%$ & $-7,2$ & $0\,\%$ & $+0,74$ & $19\,\%$ & $+4,4$ & $26\,\%$ \\
\midrule
tubo & $45^\circ$ & $-2,4$ & $-1,1$ & $3\,\%$ & $+0,4$ & $96\,\%$ & $-0,00$ & $50\,\%$ & $+0,2$ & $11\,\%$ \\
tubo & $90^\circ$ & $-19,5$ & $-5,4$ & $0\,\%$ & $+0,6$ & $94\,\%$ & $+0,02$ & $46\,\%$ & $+0,1$ & $23\,\%$ \\
tubo & $135^\circ$ & $-3,7$ & $-5,6$ & $16\,\%$ & $-2,0$ & $28\,\%$ & $+0,42$ & $27\,\%$ & $+2,3$ & $24\,\%$ \\
tubo & $180^\circ$ & $-0,1$ & $-2,6$ & $26\,\%$ & $-7,9$ & $0\,\%$ & $+0,51$ & $22\,\%$ & $+4,5$ & $31\,\%$ \\
\bottomrule
\end{tabular}}
\end{table}

\begin{table}[H]
\centering\small
\caption{Monte Carlo em $SO(3)$ por aquisição com receptor alinhado de dupla
polarização (200 orientações por alvo).}
\label{tab:alv_mc_al}
\resizebox{\textwidth}{!}{\begin{tabular}{@{}lrr|rrrrrrrr@{}}
\toprule
Alvo & $\Delta\phi$ & $\bar P_{bi}/\bar P_{mono}$ & \multicolumn{2}{c}{$\Delta$SNR [dB]} & \multicolumn{2}{c}{$\Delta$Contr. [dB]} & \multicolumn{2}{c}{$\Delta$PSLR [dB]} & \multicolumn{2}{c}{$\Delta V$ [dB]} \\
 & & [dB] & med. & $P_{\text{mel}}$ & med. & $P_{\text{mel}}$ & med. & $P_{\text{mel}}$ & med. & $P_{\text{mel}}$ \\
\midrule
cubo & $45^\circ$ & $-0,5$ & $-0,5$ & $25\,\%$ & $+0,3$ & $80\,\%$ & $+0,02$ & $46\,\%$ & $+0,2$ & $35\,\%$ \\
cubo & $90^\circ$ & $-1,4$ & $-2,1$ & $24\,\%$ & $\mathbf{+1,8}$ & $\mathbf{90\,\%}$ & $-0,05$ & $53\,\%$ & $+0,8$ & $33\,\%$ \\
cubo & $135^\circ$ & $-0,5$ & $-2,9$ & $30\,\%$ & $+1,2$ & $78\,\%$ & $-0,06$ & $59\,\%$ & $+0,8$ & $40\,\%$ \\
\midrule
caixa & $45^\circ$ & $-0,6$ & $-0,4$ & $22\,\%$ & $+0,5$ & $94\,\%$ & $-0,04$ & $62\,\%$ & $+0,1$ & $28\,\%$ \\
caixa & $90^\circ$ & $-1,6$ & $-1,5$ & $20\,\%$ & $\mathbf{+2,6}$ & $\mathbf{98\,\%}$ & $-0,28$ & $70\,\%$ & $+0,6$ & $30\,\%$ \\
caixa & $135^\circ$ & $-1,0$ & $-2,7$ & $21\,\%$ & $+1,0$ & $84\,\%$ & $-0,36$ & $70\,\%$ & $+0,7$ & $33\,\%$ \\
\midrule
cilindro & $45^\circ$ & $-0,5$ & $-0,2$ & $29\,\%$ & $+0,5$ & $100\,\%$ & $-0,02$ & $57\,\%$ & $+0,2$ & $24\,\%$ \\
cilindro & $90^\circ$ & $-1,6$ & $-0,9$ & $30\,\%$ & $\mathbf{+2,7}$ & $\mathbf{96\,\%}$ & $-0,39$ & $66\,\%$ & $+0,8$ & $23\,\%$ \\
cilindro & $135^\circ$ & $-1,6$ & $-1,8$ & $26\,\%$ & $+0,9$ & $80\,\%$ & $-0,59$ & $76\,\%$ & $+1,4$ & $24\,\%$ \\
\midrule
tubo & $45^\circ$ & $-0,5$ & $-0,3$ & $36\,\%$ & $+0,5$ & $98\,\%$ & $-0,04$ & $64\,\%$ & $+0,3$ & $15\,\%$ \\
tubo & $90^\circ$ & $-1,0$ & $-1,1$ & $39\,\%$ & $\mathbf{+2,6}$ & $\mathbf{98\,\%}$ & $-0,08$ & $65\,\%$ & $+1,0$ & $10\,\%$ \\
tubo & $135^\circ$ & $-0,0$ & $-1,6$ & $35\,\%$ & $+1,0$ & $92\,\%$ & $-0,07$ & $58\,\%$ & $+1,5$ & $18\,\%$ \\
\bottomrule
\end{tabular}}
\end{table}

\begin{figure}[H]
\centering
\includegraphics[width=\textwidth]{E2_montecarlo_rx_alineado.pdf}
\caption{Monte Carlo em $SO(3)$ com receptor alinhado: distribuição de (fusão
$-$ monoestático) por alvo e separação; os percentuais são as probabilidades de
melhoria.}
\label{fig:alv_mc_al}
\end{figure}

\paragraph{Potência (controle).} Com orientação aleatória, a potência média do
canal biestático a $180^\circ$ é a mesma do monoestático ($-2,9$ a $+0,9\dB$
conforme o alvo), como no Capítulo~\ref{cap:potencia}. Com receptor co-polar,
cai cerca de $20\dB$ a $90^\circ$ e de $2,4$ a $5,2\dB$ a $45^\circ$ e $135^\circ$,
aproximadamente como $\cos^2\Delta\phi$ --- o descasamento da
Proposição~\ref{prop:descasamento}. Com o receptor alinhado, fica entre $0$ e
$-1,6\dB$ do monoestático em todas as separações.

\paragraph{Receptor co-polar.} A $45^\circ$--$90^\circ$, a fusão ganha um contraste
pequeno, porém sistemático ($+0,3$ a $+0,6\dB$, com melhoria em $78$ a $100\,\%$ das
aquisições), e perde de $1,1$ a $5,7\dB$ de SNR. A $135^\circ$--$180^\circ$, piora o
contraste (até $-7,9\dB$) e o volume.

\paragraph{Receptor alinhado.} A $90^\circ$, o contraste melhora $1,8$ a $2,7\dB$
com probabilidade de $90$ a $98\,\%$ nos quatro alvos; o PSLR melhora ligeiramente
(mediana de $-0,05$ a $-0,39\dB$, $53$ a $70\,\%$), a perda de SNR cai para $0,9$ a
$2,1\dB$ e a energia na região do objeto aumenta. O preço é um lóbulo $0,6$ a
$1,0\dB$ mais volumoso. A $135^\circ$, o ganho de contraste reduz-se a $0,9$ a
$1,2\dB$. Em nenhum caso a resolução ou a localização melhoram (medianas de
diferença inferiores a $1\,$cm).

\section{Mascaramento em Cenas com Vários Objetos}
\label{sec:alv_mascaramento}

\begin{figure}[H]
\centering
\includegraphics[width=\textwidth]{D_escenaP2.pdf}
\caption{Cena com cubo, tubo inclinado em azimute e esfera de RCS distintas:
projeção de máxima intensidade ao longo de $z$, mesma escala.}
\label{fig:alv_cena}
\end{figure}

Numa cena com cubo, tubo inclinado e esfera, a esfera domina a imagem
monoestática e o tubo fica $18,1\dB$ abaixo, sob os lóbulos dos demais objetos
(relação mínima objeto--lóbulo de $-17,5\dB$). Com a fusão a $135^\circ$, a faixa
dinâmica entre objetos reduz-se de $18,1$ para $7,9\dB$, e o objeto mais fraco
melhora $10,1\dB$ em relação aos lóbulos ($-7,4\dB$). A diversidade de aspecto do
canal biestático reduz o mascaramento entre objetos de RCS muito distintas; é o
benefício de complementaridade mais nítido observado com alvos reais. A
$90^\circ$, a fusão com receptor co-polar não altera a cena ($-17,7\dB$).

\begin{sintese}
Com alvos reais, a fusão só ajuda quando o canal biestático chega com potência
comparável à do monoestático. Um receptor TM co-polar perde cerca de $20\dB$ a
$90^\circ$ por descasamento geométrico de polarização; um receptor alinhado de
dupla polarização recupera de $13$ a $28\dB$ e torna sistemática a melhoria de
contraste com orientação aleatória ($+1,8$ a $+2,7\dB$ em $90$ a $98\,\%$ das
orientações). A diversidade de aspecto reduz o mascaramento entre objetos em
$10\dB$. A resolução e a localização não melhoram em nenhum caso.
\end{sintese}

%%═══════════════════════════════════════════════════════════════════════════════
\chapter{Retorno da Superfície, Pseudo-Brewster e o Nulo Biestático}
\label{cap:superficie}
%%═══════════════════════════════════════════════════════════════════════════════

\section{Introdução}
\label{sec:sup_intro}

O Capítulo~\ref{cap:potencia} concluiu que o ângulo de Brewster só pode ser útil
como mecanismo de redução do retorno da interface, e não de aumento do eco do
alvo. Este capítulo testa essa possibilidade nas condições de um sistema real
--- solo com perdas, rugosidade, erros de incidência e de atitude, banda
finita --- e responde a duas perguntas: a operação próxima do pseudo-Brewster
reduz o retorno da superfície de forma robusta? Qual é o custo, em qualidade de
imagem, de impor essa restrição à trajetória? A métrica principal é a relação
alvo--clutter (SCR); como a potência média do alvo não varia sistematicamente
com a geometria, qualquer melhoria de SCR deve ser interpretada como redução do
retorno da superfície.

\section{Delineamento}
\label{sec:sup_delineamento}

Comparam-se duas famílias de trajetória do transmissor, construídas ponto a
ponto a partir da hélice de referência. Cada pulso conserva seu azimute e sua
distância ao centro da cena (praticamente constante); muda apenas a incidência
$\theta$ vista do centro da superfície.
\begin{itemize}
  \item \textbf{Família helicoidal}: cada ponto é girado no seu plano vertical de
  $\theta_c-\langle\theta_0\rangle$; a trajetória continua sendo uma hélice e
  conserva sua própria excursão de incidência ($\theta_c\pm7^\circ$).
  \item \textbf{Família de incidência constante}: cada ponto é colocado exatamente
  em $\theta_c$; a trajetória torna-se uma órbita de duas voltas a altura e raio
  constantes (no pseudo-Brewster do solo S1, $85\,$m de altura e $169\,$m de
  raio), isto é, um círculo sobre o cone de incidência $\theta_c$.
\end{itemize}
O receptor acompanha o transmissor à mesma altura e no mesmo raio, desfasado
de $\Delta\phi\in\{0,\,90^\circ,\,135^\circ,\,180^\circ\}$. Variam-se a incidência
central ($35^\circ$, $45^\circ$, $53,4^\circ$, $\theta_{pB}$ e $75^\circ$), o erro
angular ($0$, $\pm0,5^\circ$, $\pm1^\circ$, $\pm2^\circ$, $\pm5^\circ$), a polarização
(TM e TE, co-polares), o rolamento da antena ($0$ a $5^\circ$), o solo (S0 a S3) e a
rugosidade ($s=0$, $1$, $2$ e $4\,$cm). O SCR é medido no voxel do alvo (ponto de
$1\,$m$^2$ a $2\,$m de profundidade), com a superfície congelada e compartilhada
pelos dois canais, três realizações, janela de Hann e janela de recepção
(Seção~\ref{ssec:bp_algo}). Comparar TM com TE na mesma geometria separa o efeito
de Brewster, exclusivo da TM, dos efeitos puramente geométricos.

\begin{figure}[H]
\centering
\includegraphics[width=\textwidth]{C_trayectorias_3d.pdf}
\caption{Trajetórias do transmissor. (a) Hélice de referência, entre os cones
de incidência de $46,3^\circ$ e $60,5^\circ$. (b) Família helicoidal no
pseudo-Brewster: a mesma hélice girada para baixo, que cruza o cone de
$\theta_{pB}$ a meio da abertura. (c) Família de incidência constante: todos os
pulsos sobre o cone de $63,45^\circ$.}
\label{fig:sup_traj}
\end{figure}

\begin{figure}[H]
\centering
\includegraphics[width=\textwidth]{C_geometria.pdf}
\caption{Geometria do estudo. (a) Plano vertical $(\rho,z)$ das famílias. (b)
Incidência do transmissor em cada pulso. (c) Planta da órbita de incidência
constante, com o receptor a $90^\circ$, $135^\circ$ e $180^\circ$ e o ponto
especular de cada par. (d) Incidência no ponto especular em função de
$\Delta\phi$: expressão~\eqref{eq:theta_esp} e simulação.}
\label{fig:sup_geom}
\end{figure}

A Fig.~\ref{fig:sup_geom}(d) traduz a expressão~\eqref{eq:theta_esp}: a
reflexão especular não ocorre na incidência das plataformas, mas no ponto
especular do par Tx--Rx. Só na configuração diametral as duas coincidem.
\textbf{Escolher $\theta_c=\theta_{pB}$ coloca a especular no pseudo-Brewster
apenas na configuração diametral}; nas demais, o que a incidência $\theta_c$
controla é o retorno difuso.

\section{Robustez da Reflectância}
\label{sec:sup_enlace}

\begin{table}[H]
\centering\small
\caption{Reflectância TM efetiva [dB] em $\theta_{pB}+\delta$. Rolamento nulo: média
de $\pm\delta$ (simétrica a menos de $1\dB$); rolamento de $5^\circ$: pares
$+\delta/-\delta$. Última coluna: incidência de $53,4^\circ$. Cores segundo a supressão
em relação à incidência de $53,4^\circ$ no pior sinal de $\delta$:
\colorbox{melhor}{$\ge10\dB$}, \colorbox{regular}{$5$--$10\dB$}, \colorbox{pior}{$<5\dB$}.}
\label{tab:sup_refl}
\resizebox{\textwidth}{!}{\begin{tabular}{@{}lrrrrr|r@{}}
\toprule
Solo / rolamento & $\delta=0$ & $\pm0,5^{\circ}$ & $\pm1^{\circ}$ & $\pm2^{\circ}$ & $\pm5^{\circ}$ & $53,4^\circ$ \\
\midrule
S0, $0^\circ$ & \cellcolor{melhor}$-189,7$ & \cellcolor{melhor}$-41,7$ & \cellcolor{melhor}$-35,7$ & \cellcolor{melhor}$-29,7$ & \cellcolor{pior}$-21,7$ & $-17,5$ \\
S0, $5^\circ$ & \cellcolor{melhor}$-46,8$ & \cellcolor{melhor}$-37,8$/$-49,1$ & \cellcolor{melhor}$-33,4$/$-38,9$ & \cellcolor{melhor}$-28,2$/$-31,5$ & \cellcolor{pior}$-20,3$/$-23,3$ & --- \\
S1, $0^\circ$ & \cellcolor{melhor}$-34,5$ & \cellcolor{melhor}$-33,8$ & \cellcolor{melhor}$-32,1$ & \cellcolor{melhor}$-28,5$ & \cellcolor{pior}$-21,5$ & $-17,5$ \\
S1, $5^\circ$ & \cellcolor{melhor}$-34,4$ & \cellcolor{melhor}$-32,9$/$-34,5$ & \cellcolor{melhor}$-31,0$/$-33,3$ & \cellcolor{regular}$-27,3$/$-29,8$ & \cellcolor{pior}$-20,1$/$-23,0$ & --- \\
S2 ($\theta_{pB}=68,66^{\circ}$) & \cellcolor{melhor}$-30,7$ & & & & & $-12,6$ \\
S3 ($\theta_{pB}=74,05^{\circ}$) & \cellcolor{melhor}$-26,5$ & & & & & $-8,7$ \\
\bottomrule
\end{tabular}}
\end{table}

No solo sem perdas, o mínimo é um zero exato que \textbf{desaparece com
$0,5^\circ$} de erro ($-190\to-42\dB$): é um mínimo estreito demais para
justificar uma trajetória real. Com perdas, o mínimo é \textbf{raso, porém
largo}: em S1, apenas $-34,5\dB$, perdendo $0,7$, $2,4$, $6,0$ e $13,0\dB$ com
erros de $0,5^\circ$, $1^\circ$, $2^\circ$ e $5^\circ$. Em relação à incidência de
$53,4^\circ$, a supressão disponível é de $17\dB$. O mínimo torna-se mais raso com a
umidade ($-30,7$ e $-26,5\dB$ em S2 e S3) e desloca-se para $68,7^\circ$ e $74,0^\circ$.
O rolamento de $1$ a $5^\circ$ é irrelevante em S1 (variação inferior a $1,5\dB$).

\section{O Nulo Biestático do Espalhamento Difuso}
\label{sec:sup_nulo}

A expressão~\eqref{eq:alpha_pp} mostra que, para o retorno difuso, o
pseudo-Brewster não é um ponto mas uma \emph{curva} no espaço
(incidência, separação).

\begin{proposicao}[Nulo biestático VV de primeira ordem]
\label{prop:nulo_vv}
Sejam o transmissor e o receptor à mesma elevação $\theta$ vista de uma faceta
da superfície, separados pela distância azimutal $\Delta\phi$, e seja
$\varepsilon$ real. O coeficiente $\alpha_{vv}$ do espalhamento difuso de
primeira ordem anula-se se, e somente se,
\begin{equation}
  \boxed{\cos\Delta\phi^\ast = -\frac{\varepsilon\sin^2\theta}{\varepsilon-\sin^2\theta}},
  \label{eq:nulo_vv}
\end{equation}
o que admite solução real se, e somente se, $\theta\le\theta_B$. A curva
$\Delta\phi^\ast(\theta)$ é monótona, parte de $90^\circ$ em incidência normal,
termina no ângulo de Brewster clássico em $\Delta\phi^\ast=180^\circ$ e nunca
passa pela configuração monoestática. O coeficiente $\alpha_{hh}$ anula-se em
$\Delta\phi=90^\circ$ para qualquer $\theta$ e qualquer $\varepsilon$.
\end{proposicao}

\begin{proof}
Com $\theta_i=\theta_s=\theta$, o numerador de $\alpha_{vv}$
em~\eqref{eq:alpha_pp} é
$(\varepsilon-1)[\varepsilon\sin^2\theta-\cos\phi_s\,(\varepsilon-\sin^2\theta)]$,
que se anula para $\cos\phi_s=\varepsilon\sin^2\theta/(\varepsilon-\sin^2\theta)$.
O azimute de espalhamento é medido a partir da direção especular, que fica a
$180^\circ$ do transmissor; logo $\phi_s=\Delta\phi-\pi$ e
$\cos\phi_s=-\cos\Delta\phi$, o que fornece~\eqref{eq:nulo_vv}. A solução existe
se o lado direito tem módulo menor ou igual a $1$, isto é,
$\varepsilon\sin^2\theta\le\varepsilon-\sin^2\theta$, ou
$\tan^2\theta\le\varepsilon=\tan^2\theta_B$. Em $\theta=\theta_B$, o lado direito vale
$-1$ e $\Delta\phi^\ast=180^\circ$ (espalhamento na direção especular, onde
$\alpha_{vv}$ se reduz ao numerador de $r_{\text{TM}}$); em $\theta\to0$, vale
$0$ e $\Delta\phi^\ast\to90^\circ$. O lado direito é crescente em $\sin^2\theta$,
o que dá a monotonicidade; como é sempre negativo, $\Delta\phi^\ast>90^\circ$ e a
curva nunca atinge $\Delta\phi=0$. Para $\alpha_{hh}\propto\cos\phi_s$, o zero
ocorre em $\phi_s=\pm90^\circ$, ou seja, $\Delta\phi=90^\circ$.
\end{proof}

A Proposição~\ref{prop:nulo_vv} generaliza o ângulo de Brewster ao espalhamento
biestático difuso: o zero de reflexão especular é o extremo $\Delta\phi=180^\circ$ de
uma família contínua de nulos difusos. Com $\varepsilon\approx\mathrm{Re}\,\varepsilon_r$,
na incidência média da hélice ($53,4^\circ$), a expressão~\eqref{eq:nulo_vv} fornece
$\Delta\phi^\ast=140,2^\circ$ (S0, S1), $135,7^\circ$ (S2) e $132,9^\circ$ (S3),
em acordo com os mínimos numéricos de $140^\circ$, $136^\circ$ e $133^\circ$
(Fig.~\ref{fig:sup_sigma0}). Nos solos com perdas, o nulo fica de $34$ a $39\dB$
abaixo do retroespalhamento e mantém-se a menos de $10\dB$ do mínimo numa faixa
de $4$ a $8^\circ$ de $\Delta\phi$. O nulo HH, ao contrário, é puramente geométrico:
é o desacoplamento entre os vetores TE do transmissor e do receptor, ortogonais a
$90^\circ$.

\begin{figure}[H]
\centering
\includegraphics[width=\textwidth]{C_sigma0_mapa_suelos.pdf}
\caption{Coeficiente de espalhamento biestático de primeira ordem com o receptor
à mesma elevação do transmissor, para os quatro solos ($s=2\,$cm). Acima VV (TM),
abaixo HH (TE). Curva vermelha: nulo VV~\eqref{eq:nulo_vv}; vertical vermelha:
nulo HH. Tracejado: $\theta_{pB}$ de cada solo; faixa clara: incidências da
hélice de referência.}
\label{fig:sup_sigma0}
\end{figure}

\section{Relação Alvo--Clutter na Imagem}
\label{sec:sup_scr}

\begin{table}[H]
\centering\small
\caption{SCR na imagem [dB], no voxel do alvo pontual (S1, $s=2\,$cm, três
realizações). \colorbox{melhor}{\textbf{verde}}: melhor canal da linha;
\colorbox{pior}{vermelho}: SCR $<0$. $\leftarrow$: linha no pseudo-Brewster.}
\label{tab:sup_scr}
\resizebox{\textwidth}{!}{\begin{tabular}{@{}llrrrrrr@{}}
\toprule
Pol. & Família, $\theta_c$ & Mono & Bi $90^\circ$ & Bi $135^\circ$ & Diametral & Mono+bi $90^\circ$ & Mono+bi $135^\circ$ \\
\midrule
TM & helicoidal, $35^\circ$ & \cellcolor{pior}$-0,5$ & 3,3 & 0,8 & \cellcolor{pior}$-59,6$ & 2,8 & \cellcolor{melhor}\textbf{5,8} \\
 & helicoidal, $45^\circ$ & \cellcolor{pior}$-0,2$ & 1,3 & \cellcolor{melhor}\textbf{13,0} & \cellcolor{pior}$-58,2$ & 4,4 & 5,6 \\
 & helicoidal, $53,4^\circ$ & 3,1 & 3,1 & \cellcolor{melhor}\textbf{9,8} & \cellcolor{pior}$-55,3$ & 4,3 & 9,5 \\
 & \textbf{helicoidal, $\theta_{pB}$} $\leftarrow$ & 8,0 & 6,0 & 7,4 & \cellcolor{pior}$-41,5$ & 9,4 & \cellcolor{melhor}\textbf{12,0} \\
 & helicoidal, $75^\circ$ & 14,8 & 11,6 & 12,4 & \cellcolor{pior}$-66,3$ & 15,6 & \cellcolor{melhor}\textbf{18,0} \\
\cmidrule(l){2-8}
 & inc.\ constante, $35^\circ$ & \cellcolor{pior}$-4,7$ & \cellcolor{pior}$-1,6$ & \cellcolor{pior}$-13,4$ & \cellcolor{pior}$-61,2$ & \cellcolor{melhor}\textbf{$-0,2$} & \cellcolor{pior}$-8,3$ \\
 & inc.\ constante, $45^\circ$ & 0,5 & \cellcolor{pior}$-0,9$ & 3,4 & \cellcolor{pior}$-60,1$ & 1,6 & \cellcolor{melhor}\textbf{5,6} \\
 & inc.\ constante, $53,4^\circ$ & 0,2 & 3,0 & \cellcolor{melhor}\textbf{6,9} & \cellcolor{pior}$-57,3$ & 2,9 & 4,2 \\
 & \textbf{inc.\ constante, $\theta_{pB}$} $\leftarrow$ & 2,5 & 3,5 & 5,1 & \cellcolor{pior}$-41,5$ & \cellcolor{melhor}\textbf{7,8} & 7,3 \\
 & inc.\ constante, $75^\circ$ & 8,4 & 5,7 & 8,3 & \cellcolor{pior}$-66,0$ & \cellcolor{melhor}\textbf{12,6} & 12,1 \\
\midrule
TE & helicoidal, $35^\circ$ & 1,4 & \cellcolor{melhor}\textbf{14,4} & 1,5 & \cellcolor{pior}$-64,4$ & 6,7 & 4,2 \\
 & helicoidal, $45^\circ$ & 2,4 & \cellcolor{melhor}\textbf{18,6} & \cellcolor{pior}$-4,9$ & \cellcolor{pior}$-66,9$ & 8,8 & 0,7 \\
 & helicoidal, $53,4^\circ$ & 6,2 & \cellcolor{melhor}\textbf{23,8} & \cellcolor{pior}$-2,3$ & \cellcolor{pior}$-69,8$ & 11,6 & 4,6 \\
 & \textbf{helicoidal, $\theta_{pB}$} $\leftarrow$ & 11,5 & \cellcolor{melhor}\textbf{28,9} & \cellcolor{pior}$-1,3$ & \cellcolor{pior}$-74,4$ & 17,1 & 4,7 \\
 & helicoidal, $75^\circ$ & 18,5 & \cellcolor{melhor}\textbf{35,0} & 6,9 & \cellcolor{pior}$-81,3$ & 24,6 & 13,6 \\
\cmidrule(l){2-8}
 & inc.\ constante, $35^\circ$ & \cellcolor{pior}$-2,6$ & \cellcolor{melhor}\textbf{11,7} & \cellcolor{pior}$-18,1$ & \cellcolor{pior}$-65,8$ & 3,6 & \cellcolor{pior}$-12,0$ \\
 & inc.\ constante, $45^\circ$ & 3,5 & \cellcolor{melhor}\textbf{18,5} & \cellcolor{pior}$-9,1$ & \cellcolor{pior}$-68,6$ & 9,0 & \cellcolor{pior}$-3,3$ \\
 & inc.\ constante, $53,4^\circ$ & 3,8 & \cellcolor{melhor}\textbf{24,6} & \cellcolor{pior}$-3,9$ & \cellcolor{pior}$-71,5$ & 9,5 & 0,2 \\
 & \textbf{inc.\ constante, $\theta_{pB}$} $\leftarrow$ & 6,7 & \cellcolor{melhor}\textbf{28,7} & \cellcolor{pior}$-1,9$ & \cellcolor{pior}$-75,5$ & 12,9 & 4,2 \\
 & inc.\ constante, $75^\circ$ & 13,1 & \cellcolor{melhor}\textbf{36,0} & 4,4 & \cellcolor{pior}$-81,6$ & 19,4 & 10,2 \\
\bottomrule
\end{tabular}}
\end{table}

\begin{figure}[H]
\centering
\includegraphics[width=\textwidth]{C_scr_incidencia_rugosa.pdf}
\caption{SCR no voxel do alvo em função da incidência central (S1, $s=2\,$cm),
famílias helicoidal e de incidência constante, TM e TE. Linha vermelha:
pseudo-Brewster.}
\label{fig:sup_scr}
\end{figure}

A Tabela~\ref{tab:sup_scr} sustenta quatro conclusões.
\textbf{(i) O pseudo-Brewster só se manifesta na configuração diametral, que é
inútil.} É a única coluna com pico em $\theta_{pB}$: em TM passa de $-55,3$ para
$-41,5\dB$, e em TE não há pico (controle). Ainda assim, fica mais de $40\dB$ abaixo
de qualquer outro canal. \textbf{(ii) Nos demais canais, maior incidência é melhor,
sem pico em $\theta_{pB}$.} O SCR cresce de forma quase monótona de $35^\circ$ a
$75^\circ$, igualmente em TM e TE (por exemplo, monoestático helicoidal:
$3,1\to8,0\to14,8\dB$ em TM e $6,2\to11,5\to18,5\dB$ em TE, em $53,4^\circ$,
$63,45^\circ$ e $75^\circ$), porque o retorno difuso decai com a incidência pelo
espectro de rugosidade $W$, e não por Brewster. \textbf{(iii) A TE supera a TM}
em $2$ a $5\dB$ no monoestático, ao contrário do que se esperaria se Brewster
importasse: em retroespalhamento, o difuso VV é maior que o HH. \textbf{(iv) Com o
alvo pontual}, o canal biestático TE a $90^\circ$ é o melhor em todas as linhas,
pelo nulo HH; em TM, um canal a $135^\circ$ é o melhor em sete de dez linhas. A
Seção~\ref{sec:sup_alvos_reais} mostra que apenas a segunda conclusão sobrevive
com alvos reais.

\paragraph{Hélice ou órbita de incidência constante.} A família helicoidal supera
a de incidência constante em 54 de 60 comparações de mesma incidência,
polarização e canal, com $3,5\dB$ de SCR em média ($6,2\dB$ no canal biestático a
$135^\circ$). O eco do alvo é idêntico nas duas famílias (diferença inferior a
$0,2\dB$); muda a quantidade de clutter que chega ao voxel do alvo. A superfície
está $2\,$m acima do alvo, e a órbita de incidência constante tem pior resolução
vertical ($2,34$ contra $2,08\,$m com janela de Hann) e pior ISLR, de modo que
vaza mais energia da superfície para o alvo. A órbita que impõe o
pseudo-Brewster é, portanto, pior também em SCR, a métrica para a qual foi
proposta.

\paragraph{Robustez ao erro angular.} Em S1, na órbita de incidência constante,
erros de até $\pm5^\circ$ alteram o SCR do monoestático, do biestático e da fusão em
cerca de $\pm3\dB$, sem tendência e da ordem da dispersão entre realizações; só a
configuração diametral cai ao afastar-se de $\theta_{pB}$ (até $15\dB$ com
$\pm5^\circ$). Fora da diametral, o SCR é robusto ao erro angular \emph{porque não
depende de Brewster}. No solo sem perdas, a diametral alcança $61\dB$ em
$\theta_{pB}$ exato e cai para $-12\dB$ com $\pm2^\circ$.

\paragraph{Outros fatores.} O SCR do monoestático e dos biestáticos
intermediários cai cerca de $6\dB$ cada vez que $s$ dobra (o difuso escala com
$s^2$); o da diametral quase não muda, pois é dominado pela especular. Nos solos
úmidos (S2 e S3), o SCR é negativo em todas as geometrias ($-14$ a $-28\dB$ em S2 e
$-58$ a $-84\dB$ em S3): domina a atenuação do alvo a $2\,$m, e o pseudo-Brewster
não a remedeia. Um rolamento de $0$ a $5^\circ$ não altera o SCR em S1
($\le0,2\dB$).

\section{Custo em Imagem da Restrição de Incidência}
\label{sec:sup_custo}

\begin{table}[H]
\centering\small
\caption{Hélice de referência ($53,4^\circ$) frente à órbita restrita ao
pseudo-Brewster, com janela de Hann. \colorbox{melhor}{Verde}: trajetória que
vence cada métrica no canal. Menor é melhor, exceto contraste e SCR.}
\label{tab:sup_custo}
\resizebox{\textwidth}{!}{\begin{tabular}{@{}llrrrrrr@{}}
\toprule
Trajetória & Canal & $\delta_x$ [m] & $\delta_z$ [m] & $V_{-3}$ [l] & ISLR [dB] & Contraste [dB] & SCR [dB] \\
\midrule
hélice ($53,4^\circ$) & mono & 0,157 & \cellcolor{melhor}2,08 & 29,0 & \cellcolor{melhor}6,5 & \cellcolor{melhor}19,6 & \cellcolor{melhor}3,1 \\
órbita em $\theta_{pB}$ & mono & \cellcolor{melhor}0,140 & 2,34 & \cellcolor{melhor}25,0 & 8,3 & 18,4 & 2,5 \\
\midrule
hélice ($53,4^\circ$) & mono+bi $135^\circ$ & 0,213 & \cellcolor{melhor}2,08 & 52,0 & \cellcolor{melhor}2,4 & \cellcolor{melhor}20,4 & \cellcolor{melhor}9,5 \\
órbita em $\theta_{pB}$ & mono+bi $135^\circ$ & \cellcolor{melhor}0,191 & 2,34 & \cellcolor{melhor}46,0 & 3,6 & 19,9 & 7,3 \\
\bottomrule
\end{tabular}}
\end{table}

Impor a órbita de pseudo-Brewster só traz ganho no plano horizontal: $\delta_x$
melhora $11\,\%$ (maior $\sin\psi$) e o volume do lóbulo, $12$ a $14\,\%$. Perde-se
em todo o resto: $\delta_z$ piora $12,5\,\%$ (perde-se a excursão de elevação da
hélice), o ISLR sobe $1,2$ a $1,8\dB$, o contraste cai $0,5$ a $1,2\dB$ e o SCR
cai $0,6$ a $2,2\dB$. A restrição obriga ainda a baixar a plataforma a cerca de
$85\,$m. O custo em imagem é pequeno, mas o benefício é nulo.

\section{Os Nulos com Alvos Reais}
\label{sec:sup_alvos_reais}

O SCR da Tabela~\ref{tab:sup_scr} usa um espalhador escalar, sem polarização:
só a superfície sente os nulos VV e HH. Um alvo real também espalha com
polarização, e o mesmo desacoplamento geométrico pode anular seu eco. Para
separar os dois efeitos, repetiu-se a análise na hélice de referência (S1,
$s=2\,$cm, três realizações) com o receptor na mesma hélice e
$\Delta\phi$ de $0^\circ$ a $180^\circ$ em passos de $15^\circ$, mais $140^\circ$, em TM e TE
co-polares, para quatro alvos: o ponto escalar, uma esfera condutora de raio
$0,5\,$m, um cubo de $1\,$m em orientação nominal e o mesmo cubo em 24 orientações
aleatórias em $SO(3)$. Além do SCR no voxel, mediu-se a relação entre o pico do
alvo e o pico da superfície em toda a grade (PSCR).

\begin{figure}[H]
\centering
\includegraphics[width=\textwidth]{C2_nulos_scr.pdf}
\caption{Nulos com alvos reais em função de $\Delta\phi$ (hélice de referência,
S1, $s=2\,$cm). (a, d) Potência na imagem do canal biestático relativa ao
monoestático: superfície no voxel do alvo (vermelho) e eco de cada alvo. (b, e)
Ganho de SCR no voxel. (c, f) Ganho de PSCR. Contínuo: biestático isolado;
tracejado: fusão a orçamento constante. Linhas cinzas: nulo HH ($90^\circ$) e
nulo VV ($140^\circ$).}
\label{fig:sup_nulos_reais}
\end{figure}

\begin{table}[H]
\centering\small
\caption{SCR no voxel do alvo [dB] com alvos reais. ``Superfície'': potência da
superfície no voxel relativa ao monoestático. \colorbox{melhor}{Verde}: melhoria
superior a $3\dB$ sobre o monoestático da mesma polarização;
\colorbox{pior}{vermelho}: piora superior a $3\dB$.}
\label{tab:sup_nulos_reais}
\begin{tabular}{@{}llrrrrr@{}}
\toprule
Pol. & Canal & Superfície & Ponto & Esfera & Cubo nominal & Cubo $SO(3)$ \\
\midrule
TM & mono & $0,0$ & $0,7$ & $-0,9$ & $-9,5$ & $-12,0$ \\
 & bi $90^\circ$ & $-2,0$ & $2,6$ & \cellcolor{pior}$-13,3$ & \cellcolor{pior}$-25,3$ & \cellcolor{pior}$-27,9$ \\
 & bi $120^\circ$ & $-4,6$ & \cellcolor{melhor}$5,2$ & \cellcolor{pior}$-6,1$ & \cellcolor{pior}$-19,1$ & \cellcolor{melhor}$-6,2$ \\
 & bi $135^\circ$ & $-12,5$ & \cellcolor{melhor}$13,1$ & \cellcolor{melhor}$7,5$ & \cellcolor{melhor}$14,7$ & \cellcolor{melhor}$6,2$ \\
 & bi $140^\circ$ & $-13,6$ & \cellcolor{melhor}$14,2$ & \cellcolor{melhor}$\mathbf{9,8}$ & \cellcolor{melhor}$\mathbf{20,1}$ & \cellcolor{melhor}$\mathbf{8,3}$ \\
 & bi $150^\circ$ & $-1,2$ & $1,8$ & $-0,9$ & \cellcolor{melhor}$13,7$ & \cellcolor{melhor}$-2,7$ \\
 & mono+bi $140^\circ$ & $-5,7$ & \cellcolor{melhor}$6,4$ & $-0,9$ & \cellcolor{melhor}$4,7$ & \cellcolor{melhor}$-4,2$ \\
\midrule
TE & mono & $0,0$ & $3,9$ & $2,3$ & $-6,4$ & $-8,8$ \\
 & bi $90^\circ$ & $-17,8$ & \cellcolor{melhor}$\mathbf{21,7}$ & \cellcolor{pior}$-4,8$ & $-6,2$ & $-9,0$ \\
 & mono+bi $90^\circ$ & $-5,5$ & \cellcolor{melhor}$9,4$ & $2,2$ & $-7,4$ & $-9,3$ \\
\bottomrule
\end{tabular}
\end{table}

\paragraph{O nulo HH a $90^\circ$ é um artefato do espalhador escalar.} A
superfície cai $17,8\dB$, mas o eco da esfera cai $25\dB$ e o do cubo em
orientação aleatória, $18\dB$. O SCR da esfera piora $7,1\dB$ em relação ao
monoestático, e o do cubo não muda. É o mesmo mecanismo da
Proposição~\ref{prop:descasamento}: a $90^\circ$, os vetores TE do transmissor e
do receptor são ortogonais, o que desacopla tanto a superfície quanto o alvo.

\paragraph{O nulo VV a $135^\circ$--$140^\circ$ sobrevive.} A superfície cai $12,5$ a
$13,6\dB$, enquanto o eco dos alvos cai pouco (esfera, $-4,1$ e $-2,9\dB$) ou até
sobe (cubo em orientação aleatória, $+5,8$ e $+6,7\dB$, pela especular de suas
faces ao longo da bissetriz). O SCR do canal biestático isolado melhora, em
relação ao monoestático, de $8,4$ a $10,7\dB$ na esfera, de $24$ a $30\dB$ no cubo
nominal e de $18$ a $20\dB$ no cubo em orientação aleatória. O PSCR melhora menos,
mas com o mesmo sinal: $+2,5$ a $+3,7\dB$ na esfera e $+12$ a $+14\dB$ no cubo.

\paragraph{A janela é estreita.} A $120^\circ$ e a $150^\circ$, a esfera já não
melhora ($-6,1$ e $-0,9\dB$), e a partir de $150^\circ$ a especular do nadir domina o
PSCR. O receptor deve manter-se a cerca de $\pm5^\circ$ de $140^\circ$, coerentemente
com a largura de $4$ a $8^\circ$ do nulo VV.

\paragraph{A fusão de pesos iguais desperdiça o canal limpo.} A imagem fundida é
a média das duas, e o clutter do monoestático a domina: sua superfície cai
apenas $5,7\dB$, próximo do limite de $6\dB$ da média de dois canais. A $140^\circ$,
a fusão deixa a esfera igual ao monoestático e fica $10$ a $15\dB$ abaixo do
biestático isolado. Para explorar o nulo, deve-se usar o canal biestático isolado
ou uma fusão ponderada pelo clutter de cada canal.

\paragraph{O preço é a resolução.} A $135^\circ$, o canal biestático tem
$\delta_x=0,41\,$m, contra $0,157\,$m no monoestático, com a mesma $\delta_z$. Trata-se
de um modo de \emph{detecção} em cenários limitados por clutter, não de um modo de
imageamento fino.

\section{Síntese do Capítulo}
\label{sec:sup_sintese}

\begin{table}[H]
\centering\small
\caption{Respostas diretas sobre a supressão do retorno da superfície (S1,
$s=2\,$cm). \colorbox{melhor}{Verde}: opção recomendada.}
\label{tab:sup_respostas}
\begin{tabular}{@{}L{3.5cm}L{2.9cm}L{3.5cm}L{3.9cm}@{}}
\toprule
\textbf{Pergunta} & \textbf{Melhor opção} & \textbf{Margem} & \textbf{Razão} \\
\midrule
Restringir a trajetória ao pseudo-Brewster? & \cellcolor{melhor}Não & a órbita restrita perde $0,6$ a $2,2\dB$ de SCR & só suprime a especular, que só pesa na diametral \\
Hélice ou incidência constante? & \cellcolor{melhor}Hélice & $+3,5\dB$ de SCR médio; vence 54/60 & melhor $\delta_z$ e ISLR \\
Qual incidência? & \cellcolor{melhor}A maior viável & mono TM: $+11,7\dB$ de $53,4^\circ$ a $75^\circ$ & o difuso cai com $\theta$, igualmente em TM e TE \\
Qual canal em TM? & \cellcolor{melhor}Bi isolado a $135^\circ$--$140^\circ$ & $8$ a $20\dB$ sobre o mono com alvos reais & nulo VV biestático \\
TM ou TE? & \cellcolor{melhor}TM & TE a $90^\circ$ só vence com o ponto escalar & o desacoplamento anula também o alvo \\
Fusão de pesos iguais? & \cellcolor{pior}Não para clutter & $10$ a $15\dB$ abaixo do bi isolado & o clutter do mono domina a média \\
Configuração diametral? & \cellcolor{pior}Nunca & $-41$ a $-81\dB$ & a especular cai sobre o alvo \\
\bottomrule
\end{tabular}
\end{table}

\begin{sintese}
O pseudo-Brewster não é um critério útil de trajetória: num solo com perdas é um
mínimo raso, que só atua sobre a reflexão especular, e esta só alcança o
receptor na configuração diametral, inútil para imageamento; a órbita que o
impõe é pior que a hélice em SCR, resolução vertical e lóbulos. O que eleva o SCR
é a geometria: maior incidência e, sobretudo, o receptor a $135^\circ$--$140^\circ$
em TM usado como canal isolado, no nulo biestático VV de primeira ordem, que eleva
de $8$ a $20\dB$ o SCR de alvos reais. O nulo HH a $90^\circ$ é um artefato do alvo
escalar. Como o nulo VV é de primeira ordem, ele deve ser confirmado com modelos de
espalhamento de ordem superior e com medidas.
\end{sintese}

%%═══════════════════════════════════════════════════════════════════════════════
\chapter{Projeto da Trajetória do Receptor e Requisitos de Sincronização}
\label{cap:trajetoria}
%%═══════════════════════════════════════════════════════════════════════════════

\section{Introdução}
\label{sec:traj_intro}

Os capítulos anteriores identificaram a região promissora
($\Delta\phi\approx90^\circ$--$120^\circ$ para imageamento,
$\Delta\phi\approx140^\circ$ para supressão de clutter) por varredura direta, sem
otimizador, como recomendam a boa prática e a própria natureza exploratória do
problema. Este capítulo converte essa compreensão física num projeto de
trajetória e quantifica as condições operacionais que o tornam viável:
tolerâncias de posição, de relógio e de fase entre os canais.

\section{Espaço de Projeto e Função de Custo}
\label{sec:traj_custo}

A trajetória do transmissor é conhecida e fixa, e o receptor voa à mesma altura.
O espaço de projeto é parametrizado por uma separação de partida
$\Delta\phi_0\in\{60^\circ,\dots,165^\circ\}$, uma varredura linear
\begin{equation}
  \Delta\phi(t)=\Delta\phi_0+A\Big(\frac{2t}{T}-1\Big),\qquad A\in\{0,\,15^\circ,\,30^\circ,\,45^\circ\},
\end{equation}
e uma razão de raios $\kappa\in\{0,85;\,1;\,1,15\}$: 96 candidatos. Nenhuma
métrica contém a potência do eco do alvo. Avaliam-se separadamente
\begin{itemize}
  \item a qualidade de imagem, $Q_{\text{img}}=\text{ISLR}_{\text{mono}}-\text{ISLR}_{\text{fus}}$;
  \item o crescimento da PSF, $A_{\text{PSF}}=10\log_{10}(V_{\text{fus}}/V_{\text{mono}})$;
  \item o ganho de clutter, $\Delta\text{SCR}=\text{SCR}_{\text{fus}}-\text{SCR}_{\text{mono}}$
  (S1, $s=2\,$cm, TM, duas realizações), usado exclusivamente como medida de
  supressão do retorno da superfície;
  \item o custo de voo, $C_{\text{voo}}=(v^{Rx}_{\max}/v_{Tx}-1)_++|L_{Rx}/L_{Tx}-1|$,
  com o candidato declarado inviável se a separação entre os VANTs for inferior a
  $15\,$m ou se $v^{Rx}_{\max}>1,5\,v_{Tx}$.
\end{itemize}
Com as métricas normalizadas ao intervalo $[0,1]$ dos candidatos viáveis, a
função de custo ponderada é
\begin{equation}
  J = w_1\,\hat Q_{\text{img}} + w_2\,\widehat{\Delta\text{SCR}} - w_3\,\hat A_{\text{PSF}} - w_4\,\hat C_{\text{voo}},
  \qquad (w_1,w_2,w_3,w_4)=(1;\,0,5;\,1;\,1),
  \label{eq:custo_imagem}
\end{equation}
a ser maximizada. Os pesos são explicitados e acompanhados da fronteira de
Pareto, e as métricas são reportadas separadamente, para que um único valor de
custo não oculte os compromissos. A janela espectral é a de Hann.

\begin{figure}[H]
\centering
\includegraphics[width=\textwidth]{F_pareto.pdf}
\caption{Fronteira de Pareto (círculos vermelhos, 38 de 96 candidatos) nas
projeções $Q_{\text{img}}$--$C_{\text{voo}}$, $Q_{\text{img}}$--$\Delta$SCR e
$Q_{\text{img}}$--$A_{\text{PSF}}$; a cor indica $\Delta\phi_0$, e a estrela, o
máximo de $J$.}
\label{fig:traj_pareto}
\end{figure}

\begin{table}[H]
\centering\small
\caption{Escolha segundo os pesos (todas Pareto-ótimas).}
\label{tab:traj_pesos}
\resizebox{\textwidth}{!}{\begin{tabular}{@{}lcccc|ccc|rrrr@{}}
\toprule
Pesos & $w_1$ & $w_2$ & $w_3$ & $w_4$ & $\Delta\phi_0$ & $A$ & $\kappa$ & $Q_{\text{img}}$ & $\Delta$SCR & $A_{\text{PSF}}$ & $C_{\text{voo}}$ \\
\midrule
principal & 1 & 0,5 & 1 & 1 & $105^\circ$ & $45^\circ$ & 0,85 & 7,72 & 4,65 & 1,83 & 0,04 \\
sem SCR nem voo & 1 & 0 & 1 & 0 & $105^\circ$ & $45^\circ$ & 1,15 & 7,36 & 0,66 & 1,55 & 0,60 \\
mais SCR & 1 & 1 & 0,5 & 1 & $105^\circ$ & $45^\circ$ & 0,85 & 7,72 & 4,65 & 1,83 & 0,04 \\
voo prioritário & 0,5 & 0,5 & 0,5 & 2 & $75^\circ$ & $0^\circ$ & 1,00 & 4,00 & 3,12 & 0,96 & 0,00 \\
\bottomrule
\end{tabular}}
\end{table}

\section{Trajetória Proposta}
\label{sec:traj_proposta}

Com três dos quatro conjuntos de pesos, a escolha é a mesma: receptor à mesma
altura do transmissor, raio $0,85\,\rho_{Tx}$ e separação que cresce linearmente de
$60^\circ$ a $150^\circ$ ao longo das duas voltas. Somente quando o custo de voo
domina escolhe-se uma separação constante de $75^\circ$. A escolha é estável.

\begin{table}[H]
\centering\small
\caption{Trajetória proposta frente ao monoestático, à melhor separação constante e
à região da configuração diametral (janela de Hann).}
\label{tab:traj_final}
\resizebox{\textwidth}{!}{\begin{tabular}{@{}lrrrrrrr@{}}
\toprule
 & PSLR & ISLR & $V_{-3}$ & $\delta_x$ & $\delta_z$ & Contr. & SCR \\
 & [dB] & [dB] & [l] & [m] & [m] & [dB] & [dB] \\
\midrule
Monoestático & $-7,73$ & 6,71 & 29,4 & 0,158 & 2,08 & 19,5 & 6,9 \\
Mono+bi, $\Delta\phi=105^{\circ}$ constante & $-10,40$ & 1,86 & 43,2 & 0,19 & 2,08 & 22,4 & 10,0 \\
\textbf{Mono+bi proposta} ($60^{\circ}\to150^{\circ}$, $\kappa=0,85$) & $\mathbf{-10,10}$ & $\mathbf{-1,01}$ & 44,8 & 0,20 & 2,06 & \textbf{25,0} & 11,5 \\
Mono+bi, $\Delta\phi=165^{\circ}$ constante & $-5,2$ & 7,5 & 61,7 & 0,23 & 2,08 & 14,5 & $-32,8$ \\
\bottomrule
\end{tabular}}
\end{table}

\begin{figure}[H]
\centering
\includegraphics[width=\textwidth]{F_trayectoria_final.pdf}
\caption{Trajetória proposta do receptor comparada com a configuração diametral e
com a monoestática; à direita, $\Delta\phi(t)$ com o transmissor a $10\,$m/s
($190\,$s de aquisição).}
\label{fig:traj_final}
\end{figure}

Frente à imagem monoestática, a trajetória proposta reduz o ISLR em $7,7\dB$ e o
PSLR em $2,4\dB$ e eleva o contraste em $5,5\dB$, à custa de um lóbulo principal
$52\,\%$ mais volumoso ($\delta_x$ de $0,158$ para $0,20\,$m). Frente à melhor
separação constante ($105^\circ$), a varredura ganha $2,9\dB$ de ISLR: anéis de raios
diferentes preenchem o disco de números de onda (Seção~\ref{sec:geo_interp}). O
custo de voo é baixo: velocidade máxima do receptor de $0,97$ vez a do
transmissor, percurso relativo de $0,96$, aceleração máxima de $0,77\,$m/s$^2$ a
$10\,$m/s e separação mínima de $128\,$m entre os VANTs.

Dois resultados negativos acompanham o projeto. \textbf{O raio não acrescenta
resolução vertical}: à mesma altura, $\kappa$ entre $0,85$ e $1,15$ altera $\delta_z$
em cerca de $1\,\%$, porque a extensão em $k_z$ é fixada pela banda, e não pela
pequena diferença de elevação ($\pm3^\circ$), em acordo com a
Proposição~\ref{prop:limite_z}. \textbf{O SCR não ordena as separações
intermediárias}: o ganho de SCR da fusão fica entre $-0,5$ e $+5,6\dB$ para
$\Delta\phi_0\le150^\circ$, da ordem da dispersão entre realizações ($\pm3\dB$); só a
região $\Delta\phi_0\ge165^\circ$ é claramente ruim ($-33$ a $-40\dB$).

\section{Robustez, Sincronização e Calibração de Fase}
\label{sec:traj_sincro}

A fusão coerente pressupõe que os dois canais compartilhem uma referência de
fase. Num sistema com plataformas separadas, a sincronização de relógio e de fase
é o principal desafio tecnológico do SAR biestático~\citar{krieger2006}, \citar{younis2006};
sistemas UAV recentes resolvem-no com cadeias de sincronização bidirecionais ou
com receptores digitais coerentes~\citar{liu2024}, \citar{kanz2025}. Avaliou-se a
sensibilidade da trajetória proposta a cada fonte de erro do canal biestático,
sem calibração.

\begin{figure}[H]
\centering
\includegraphics[width=\textwidth]{F_robustez.pdf}
\caption{Robustez da fusão com a trajetória proposta frente a erros do canal
biestático não calibrado.}
\label{fig:traj_robustez}
\end{figure}

\begin{table}[H]
\centering\small
\caption{Tolerâncias do canal biestático (fusão proposta; nominal: PSLR
$-10,1\dB$, ISLR $-1,0\dB$; monoestático: $-7,7$ e $6,7\dB$).}
\label{tab:traj_tolerancias}
\begin{tabular}{@{}L{4.8cm}L{3.6cm}L{6.0cm}@{}}
\toprule
Perturbação & Tolerância & Observação \\
\midrule
Ruído de posição do receptor (por pulso) & $\le10\,$cm ($\lambda_0/7$) & ISLR $-0,2\dB$ a $10\,$cm; média na abertura \\
Deriva de posição (passeio aleatório, rms) & $\le5\,$cm & $20\,$cm destrói o ganho (PSLR $-2,0\dB$) \\
Ruído de fase do oscilador (por pulso) & $\le45^\circ$ rms & apenas perda de amplitude \\
Desvio de sincronismo constante & \textbf{calibrável} & $1\,$ns $=153^\circ$ de fase: sem calibração anula a fusão; calibrado, PSLR $-9,7$ e ISLR $-1,0\dB$ \\
Deriva de sincronismo (linear) & $<1\,$ns na aquisição & $1\,$ns: PSLR $-6,5\dB$; $5\,$ns: $-2,1\dB$ \\
Erro do índice do solo ($\pm10\,\%$) & sim & afeta igualmente mono e fusão (localização $0,17\,$m) \\
\bottomrule
\end{tabular}
\end{table}

\begin{figure}[H]
\centering
\includegraphics[width=\textwidth]{F2_fase.pdf}
\caption{Fase residual constante entre canais: sem calibração (laranja), calibrada
sobre o lóbulo principal do espalhador mais intenso (verde), fusão incoerente
$\tfrac12(|I_m|^2+|I_b|^2)$ (tracejado) e monoestático (azul).}
\label{fig:traj_fase}
\end{figure}

\paragraph{Fase constante.} A fusão coerente exige conhecer uma fase constante
$\varphi_0$ entre canais com erro inferior a cerca de $30^\circ$: a $60^\circ$ já se
perdem $2,1\dB$ de ISLR, e a $150^\circ$ a fusão é pior que o monoestático. Um
desvio de relógio constante $\Delta t$ produz exatamente uma fase
$\varphi_0=2\pi f_c\Delta t$, mais um deslocamento em alcance desprezível,
$c\,\Delta t/(2n)$. Ela é estimada sem ambiguidade sobre o lóbulo principal de um
refletor intenso, ou a partir do sinal direto transmissor$\to$receptor, cuja
geometria é conhecida por GNSS-RTK; calibrada, a fusão recupera exatamente o
desempenho nominal. Estimar a fase sobre a imagem inteira falha (erro de
cerca de $180^\circ$), porque os anéis de lóbulos em contrafase dominam a
estimativa.

\paragraph{Derivas.} As derivas de relógio e de posição não são calibráveis por
uma constante. A deriva de relógio deve permanecer muito abaixo de $1\,$ns nos
$190\,$s de aquisição, isto é, uma estabilidade relativa da ordem de
$5\times10^{-12}$, compatível com osciladores disciplinados por GPS ou com um
enlace de sincronização bidirecional; a deriva de posição deve ficar abaixo de
$5\,$cm rms, compatível com RTK. Os ruídos brancos de posição e de fase são bem
tolerados, porque se promediam ao longo da abertura.

\paragraph{Fusão incoerente.} Somar intensidades dispensa a fase, mas rende
$2,0\dB$ menos de ISLR e $4,4\dB$ menos de contraste que a fusão coerente calibrada,
com PSLR semelhante ($-10,7\dB$). É uma alternativa de reserva quando não houver
calibração confiável, não um substituto.

\paragraph{Modelo de erro realista.} A fase constante é o caso ideal. Um modelo
de erro mais completo entre os canais é
\begin{equation}
  \varphi(f,t)=\varphi_0+2\pi f\,\Delta\tau(t)+H(f),
  \label{eq:erro_fase}
\end{equation}
em que $\Delta\tau(t)$ reúne a deriva de relógio e o erro de posição do receptor e
$H(f)$ é a diferença entre as respostas em frequência das duas cadeias de
recepção. Os resultados acima indicam a escala dos termos variáveis; sua
estimação pulso a pulso pelo sinal direto e a calibração de $H(f)$ com um refletor
de referência fazem parte do protocolo experimental do
Capítulo~\ref{cap:dados_reais}.

\section{Dois Modos de Operação do Receptor}
\label{sec:traj_modos}

Os resultados dos Capítulos~\ref{cap:geometria} a~\ref{cap:trajetoria} conduzem
a dois modos de operação distintos, cujos objetivos são incompatíveis numa única
geometria:

\begin{description}
  \item[Modo de imageamento.] Receptor à mesma altura, raio $0,85\,\rho_{Tx}$,
  separação varrida de $60^\circ$ a $150^\circ$, dupla polarização alinhada ao campo
  do transmissor, fusão coerente com calibração de fase. Melhora lóbulos e
  contraste à custa de um lóbulo principal mais largo.
  \item[Modo de detecção limitada por clutter.] Receptor a
  $\Delta\phi\approx140^\circ\pm5^\circ$ em polarização TM, canal biestático usado
  isoladamente ou numa fusão ponderada pelo clutter. Eleva o SCR de alvos reais em
  $8$ a $20\dB$, com uma resolução horizontal $2,6$ vezes pior.
\end{description}

A trajetória do modo de imageamento foi selecionada com o alvo pontual; sua
reavaliação com alvos anisotrópicos em orientação aleatória, com o receptor
alinhado e sob clutter, usando a mediana do contraste como métrica de imagem, é
um passo necessário identificado no Capítulo~\ref{cap:conclusoes}.

\begin{sintese}
A trajetória recomendada para imageamento --- mesma altura, raio $0,85\,\rho_{Tx}$ e
separação varrida de $60^\circ$ a $150^\circ$ --- reduz o ISLR em $7,7\dB$ e eleva o
contraste em $5,5\dB$ frente ao monoestático, com custo de voo quase nulo. A
condição é a coerência: uma fase constante entre canais deve ser calibrada com erro
inferior a $30^\circ$, e a deriva de relógio deve ficar muito abaixo de $1\,$ns na
abertura.
\end{sintese}

%%═══════════════════════════════════════════════════════════════════════════════
\chapter{Avaliação com Dados Reais}
\label{cap:dados_reais}
%%═══════════════════════════════════════════════════════════════════════════════

% NOTA: Este capítulo está estruturado para receber os resultados da campanha
%       experimental. As seções de protocolo foram redigidas a partir das
%       conclusões dos Capítulos 7 a 10; as seções de resultados e discussão
%       devem ser preenchidas após o processamento dos dados medidos.

\section{Objetivos}
\label{sec:real_objetivos}

As conclusões desta tese foram obtidas por simulação eletromagnética coerente,
com modelos cujos limites estão listados na Seção~\ref{sec:concl_lim}: Óptica
Física para os alvos, espalhamento de superfície de primeira ordem, solo homogêneo
e sincronização idealizada ou perturbada de forma controlada. Este capítulo
reserva o espaço para a avaliação dessas conclusões com dados reais e define o
protocolo correspondente. A avaliação tem dois componentes complementares: o
processamento de dados medidos em voo e a \emph{simulação informada por dados
reais}, em que o simulador do Capítulo~\ref{cap:modelo} é alimentado com as
trajetórias, os parâmetros do solo e as respostas das cadeias de recepção
efetivamente medidos, de modo que a diferença entre a imagem medida e a
simulada isole os efeitos não modelados.

As hipóteses a testar, derivadas dos resultados de simulação, são:

\begin{enumerate}[label=\textbf{E\arabic*.},leftmargin=*]
  \item \textbf{Invariância da potência.} A potência média recebida de alvos de
  orientação não controlada não varia sistematicamente com a separação
  azimutal (Capítulo~\ref{cap:potencia}).
  \item \textbf{Ganho da fusão.} A fusão coerente calibrada com separação na faixa
  de $60^\circ$ a $150^\circ$ reduz o ISLR e eleva o contraste da imagem monoestática,
  sem melhorar a resolução a $-3\dB$ (Capítulos~\ref{cap:geometria}
  e~\ref{cap:trajetoria}).
  \item \textbf{Regra de polarização.} Um receptor de dupla polarização alinhado
  recupera o descasamento co-polar a $\Delta\phi\approx90^\circ$
  (Capítulo~\ref{cap:alvos}).
  \item \textbf{Nulo biestático VV.} O retorno difuso da superfície em TM
  apresenta um mínimo próximo de $\Delta\phi^\ast$ dado por~\eqref{eq:nulo_vv},
  e o canal biestático nessa separação eleva a relação alvo--clutter
  (Capítulo~\ref{cap:superficie}).
  \item \textbf{Calibração de fase.} Uma fase constante entre canais, estimada
  pelo sinal direto ou por um refletor de referência, é suficiente para
  recuperar o ganho da fusão; as derivas residuais permanecem dentro das
  tolerâncias da Tabela~\ref{tab:traj_tolerancias}.
\end{enumerate}

\section{Sistema de Aquisição}
\label{sec:real_sistema}

% TODO: Descrever o hardware efetivamente utilizado: radar do transceptor,
%       receptor remoto, antenas (padrão, polarização, dupla polarização),
%       referência de tempo e de frequência, registro de dados, VANTs e
%       sistema de posicionamento.

\begin{table}[H]
\centering\small
\caption{Parâmetros do sistema experimental frente aos valores usados nas
simulações.}
\label{tab:real_parametros}
\begin{tabular}{@{}lll@{}}
\toprule
Parâmetro & Simulação & Experimento \\
\midrule
Frequência portadora & $425\,$MHz & [a definir] \\
Largura de banda & $50\,$MHz & [a definir] \\
Potência transmitida & $100\,$kW (pico) & [a definir] \\
Ganho das antenas & $6\,$dBi & [a definir] \\
Polarização do transmissor & TM & [a definir] \\
Polarização do receptor remoto & TM co-polar / dupla alinhada & [a definir] \\
Referência de tempo & ideal ou perturbada & [a definir] \\
Posicionamento & ideal ou perturbado & GNSS-RTK [a confirmar] \\
Trajetória do transceptor & hélice de duas voltas & [a definir] \\
Velocidade do transceptor & $10\,$m/s & [a definir] \\
\bottomrule
\end{tabular}
\end{table}

\section{Cenário de Ensaio}
\label{sec:real_cenario}

% TODO: Descrever o sítio de ensaio e os alvos.

O cenário de ensaio deve conter:
\begin{itemize}
  \item \textbf{Alvos enterrados} de geometria conhecida, a profundidades
  conhecidas, que reproduzam as classes da biblioteca simulada: ao menos um alvo
  de resposta quase isotrópica (esfera ou cilindro vertical curto), um tubo
  horizontal e um objeto com face plana horizontal, com espaçamento suficiente
  para avaliar o mascaramento entre objetos de RCS distintas.
  \item \textbf{Refletores de referência} na superfície (refletores de canto
  triédricos), fora da região de interesse, para a calibração de fase entre
  canais e da resposta em frequência $H(f)$ de~\eqref{eq:erro_fase}.
  \item \textbf{Caracterização do solo}: permissividade complexa em função da
  profundidade e da umidade, medida no local na data da campanha; perfil de
  rugosidade (rugosidade rms e comprimento de correlação); levantamento
  topográfico da superfície.
  \item \textbf{Região de clutter controlado}: uma área sem alvos, com a mesma
  superfície, para medir o retorno da superfície em função de $\Delta\phi$ e
  testar a hipótese E4.
\end{itemize}

\section{Plano de Voo}
\label{sec:real_voo}

% TODO: Detalhar as trajetórias efetivamente voadas, o número de repetições e
%       as condições meteorológicas.

Para cada cenário, propõem-se as seguintes aquisições, todas com o mesmo voo do
transceptor:
\begin{enumerate}
  \item canal monoestático isolado (referência permanente);
  \item receptor remoto em separações constantes de $45^\circ$, $90^\circ$,
  $120^\circ$, $140^\circ$ e $180^\circ$, para reproduzir a varredura do
  Capítulo~\ref{cap:geometria} e testar E1 e E4;
  \item receptor remoto na trajetória proposta do modo de imageamento
  (Seção~\ref{sec:traj_proposta}), para testar E2;
  \item a separação de $90^\circ$ com receptor co-polar e com receptor alinhado,
  para testar E3.
\end{enumerate}
Cada aquisição deve ser repetida ao menos três vezes, para estimar a dispersão
entre voos e permitir testes estatísticos pareados entre configurações.

\section{Sincronização, Calibração e Compensação de Movimento}
\label{sec:real_calibracao}

% TODO: Descrever o procedimento efetivamente implementado.

O procedimento de calibração decorre dos requisitos da
Seção~\ref{sec:traj_sincro}:
\begin{enumerate}
  \item estimação, pulso a pulso, do retardo $\Delta\tau(t)$ a partir do sinal
  direto transmissor$\to$receptor, cuja geometria é conhecida pelas trajetórias
  medidas;
  \item calibração da resposta em frequência relativa $H(f)$ das cadeias de
  recepção sobre um refletor de referência;
  \item estimação da fase residual constante $\varphi_0$ sobre o lóbulo principal
  do refletor de referência, e não sobre a imagem inteira;
  \item retroprojeção com refração usando as trajetórias medidas de ambas as
  plataformas e o índice do solo estimado no local, com verificação do foco sobre
  os alvos de posição conhecida.
\end{enumerate}

\section{Simulação Informada por Dados Reais}
\label{sec:real_simulacao}

% TODO: Executar o simulador do Capítulo 6 com os parâmetros medidos e comparar.

Para cada aquisição, o simulador deve ser executado com as trajetórias medidas,
a permissividade e a rugosidade medidas, os padrões de antena medidos e os erros
de sincronização residuais estimados. A comparação entre a imagem medida e a
simulada, pelas mesmas métricas da Tabela~\ref{tab:metricas}, permite distinguir
três fontes de discrepância: limitações do modelo de alvo (Óptica Física),
limitações do modelo de superfície (primeira ordem) e erros não calibrados do
sistema.

\section{Métricas e Critério de Sucesso}
\label{sec:real_criterio}

As métricas são as da Tabela~\ref{tab:metricas}, calculadas sobre as respostas
dos alvos de posição conhecida, com a mesma referência e a mesma faixa dinâmica
para todas as imagens. A contribuição desta tese estará demonstrada
experimentalmente se a imagem fundida superar objetivamente a imagem monoestática
em ISLR e contraste, de forma estatisticamente significativa entre repetições, com
a fusão normalizada a orçamento constante, e se a melhoria desaparecer quando a
fase entre canais não for calibrada.

\section{Resultados}
\label{sec:real_resultados}

% TODO: Inserir os resultados.

\textit{[Resultados a serem inseridos após a campanha experimental.]}

\section{Discussão}
\label{sec:real_discussao}

% TODO: Confrontar os resultados com as hipóteses E1--E5 e com a simulação
%       informada por dados reais.

\textit{[Discussão a ser inserida após a campanha experimental.]}

%%═══════════════════════════════════════════════════════════════════════════════
\chapter{Conclusões e Perspectivas}
\label{cap:conclusoes}
%%═══════════════════════════════════════════════════════════════════════════════

\section{Respostas às Perguntas de Pesquisa}
\label{sec:concl_respostas}

\begin{longtable}{@{}L{4.4cm}L{10.4cm}@{}}
\caption{Perguntas de pesquisa e respostas, com a evidência que as sustenta.}
\label{tab:concl_respostas}\\
\toprule
\textbf{Pergunta} & \textbf{Resposta} \\
\midrule
\endfirsthead
\toprule
\textbf{Pergunta} & \textbf{Resposta} \\
\midrule
\endhead
A potência recebida, ou o ângulo de Brewster, determinam a melhor posição do
receptor? &
\textbf{Não.} A transmitância de Brewster acrescenta cerca de $1\dB$ ao enlace e
é dominada pela perda geométrica; a potência é invariante à separação azimutal
para alvos centrados e, promediada sobre a orientação, é a mesma em qualquer
geometria ($-0,37$ a $+0,21\dB$). Capítulo~\ref{cap:potencia}. \\
Como a separação Tx--Rx altera o suporte em $k$ e a PSF? &
À mesma altura e no mesmo raio, o canal biestático ocupa um anel horizontal de raio
$|\cos(\Delta\phi/2)|$ vezes o monoestático, com a mesma extensão em $k_z$; a PSF
alarga-se como $1/|\cos(\Delta\phi/2)|$ e a resolução vertical não muda. Na
configuração diametral o anel colapsa e não há foco horizontal.
Capítulos~\ref{cap:resolucao} e~\ref{cap:geometria}. \\
Qual separação produz a melhor imagem biestática isolada? &
A menor viável; não há ótimo intermediário, e o canal mais nítido é praticamente
um segundo monoestático. Em SCR, porém, o canal isolado em TM a
$135^\circ$--$140^\circ$ é o melhor com alvos reais. Capítulos~\ref{cap:geometria}
e~\ref{cap:superficie}. \\
Qual separação acrescenta mais informação à imagem monoestática? &
$\Delta\phi\approx90^\circ$--$120^\circ$ constante (posto efetivo de $28,4$ para
$33,8$; área ocupada em $(k_x,k_y)$ de $294$ para $475\,$rad$^2$/m$^2$) e, melhor
ainda, uma separação varrida de $60^\circ$ a $150^\circ$. Capítulos~\ref{cap:geometria}
e~\ref{cap:trajetoria}. \\
A fusão coerente melhora resolução, volume, PSLR, ISLR, contraste ou
localização? &
\textbf{Melhora lóbulos e contraste, não a resolução.} Com o alvo pontual: ISLR até
$-7,7\dB$, PSLR até $-5,2\dB$, contraste até $+5,5\dB$; a resolução horizontal
piora de $16$ a $30\,\%$, o volume de $35$ a $58\,\%$, e a separabilidade de alvos
próximos cai. A localização não muda. Todo o ganho é geométrico. \\
A melhoria persiste com RCS anisotrópica e orientação aleatória? &
\textbf{Sim, com receptor de dupla polarização alinhado.} Contraste $+1,8$ a
$+2,7\dB$ em $90$ a $98\,\%$ das orientações a $90^\circ$. Com receptor TM
co-polar, o canal biestático perde cerca de $20\dB$ e o ganho é marginal.
Capítulo~\ref{cap:alvos}. \\
O pseudo-Brewster reduz o retorno da superfície de forma robusta? &
\textbf{Apenas parcialmente, e só onde não interessa.} Num solo com perdas o
mínimo é largo, mas raso ($-34,5\dB$), e só atua sobre a especular, que só chega
ao receptor na configuração diametral; nas demais geometrias o clutter é difuso e
o SCR não depende de Brewster. O mecanismo útil é o nulo biestático VV
($\Delta\phi^\ast\approx140^\circ$), que eleva de $8$ a $20\dB$ o SCR de alvos reais.
Capítulo~\ref{cap:superficie}. \\
Qual é o custo de impor o pseudo-Brewster à trajetória? &
\textbf{Pequeno, mas sem benefício}: $\delta_x$ $-11\,\%$, $\delta_z$ $+12,5\,\%$,
ISLR $+1,2$ a $+1,8\dB$, contraste $-0,5$ a $-1,2\dB$ e SCR $-0,6$ a $-2,2\dB$; a
órbita de incidência constante é pior que a hélice em 54 de 60 comparações de SCR.
Capítulo~\ref{cap:superficie}. \\
\bottomrule
\end{longtable}

\section{Contribuições}
\label{sec:concl_contrib}

\begin{enumerate}
  \item \textbf{Resolução 3D com refração em forma fechada para geometrias
  arbitrárias.} As expressões~\eqref{eq:delta_xy} e~\eqref{eq:delta_z}, em termos
  dos fatores de abertura $S$ e $q$, reproduzem a imagem com erro inferior a
  $1,6\,\%$ em 34 geometrias, dentro e fora do eixo. As
  Proposições~\ref{prop:limite_xy} e~\ref{prop:limite_z} estabelecem que um receptor
  à mesma altura nunca melhora a resolução a $-3\dB$ e que o ganho vertical de
  qualquer geometria biestática está limitado a $4,6\,\%$ pela refração.

  \item \textbf{Refutação documentada do critério de Brewster.} Três evidências
  independentes --- domínio do alcance sobre a transmitância, invariância azimutal
  da potência e equivalência estatística sob orientação aleatória --- mostram que a
  potência recebida não discrimina posições do receptor.

  \item \textbf{Fusão coerente normalizada e mecanismo do ganho.} A fusão a
  orçamento constante separa o ganho geométrico do ganho por número de amostras; o
  ganho é explicado pelo preenchimento do disco de números de onda por anéis
  concêntricos cujos lóbulos de Bessel se cancelam parcialmente.

  \item \textbf{Regra de polarização do receptor remoto.} O descasamento co-polar
  (Proposição~\ref{prop:descasamento}) anula a complementaridade a
  $\Delta\phi\approx90^\circ$ com alvos reais; um receptor alinhado de dupla
  polarização torna a melhoria de contraste robusta à orientação.

  \item \textbf{Nulo biestático do clutter difuso.} A
  Proposição~\ref{prop:nulo_vv} generaliza o ângulo de Brewster ao espalhamento
  biestático difuso; o nulo VV resultante, ao contrário do nulo HH, sobrevive com
  alvos reais e define um modo de detecção limitado por clutter.

  \item \textbf{Projeto da trajetória por métricas de imagem e requisitos de
  sincronização.} Uma função de custo sem termo de potência, com fronteira de Pareto
  explícita, conduz a uma trajetória de separação varrida com custo de voo quase
  nulo; quantificam-se as tolerâncias de fase ($<30^\circ$), de deriva de relógio
  ($\ll1\,$ns) e de deriva de posição ($<5\,$cm).

  \item \textbf{Correções ao formalismo de base.} A constante da largura de
  Bessel a $-3\dB$ é $1,1264$, e não $1,20$ ($6,5\,\%$ de superestimação); as
  contribuições de banda e de abertura à resolução vertical somam-se em quadratura;
  o coeficiente de transmissão da formulação ABCD com impedâncias assimétricas foi
  corrigido e verificado pela conservação de energia.
\end{enumerate}

\section{Resultados Negativos}
\label{sec:concl_negativos}

A contribuição desta tese não é a afirmação de que ``o biestático é melhor'', mas a
delimitação de quando, onde e por que ele acrescenta informação. Os resultados
negativos fazem parte dessa delimitação:
\begin{enumerate}
  \item A posição de Brewster não maximiza a potência recebida, e a potência não
  discrimina geometrias.
  \item A configuração diametral não serve para imageamento, e sua fusão com o
  canal monoestático é pior que o monoestático isolado.
  \item A fusão não melhora a resolução a $-3\dB$ em nenhuma separação, e o receptor
  à mesma altura não acrescenta resolução vertical, nem pela separação nem pelo raio.
  \item A separabilidade de alvos próximos piora com a fusão.
  \item A fusão de pesos iguais dilui a SNR quando o canal biestático é fraco e
  desperdiça o canal biestático quando ele tem pouco clutter.
  \item Mais eco biestático não implica imagem melhor.
  \item A fusão coerente sem calibração de fase pode ser pior que o monoestático.
  \item O pseudo-Brewster não altera o SCR do monoestático nem das separações
  intermediárias, e a órbita que o impõe é pior que a hélice.
  \item Em solos úmidos, nenhuma geometria alcança SCR positivo para um alvo de
  $1\,$m$^2$ a $2\,$m de profundidade a $425\,$MHz.
  \item O nulo HH a $90^\circ$, aparente com o alvo pontual escalar, não melhora o SCR
  de alvos reais: o alvo escalar não serve para avaliar efeitos de polarização.
  \item Uma órbita de altura constante acumula coerentemente, sobre seu eixo, os
  lóbulos distantes em alcance do reflexo do nadir se o espectro não for ponderado.
\end{enumerate}

\section{Limitações}
\label{sec:concl_lim}

\begin{itemize}
  \item \textbf{Modelo dos alvos.} Condutores perfeitos por Óptica Física, sem
  difração de arestas, ondas de superfície nem reflexões múltiplas; o tubo está no
  limite de validade da aproximação. Apenas componentes co-polares.
  \item \textbf{Modelo da superfície.} Espalhamento de primeira ordem
  ($k_0s\le0,36$) mais especular coerente; os termos de segunda ordem preenchem os
  nulos, de modo que o ganho do nulo VV é uma cota otimista. Facetas limitadas a
  $\pm12\,$m.
  \item \textbf{Modelo do solo.} Homogêneo, com permissividade constante na banda,
  sem clutter volumétrico; o estudo multicamada limitou-se à transmitância.
  \item \textbf{Abertura.} Todos os resultados usam duas voltas completas da hélice;
  com aberturas parciais, o suporte em $k$ deixa de ser anular e a separação ótima
  pode mudar.
  \item \textbf{Estatística.} O SCR na imagem foi estimado com duas a três
  realizações da superfície (dispersão de cerca de $\pm3\dB$); os ganhos do nulo VV
  superam com margem essa dispersão, mas as diferenças entre separações
  intermediárias não.
  \item \textbf{Seleção da trajetória.} A trajetória do modo de imageamento foi
  selecionada com o alvo pontual.
  \item \textbf{Validação experimental.} Os resultados são de simulação; o
  Capítulo~\ref{cap:dados_reais} define o protocolo de validação com dados reais.
\end{itemize}

\section{Perspectivas}
\label{sec:concl_fut}

\begin{enumerate}
  \item \textbf{Reavaliação da trajetória com alvos reais.} Reavaliar os melhores
  candidatos da fronteira de Pareto com alvos anisotrópicos em orientação
  aleatória, receptor alinhado e clutter, usando a mediana do contraste e a
  probabilidade de melhoria como métricas de imagem.
  \item \textbf{Métodos de combinação dos canais.} Comparar a fusão de pesos iguais,
  a fusão ponderada (pesos por canal que maximizem o SCR ou minimizem o ISLR), a
  reconstrução conjunta por inversão regularizada, que pode estimar a fase entre
  canais como incógnita, e a fusão incoerente.
  \item \textbf{Validação do nulo VV.} Avaliar o nulo com modelos de espalhamento de
  ordem superior (IEM/AIEM~\citar{fung1994}, segunda ordem), polarização cruzada da
  superfície e do alvo, correlação exponencial e rugosidade multiescala, e medir
  quanto o nulo se preenche e qual precisão de $\Delta\phi$ ele exige.
  \item \textbf{Modelo de erro de sincronização.} Simular separadamente os termos
  de~\eqref{eq:erro_fase} e estimá-los pulso a pulso a partir do sinal direto, com
  SNR realista.
  \item \textbf{Aberturas parciais.} Repetir a varredura de separação e o projeto da
  trajetória com arcos de $45^\circ$, $90^\circ$, $180^\circ$ e $360^\circ$.
  \item \textbf{Domínio de validade em solos úmidos.} Calcular a profundidade máxima
  com SCR positivo para cada solo e sua dependência com a frequência, já que a
  atenuação cresce aproximadamente com $f$.
  \item \textbf{Validação experimental} segundo o protocolo do
  Capítulo~\ref{cap:dados_reais}.
\end{enumerate}

%%═══════════════════════════════════════════════════════════════════════════════
\appendix
%%═══════════════════════════════════════════════════════════════════════════════

\chapter{Derivações Matemáticas}
\label{app:math}

\section{Conservação de Energia nos Coeficientes de Fresnel TM}
\label{app:energia_fresnel}

Da continuidade do fluxo de potência normal à interface:
\[
  \frac{|E_0|^2}{2\eta_1}\cos\theta_i(1-|r_{\text{TM}}|^2)
  = \frac{|t_{\text{TM}}|^2|E_0|^2}{2\eta_2}\cos\theta_t
  \quad\Rightarrow\quad
  T_{\text{TM}} = \frac{n_2\cos\theta_t}{n_1\cos\theta_i}|t_{\text{TM}}|^2 .
\]
Substituindo~\eqref{eq:r_TM}:
\[
  T_{\text{TM}} =
  \frac{n_2\cos\theta_t}{n_1\cos\theta_i}
  \cdot\frac{4n_1^2\cos^2\theta_i}{(n_2\cos\theta_i+n_1\cos\theta_t)^2}
  = \frac{4n_1 n_2\cos\theta_i\cos\theta_t}{(n_2\cos\theta_i+n_1\cos\theta_t)^2},
\]
e o balanço é
\[
  R_{\text{TM}} + T_{\text{TM}} =
  \frac{(n_2\cos\theta_i-n_1\cos\theta_t)^2 + 4n_1 n_2\cos\theta_i\cos\theta_t}
       {(n_2\cos\theta_i+n_1\cos\theta_t)^2} = 1 .
\]

\section{PSF Circular como Função de Bessel}
\label{app:bessel}

Seja $\alpha$ uniformemente distribuído em $[0,2\pi]$ e $k_x = R_k\cos\alpha$. A
densidade de $k_x$ em $[-R_k, R_k]$ é a distribuição arco-seno,
$p(k_x) = [\pi\sqrt{R_k^2-k_x^2}]^{-1}$, cuja transformada de Fourier é
\[
  \hat{p}(\delta) = \int_{-R_k}^{R_k}
  \frac{e^{\jj k_x\delta}}{\pi\sqrt{R_k^2-k_x^2}}\,\dd k_x
  \stackrel{k_x=R_k\cos\varphi}{=}
  \frac{1}{\pi}\int_0^\pi e^{\jj R_k\delta\cos\varphi}\,\dd\varphi
  = J_0(R_k\delta),
\]
pela representação integral $J_0(z) = \frac{1}{\pi}\int_0^\pi e^{\jj z\cos\varphi}\dd\varphi$.
Em duas dimensões, um anel uniforme de raio $R_k$ produz a PSF radial
$J_0(R_k\rho)$. A largura a $-3\dB$ da PSF em amplitude resolve
$J_0(x_{J_0})=1/\sqrt2$, cuja menor raiz positiva é $x_{J_0}=1,1264$; a largura
total é $2x_{J_0}/R_k$. O primeiro lóbulo lateral de $J_0$ tem amplitude $0,403$,
isto é, $-7,9\dB$, que é o PSLR da imagem monoestática de abertura circular
completa.

\section{Largura Vertical Equivalente}
\label{app:dz}

Seja o suporte vertical gerado por duas variáveis independentes e aproximadamente
uniformes: a frequência, com largura $B$, e o fator $q$, com desvio padrão
$\sigma_q$. Para $K_z=2\pi fq/c$, linearizando em torno de $(f_c,\bar q)$,
$\delta K_z\approx(2\pi/c)(\bar q\,\delta f+f_c\,\delta q)$, e a variância do suporte
é $(2\pi/c)^2(\bar q^2B^2/12+f_c^2\sigma_q^2)$. A curvatura do lóbulo principal da
PSF é proporcional a essa variância, de modo que a PSF vertical equivale, até
segunda ordem, à de um único suporte uniforme de largura
$\sqrt{(B\bar q)^2+12f_c^2\sigma_q^2}$ (em unidades de frequência), cuja largura a
$-3\dB$ é $0,886\,c$ dividido por essa largura. Com janela de Hann, a constante
passa a $1,443$. A soma linear das contribuições corresponderia a supor
correlação perfeita entre $\delta f$ e $\delta q$, o que não ocorre, pois a
frequência e a posição na abertura são varridas independentemente.

\section{Derivação das Matrizes ABCD a partir das Equações de Maxwell}
\label{app:abcd_deriv}

Na camada $i$, em polarização TM, $H_y(z) = C_i^+e^{\jj k_{z,i}z} + C_i^-e^{-\jj k_{z,i}z}$.
As condições de contorno TM (continuidade de $H_y$ e de
$(1/n_i^2)\partial_z H_y\propto E_x$) nas faces da camada equivalem à propagação
do vetor de estado $[H_y,\,E_x]^T$ ao longo da espessura $d_i$:
\begin{align*}
  H_y(z+d_i) &= H_y(z)\cos(k_{z,i}d_i) + \jj Z_i E_x(z)\sin(k_{z,i}d_i), \\
  E_x(z+d_i) &= \frac{\jj}{Z_i}H_y(z)\sin(k_{z,i}d_i) + E_x(z)\cos(k_{z,i}d_i),
\end{align*}
que é a matriz $\mathbf{M}_i$ de~\eqref{eq:ABCD_matrix}. Impondo, na saída, apenas
a onda transmitida no meio semi-infinito ($E_x=H_y/Z_{\text{out}}$) e, na entrada,
a superposição das ondas incidente e refletida, obtêm-se $\Gamma_{\text{TM}}$ e
$\tau_{\text{TM}}$ de~\eqref{eq:Gamma_TM}, com o denominador comum
$AZ_{\text{out}}+B+Z_{\text{in}}(CZ_{\text{out}}+D)$.

\chapter{Parâmetros das Simulações}
\label{app:parametros}

\begin{table}[H]
\centering\small
\caption{Parâmetros comuns às simulações dos Capítulos~\ref{cap:geometria}
a~\ref{cap:trajetoria}.}
\label{tab:app_parametros}
\begin{tabular}{@{}L{2.6cm}L{4.9cm}L{7.0cm}@{}}
\toprule
Grupo & Parâmetro & Valor \\
\midrule
Sinal & Frequência portadora & $425\,$MHz ($\lambda_0=0,706\,$m) \\
      & Largura de banda e amostras & $50\,$MHz; 64 frequências \\
      & Potência de pico; duração do pulso & $100\,$kW; $1,3\,\mu$s \\
      & Temperatura de ruído & $300\,$K \\
Antenas & Ganho; feixe & $6\,$dBi; gaussiano $14^\circ\times50^\circ$ apontado ao centro \\
Transceptor & Trajetória & hélice de 2 voltas \\
            & Raio; altura & $137,4\to165,5\,$m; $131,5\to93,6\,$m \\
            & Incidência no centro da cena & $46,3^\circ$ a $60,5^\circ$ (média $53,4^\circ$) \\
            & Pulsos & 1600 (800 nos estudos de Monte Carlo) \\
            & Velocidade de referência & $10\,$m/s ($190\,$s de aquisição) \\
Receptor & Altura & igual à do transceptor \\
         & Raio; separação & $\kappa\rho_{Tx}$; $\Delta\phi(t)$ \\
Solo & Permissividades & S0 $4$; S1 $3,99-\jj0,40$; S2 $6,5-\jj0,9$; S3 $12-\jj2,5$ \\
Superfície & Rugosidade de referência & $s=2\,$cm, $\ell=20\,$cm (gaussiana) \\
           & Facetas & $0,15\times0,15\,$m em $\pm12\,$m; amplitude congelada \\
Alvo de referência & Ponto isotrópico & $1\,$m$^2$ em $(0;\,0;\,-2)\,$m \\
Imagem & Retroprojeção & refração exata por voxel e pulso; sobreamostragem $32\times$ \\
       & Janela espectral & retangular (forma da PSF); Hann (clutter) \\
       & Janela de recepção & $[-30,\,+60]\,$m em torno do percurso Tx--centro--Rx \\
       & Faixa dinâmica de exibição & $30\dB$, pico monoestático em $0\dB$ \\
\bottomrule
\end{tabular}
\end{table}

\chapter{Lista de Publicações}
\label{app:pub}

% TODO: Inserir lista completa de publicações (artigos em periódico,
%       conferências e relatórios técnicos) originados desta tese.

\textit{[Lista de publicações a ser inserida.]}

%%═══════════════════════════════════════════════════════════════════════════════
\backmatter
%%═══════════════════════════════════════════════════════════════════════════════

\chapter*{Referências}
\addcontentsline{toc}{chapter}{Referências}

% As citações no corpo do texto usam \citar{<chave>} = [\ref{ref:<chave>}],
% de modo que a numeração permanece correta se a lista for reordenada.

\begin{enumerate}

  \item \label{ref:moreira2013}
  A.\ Moreira, P.\ Prats-Iraola, M.\ Younis, G.\ Krieger, I.\ Hajnsek e
  K.\ P.\ Papathanassiou, ``A tutorial on synthetic aperture radar,''
  \textit{IEEE Geoscience and Remote Sensing Magazine}, vol.\ 1, n.\ 1,
  pp.\ 6--43, Mar.\ 2013.

  \item \label{ref:munson1983}
  D.\ C.\ Munson, J.\ D.\ O'Brien e W.\ K.\ Jenkins, ``A tomographic
  formulation of spotlight-mode synthetic aperture radar,''
  \textit{Proceedings of the IEEE}, vol.\ 71, n.\ 8, pp.\ 917--925, Ago.\ 1983.

  \item \label{ref:goes2022}
  J.\ A.\ Góes, \textit{Techniques for High-Resolution 3D Images with
  Synthetic Aperture Radar}, Tese de Doutorado, Universidade Estadual de
  Campinas (UNICAMP), Campinas, SP, Brasil, 2022.

  \item \label{ref:arikan1988}
  O.\ Arikan e D.\ C.\ Munson Jr., ``A tomographic formulation of bistatic
  synthetic aperture radar,'' in \textit{Proc.\ IEEE ICASSP}, pp.\ 1099--1102,
  1988.

  \item \label{ref:bornwolf1999}
  M.\ Born e E.\ Wolf, \textit{Principles of Optics: Electromagnetic Theory
  of Propagation, Interference and Diffraction of Light}, 7ª~ed.,
  Cambridge: Cambridge University Press, 1999.

  \item \label{ref:ulaby1981}
  F.\ T.\ Ulaby, R.\ K.\ Moore e A.\ K.\ Fung, \textit{Microwave Remote
  Sensing: Active and Passive}, vol.\ 1, Reading, MA: Addison-Wesley, 1981.

  \item \label{ref:ishimaru1991}
  A.\ Ishimaru, \textit{Electromagnetic Wave Propagation, Radiation, and
  Scattering}, Englewood Cliffs, NJ: Prentice Hall, 1991.

  \item \label{ref:horne2002}
  A.\ M.\ Horne e G.\ Yates, ``Bistatic synthetic aperture radar,''
  in \textit{Proc.\ IEE Radar Conference}, QinetiQ Malvern, 2002.

  \item \label{ref:garcia2019}
  M.\ García-Fernández, Y.\ López-Fernández e F.\ Las-Heras,
  ``Autonomous airborne 3D SAR imaging system for subsurface sensing,''
  \textit{Remote Sensing}, vol.\ 11, n.\ 20, p.\ 2357, 2019.

  \item \label{ref:banda2016}
  F.\ Banda, J.\ Dall e S.\ Tebaldini,
  ``Single and multipolarimetric P-band SAR tomography of subsurface ice,''
  \textit{IEEE Transactions on Geoscience and Remote Sensing}, vol.\ 54, n.\ 5,
  pp.\ 2832--2845, 2016.

  \item \label{ref:pozar2011}
  D.\ M.\ Pozar, \textit{Microwave Engineering}, 4ª~ed.,
  Hoboken, NJ: Wiley, 2011.

  \item \label{ref:cumming2005}
  I.\ G.\ Cumming e F.\ H.\ Wong, \textit{Digital Processing of Synthetic
  Aperture Radar Data: Algorithms and Implementation},
  Norwood, MA: Artech House, 2005.

  \item \label{ref:willis1995}
  N.\ J.\ Willis, \textit{Bistatic Radar}, 2ª~ed.,
  Norwood, MA: Artech House, 1995.

  \item \label{ref:daniels2004}
  D.\ J.\ Daniels, \textit{Ground Penetrating Radar}, 2ª~ed.,
  London: IET, 2004.

  \item \label{ref:knaell1995}
  K.\ K.\ Knaell e G.\ P.\ Cardillo, ``Radar tomography for the generation of
  three-dimensional images,''
  \textit{IEE Proceedings — Radar, Sonar and Navigation}, vol.\ 142, n.\ 2,
  pp.\ 54--60, Abr.\ 1995.

  \item \label{ref:reigber2000}
  A.\ Reigber e A.\ Moreira, ``First demonstration of airborne SAR tomography
  using multibaseline L-band data,''
  \textit{IEEE Transactions on Geoscience and Remote Sensing}, vol.\ 38, n.\ 5,
  pp.\ 2142--2152, Set.\ 2000.

  \item \label{ref:ponce2014}
  O.\ Ponce, P.\ Prats-Iraola, M.\ Pinheiro, M.\ Rodriguez-Cassola,
  R.\ Scheiber, A.\ Reigber e A.\ Moreira,
  ``Fully polarimetric high-resolution 3-D imaging with circular SAR at L-band,''
  \textit{IEEE Transactions on Geoscience and Remote Sensing}, vol.\ 52, n.\ 6,
  pp.\ 3274--3290, Jun.\ 2014.

  \item \label{ref:zhu2010}
  X.\ X.\ Zhu e R.\ Bamler, ``Very high resolution spaceborne SAR tomography
  in urban environment,''
  \textit{IEEE Transactions on Geoscience and Remote Sensing}, vol.\ 48, n.\ 12,
  pp.\ 4296--4308, Dez.\ 2010.

  \item \label{ref:elachi2006}
  C.\ Elachi e J.\ van Zyl, \textit{Introduction to the Physics and Techniques
  of Remote Sensing}, 2ª~ed., Hoboken, NJ: Wiley, 2006.

  % ─── Referências adicionadas: revisão bibliográfica SAR biestático + ────
  % ─── ângulo de Brewster (setembro de 2026) ───────────────────────────────

  \item \label{ref:ding2002}
  Y.\ Ding e D.\ C.\ Munson Jr., ``A fast back-projection algorithm for
  bistatic SAR imaging,'' in \textit{Proc.\ IEEE Int.\ Conf.\ on Image
  Processing (ICIP)}, vol.\ II, pp.\ 449--452, 2002.

  \item \label{ref:homer2005}
  J.\ Homer, E.\ Donskoi, B.\ Mojarrabi, J.\ Palmer e K.\ Kubik,
  ``Three-dimensional bistatic synthetic aperture radar imaging system:
  spatial resolution analysis,'' \textit{IEE Proceedings --- Radar, Sonar
  and Navigation}, vol.\ 152, n.\ 6, pp.\ 391--394, 2005.

  \item \label{ref:garza2010}
  G.\ Garza e Z.\ Qiao, ``Resolution analysis of bistatic SAR,'' in
  \textit{Proc.\ SPIE 8021, Radar Sensor Technology XV}, 80210F, 2010.

  \item \label{ref:chen2022}
  Y.\ Chen, R.\ Chen, H.\ Liu, J.\ Guo, Y.\ Wang e J.\ Zhang, ``Research on
  configuration constraints of airborne bistatic SARs,'' \textit{Sensors},
  vol.\ 22, n.\ 17, art.\ 6534, 2022.

  \item \label{ref:zhang2018}
  H.\ Zhang, J.\ Tang, R.\ Wang, Y.\ Deng, W.\ Wang e N.\ Li, ``An
  accelerated backprojection algorithm for monostatic and bistatic SAR
  processing,'' \textit{Remote Sensing}, vol.\ 10, n.\ 1, art.\ 140, 2018.

  \item \label{ref:wang2020}
  Z.\ Wang, M.\ Liu, G.\ Ai, P.\ Wang e K.\ Lv, ``Focusing of bistatic SAR
  with curved trajectory based on extended azimuth nonlinear chirp
  scaling,'' \textit{IEEE Transactions on Geoscience and Remote Sensing},
  vol.\ 58, n.\ 6, pp.\ 4160--4179, 2020.

  \item \label{ref:sun2022}
  H.\ Sun, Z.\ Sun, T.\ Chen, Y.\ Miao, J.\ Wu e J.\ Yang, ``An efficient
  backprojection algorithm based on wavenumber-domain spectral splicing for
  monostatic and bistatic SAR configurations,'' \textit{Remote Sensing},
  vol.\ 14, n.\ 8, art.\ 1885, 2022.

  \item \label{ref:zhao2025}
  T.\ Zhao, J.\ Wang, Z.\ Cheng, Z.\ Huang e J.\ Song, ``A subimage
  autofocus bistatic ground Cartesian back-projection algorithm for passive
  bistatic SAR based on GEO satellites,'' \textit{Remote Sensing}, vol.\ 17,
  n.\ 9, art.\ 1576, 2025.

  \item \label{ref:ge2019}
  N.\ Ge e X.\ X.\ Zhu, ``Bistatic-like differential SAR tomography,''
  \textit{IEEE Transactions on Geoscience and Remote Sensing}, vol.\ 57,
  n.\ 8, pp.\ 5883--5893, 2019.

  \item \label{ref:chenz2024}
  Z.\ Chen, Y.\ Li, C.\ Hu et al., ``Repeat-pass space-surface bistatic SAR
  tomography: accurate imaging and first experiment,'' \textit{Science
  China Information Sciences}, vol.\ 67, art.\ 192304, 2024.

  \item \label{ref:wu2025}
  J.\ Wu, J.\ Yang, Z.\ Li, W.\ Pu, Z.\ Sun, H.\ An e Q.\ Yang, ``Review of
  bistatic synthetic aperture radar imaging methods,'' \textit{Journal of
  Radars}, vol.\ 14, n.\ 5, pp.\ 1115--1141, 2025.

  \item \label{ref:he2025}
  J.\ He, H.\ Xie, H.\ Liu, Z.\ Wu, B.\ Xu, N.\ Zhu, Z.\ Lu e P.\ Qin,
  ``Multi-baseline bistatic SAR three-dimensional imaging method based on
  phase error calibration combining PGA and EB-ISOA,'' \textit{Remote
  Sensing}, vol.\ 17, n.\ 3, art.\ 363, 2025.

  \item \label{ref:andre2025}
  D.\ Andre, V.\ P.\ Krishnan e C.\ Nolan, ``Artefact analysis of
  multistatic three-dimensional SAR imaging with two linear
  trajectories,'' arXiv:2503.11323, 2025.

  \item \label{ref:liu2024}
  M.\ Liu et al., ``Unmanned airborne bistatic interferometric synthetic
  aperture radar data processing method using bi-directional
  synchronization chain signals,'' \textit{Remote Sensing}, vol.\ 16, n.\ 5,
  art.\ 769, 2024.

  \item \label{ref:grathwohl2023}
  A.\ Grathwohl, B.\ Meinecke, M.\ Widmann, J.\ Kanz e C.\ Waldschmidt,
  ``UAV-based bistatic SAR-imaging using a stationary repeater,''
  \textit{IEEE Journal of Microwaves}, vol.\ 3, n.\ 2, pp.\ 625--634, 2023.

  \item \label{ref:kanz2025}
  J.\ Kanz, C.\ Gesell, C.\ Bonfert, D.\ Werbunat, A.\ Grathwohl,
  J.\ Aguilar, M.\ Vossiek e C.\ Waldschmidt, ``UAV-borne digital radar
  system for coherent multistatic SAR imaging,'' arXiv:2507.20792, 2025.

  \item \label{ref:ouchi2017}
  K.\ Ouchi e C.-S.\ Yang, ``Direct evidence and implications of Brewster's
  angle damping observed by synthetic aperture radar images,'' in
  \textit{Proc.\ IGARSS 2017}, 2017.

  \item \label{ref:ouchi2018}
  K.\ Ouchi e C.-S.\ Yang, ``Brewster angle damping observed in the
  TerraSAR-X synthetic aperture radar images of man-made targets,''
  \textit{IEEE Geoscience and Remote Sensing Letters}, vol.\ 15, n.\ 4,
  pp.\ 532--536, 2018.

  \item \label{ref:thirion2018}
  L.\ Thirion-Lefevre e R.\ Guinvarc'h, ``The double Brewster angle
  effect,'' \textit{Comptes Rendus Physique}, vol.\ 19, n.\ 1--2,
  pp.\ 43--53, 2018.

  \item \label{ref:orfeo2019}
  D.\ Orfeo, Y.\ Zhang, D.\ Burns, J.\ S.\ Miller, D.\ R.\ Huston e T.\ Xia,
  ``Bistatic antenna configurations for air-launched ground penetrating
  radar,'' \textit{Journal of Applied Remote Sensing}, vol.\ 13, n.\ 2,
  art.\ 027501, 2019.


  % ─── Referências adicionadas (outubro de 2026): espalhamento de superfície, ──
  % ─── RCS, estatística e sincronização biestática ───────────────────────────

  \item \label{ref:kell1965}
  R.\ E.\ Kell, ``On the derivation of bistatic RCS from monostatic
  measurements,'' \textit{Proceedings of the IEEE}, vol.\ 53, n.\ 8,
  pp.\ 983--988, 1965.

  \item \label{ref:knott1993}
  E.\ F.\ Knott, J.\ F.\ Shaeffer e M.\ T.\ Tuley, \textit{Radar Cross
  Section}, 2ª~ed., Norwood, MA: Artech House, 1993.

  \item \label{ref:ament1953}
  W.\ S.\ Ament, ``Toward a theory of reflection by a rough surface,''
  \textit{Proceedings of the IRE}, vol.\ 41, n.\ 1, pp.\ 142--146, 1953.

  \item \label{ref:ulaby1982}
  F.\ T.\ Ulaby, R.\ K.\ Moore e A.\ K.\ Fung, \textit{Microwave Remote
  Sensing: Active and Passive}, vol.\ II, Reading, MA: Addison-Wesley, 1982.

  \item \label{ref:fung1994}
  A.\ K.\ Fung, \textit{Microwave Scattering and Emission Models and Their
  Applications}, Norwood, MA: Artech House, 1994.

  \item \label{ref:roy2007}
  O.\ Roy e M.\ Vetterli, ``The effective rank: a measure of effective
  dimensionality,'' in \textit{Proc.\ European Signal Processing Conference
  (EUSIPCO)}, pp.\ 606--610, 2007.

  \item \label{ref:shoemake1992}
  K.\ Shoemake, ``Uniform random rotations,'' in D.\ Kirk (ed.),
  \textit{Graphics Gems III}, pp.\ 124--132, San Diego: Academic Press, 1992.

  \item \label{ref:efron1993}
  B.\ Efron e R.\ J.\ Tibshirani, \textit{An Introduction to the Bootstrap},
  New York: Chapman \& Hall, 1993.

  \item \label{ref:krieger2006}
  G.\ Krieger e A.\ Moreira, ``Spaceborne bi- and multistatic SAR: potential
  and challenges,'' \textit{IEE Proceedings --- Radar, Sonar and Navigation},
  vol.\ 153, n.\ 3, pp.\ 184--198, 2006.

  \item \label{ref:younis2006}
  M.\ Younis, R.\ Metzig e G.\ Krieger, ``Performance prediction of a phase
  synchronization link for bistatic SAR,'' \textit{IEEE Geoscience and Remote
  Sensing Letters}, vol.\ 3, n.\ 3, pp.\ 429--433, 2006.

\end{enumerate}

\end{document}
